% concatenated from main.tex, sections/01_introduction.tex, sections/02_preliminaries.tex, sections/02_preliminaries/01_notations.tex, sections/02_preliminaries/02_assumptions.tex, sections/02_preliminaries/03_noise.tex, sections/02_preliminaries/04_clt.tex, sections/02_preliminaries/05_related_works.tex, sections/03_general_results.tex, sections/03_general_results/01_last_iterates_v2.tex, sections/03_general_results/02_averaging.tex, sections/03_general_results/03_tails.tex, sections/04_applications.tex, sections/04_applications/01_weibull_gaussian.tex, sections/05_discussion.tex, sections/05_discussion/02_machinery.tex, sections/05_discussion/01_accelerated_last_vs_pr.tex, sections/06_conclusion.tex, sections/07_acknowledgements.tex, appendix_main.tex, appendix/sections/00_preliminaries.tex, appendix/sections/preliminaries/recursion_convergence_rate.tex, appendix/sections/preliminaries/bdg.tex, appendix/sections/preliminaries/weighted_clt.tex, appendix/sections/preliminaries/yk_variation.tex, appendix/sections/01_general.tex, components/02_proofs/lemmas/nonlinear_remainder.tex, appendix/sections/general/f_bound_v2.tex, components/02_proofs/lemmas/Wp_recursion.tex, components/02_proofs/theorems/general.tex, components/02_proofs/lemmas/pr_remainder.tex, components/02_proofs/corollaries/high_probability.tex, appendix/sections/02_lemmas.tex, appendix/sections/general/e_state_dependent_v2.tex, components/02_proofs/lemmas/ou_discretization.tex, components/02_proofs/lemmas/ct_ou_convergence.tex
\documentclass[11pt]{article}

\usepackage{preamble_arxiv}

\usepackage{titling}
\predate{}
\postdate{}

\title{Finite-Sample Wasserstein Error Bounds and Concentration Inequalities for Nonlinear Stochastic Approximation
}
\author{Seo Taek Kong \and R. Srikant}

\date{}


\begin{document}


\maketitle
\let\thefootnote\relax\footnotetext{University of Illinois Urbana-Champaign. Emails: \texttt{\{skong10, rsrikant\}@illinois.edu}}


\begin{abstract}
This paper derives \nonasymptotic error bounds for \nonlinear stochastic approximation algorithms in the Wasserstein-$p$ distance. To obtain explicit finite-sample guarantees for the last iterate, we develop a coupling argument that compares the discrete-time process to a limiting Ornstein-Uhlenbeck process. Our analysis applies to algorithms driven by general noise conditions, including martingale differences and functions of ergodic Markov chains. Complementing this result, we handle the convergence rate of the Polyak-Ruppert average through a direct analysis that applies under the same general setting.

Assuming the driving noise satisfies a \nonasymptotic central limit theorem, we show that the normalized last iterates converge to a Gaussian distribution in the $p$-Wasserstein distance at a rate of order $\gamma_n^{1/6}$, where $\gamma_n$ is the \stepsize. Similarly, the Polyak-Ruppert average is shown to converge in the Wasserstein distance at a rate of order $n^{-1/6}$. These distributional guarantees imply high-probability concentration inequalities that improve upon those derived from moment bounds and Markov's inequality. We demonstrate the utility of this approach by considering two applications: (1) linear stochastic approximation, where we explicitly quantify the transition from heavy-tailed to Gaussian behavior of the iterates, thereby bridging the gap between recent finite-sample analyses and asymptotic theory and (2) stochastic gradient descent, where we establish rate of convergence to the central limit theorem. 
    
\end{abstract}



\tableofcontents




\section{Introduction}
Stochastic approximation (SA) provides a versatile framework for solving fixed point problems in the presence of noisy observations. 
These algorithms have become foundational to numerous fields including statistics, control theory, reinforcement learning, and machine learning. 
The algorithm seeks a solution $x^*$ to the equation
\begin{equation}\label{eq:sa_objective}
    \mathbb{E}_\pi \left[ f (x^*, \xi)\right]  = 0 , 
\end{equation}
where $\pi$ is some probability measure over $\xi$, by initializing some $x_1$ and using the recursion
\begin{equation}\label{eq:sa}
    x_{k+1} = x_k - \gamma_k (f(x_k, \xi_k) + W_{k})  ,
\end{equation}
where $\{\gamma_k\}$ is a diminishing \stepsize sequence and $\{(\xi_k, W_{k})\}$ is a noise sequence with $\mathbb{E}[W_{k}] = 0$ for all $k$. 




The theoretical analysis of SA often begins with the ODE method.
This approach was founded in the seminal works of \citep{meerkov1972simplified,derevitskii1974two,ljung1977analysis,kushner1978stochastic}, and connects the discrete-time recursion to an associated ordinary differential equation (ODE).
This idea was later advanced in \citep{borkar2000ode,borkar2008stochastic,borkar25}, and is widely used to identify asymptotic properties of the algorithm. 
Under appropriate conditions, this method guarantees almost sure convergence of the algorithm. 


A refined statistical understanding is developed by characterizing the fluctuations around this limit.
This is established via diffusion approximations \citep{kushner_yin} or the functional central limit theorem (FCLT) \citep{borkar25}, providing a complete picture of the algorithm's asymptotic distribution.


While ODE and diffusion approximation frameworks provide asymptotic characterizations, recent applications in statistical inference and machine learning have motivated analyses of finite-time guarantees. 
Non-asymptotic analysis of SA algorithms often provides explicit bounds on the moments of the error, most commonly the mean-squared error $\lVert x_n - x^* \rVert^2_{L_2}$. 
These finite-time bounds are invaluable for providing concrete performance guarantees and estimating confidence intervals. 
However, these guarantees are typically derived as upper bounds; while they control the rate of decay, they are often not tight enough to prove that the moments of the scaled error $\gamma_n^{-1/2}(x_n - x^*)$ converge to the asymptotic limit.
Existing literature generally focuses on either asymptotic distributional results, which characterize the limiting law without convergence rates, or \nonasymptotic moment bounds, which provide rates but do not capture the limiting distribution. 


In this paper, we address this distinction by developing a \nonasymptotic analysis of weak convergence for SA. 
Our main result derives explicit bounds on the Wasserstein-$p$ distance between the law of the scaled error and the law of its limiting SDE. 
This approach provides a quantitative estimate of the CLT for stochastic approximation, with implications including: 







\begin{enumerate}
    \item 
    \textbf{Convergence Rates in Wasserstein Metrics:}
    We establish \Cref{thm:general,thm:pr_average} that translate \nonasymptotic central limit theorems for the driving noise into convergence rates in the $p$-Wasserstein metric for the stochastic approximation algorithm's iterates. 
    These bounds are derived contingent on the verification of standard moment and stability conditions, providing a modular framework for analyzing both last-iterate and Polyak-Ruppert averaged estimators. 
    

    \item 
    \textbf{Moment Convergence:}
    In \Cref{sec:applications}, we verify  the assumptions used to obtain \Cref{thm:general,thm:pr_average} for widely-used algorithms, thereby establishing weak convergence rate, convergence of moments up to order $p$, and providing explicit rates for this convergence. 
    This provides a finer characterization than previous \nonasymptotic analyses, which typically yield finite-time upper bounds that do not recover the exact asymptotic moments in the limit.     
    
    \item 
    \textbf{Tail Behavior and Phase Transitions:}
    We utilize the distributional control provided by the Wasserstein bounds to derive high-probability concentration inequalities. 
    Unlike weaker metrics (e.g., Kolmogorov or Wasserstein-1), convergence in $\mathcal{W}_p$ allows us to demonstrate that the heavy tails associated with the finite-sample distribution are transient. 
    We demonstrate the utility of this approach by quantifying the transition from heavy-tailed (Weibull) to Gaussian behavior in linear stochastic approximation, providing a unified view of the finite-sample and asymptotic regimes. 
    


\end{enumerate}






\section{Preliminaries and Problem Setup}
We state the assumptions that will be used throughout our analysis. 


\textbf{Notations:} 
We use $\lVert x \rVert$ to denote a weighted norm defined as $\lVert x \rVert = \sqrt{x^\dagger Q x}$ for a positive definite matrix $Q$ defined via the Lyapunov equation \eqref{eq:lyapunov}. 
For a matrix $A$, $\lVert A \rVert$ is used to denote the corresponding operator norm. 
The $L_p$ norm $(\mathbb{E}\lVert M \rVert^p)^{1/p}$ of a random vector $M$ is denoted by $\lVert M \rVert_{L_p}$, and similarly $\lVert A \rVert_{L_p}$ is used to denote the $L_p$ norm of the operator norm of a matrix $A$. 
For a positive definite matrix $\Gamma$, we denote by $\sqrt{\Gamma}$ its unique square root. 
We use $Z$ to denote a standard Gaussian random variable or $\{Z_m\}$ to denote an i.i.d. standard Gaussian sequence. 
Throughout the paper, we use $\mathcal{H}_k$ to denote the filtration generated by the random variables $\{x_1, \xi_1, W_1, \cdots, \xi_{k}, W_{k}\}$. 
We use $\mathbb{E}_{k-1} [\cdot]$ to denote the conditional expectation $\mathbb{E}[\cdot | \mathcal{H}_{k-1}]$.
The identity matrix is denoted by $\identity$, and $\lVert \cdot \rVert_{HS}$ is the Hilbert-Schmidt norm of a matrix. 
For a (uniformly) Lipschitz continuous map $\phi$, we use $\lip(\phi)$ to denote its Lipschitz constant. 


We state the assumptions used to establish our main results.
\begin{assumption}\label{ass:step_size}
    The \stepsize $\{\gamma_k\}_{k \geq 1}$ is chosen as $\gamma_k = \gamma_1 \cdot k^{-a}$ for some $a \in (0, 1]$ and $\gamma_1 > 0$. 
\end{assumption}
\begin{assumption}\label{ass:jacobian}
    The operator $f$ and its Jacobian $\nabla f$ are uniformly Lipschitz continuous in $x$. 
\end{assumption}
We use $A(\xi) = \nabla f(x^*, \xi)$ to denote the Jacobian of $f$ at $x^*$.
Next, we state our assumptions on the noise.
\begin{assumption}\label{ass:noise_markov}
    Let $\{\xi_k\}$ be a Markov chain on a general state space $\mathcal{S}$ with transition kernel $P$. 
    We assume:
    \begin{assumptionsub}
            \item The chain is $\psi$--irreducible and aperiodic. \label{ass:noise_markov:irreducibility}

        \item There exists a test function $V: \mathcal{S} \to [1, \infty)$ such that for some $\rho_V \in [0, 1)$ and $C_V > 0$, 
    \begin{equation*}
        P V(\xi_k) \leq \rho_V V(\xi_k) + C_V , \qquad \forall k .
    \end{equation*}
    \label{ass:noise_markov:drift}

        \item The initial distribution satisfies $\lvert V(\xi_1) \rvert_{L_{1}} < \infty$.
        \label{ass:noise_markov:initial}

    \end{assumptionsub}    
\end{assumption}
\Cref{ass:noise_markov:irreducibility,ass:noise_markov:drift} imply existence of a unique stationary measure \citep[Theorem 15.0.1]{meyn2012markov}. 
We require that the stationary measure is the probability measure $\pi$ under which the objective \eqref{eq:sa_objective} is defined. 
For the SA algorithm to be stable around the solution $x^*$, the operators $f$ and $\nabla f$ must have sufficiently well-behaved tails. 
This is imposed in \Cref{ass:noise_operator_compatibility} by relating the operators to the test function $V$ in \Cref{ass:noise_markov} as done in \citep{borkar25}. 


Let $p \geq 1$ be the largest integer such that \Crefrange{ass:noise_operator_compatibility}{ass:CLT} are satisfied.
\begin{assumption}\label{ass:noise_operator_compatibility}
    The stationary measure of the Markov chain is $\pi$. 
    The operator $f$ and its Jacobian $A$ satisfy
        \begin{align*}
            \sup_{\xi \in \mathcal{S}} \frac{\lVert f(x^*, \xi) \rVert^{2p}}{V(\xi)} < \infty , 
            \qquad
            \sup_{\xi \in \mathcal{S}} \frac{\lVert A(\xi) \rVert^{2p}}{V(\xi)} < \infty 
        \end{align*}
        for the function $V$ defined in \Cref{ass:noise_markov:drift}.
\end{assumption}
The additive noise sequence $\{W_k\}$ is used to model exogenous noise. 
\begin{assumption}\label{ass:noise_martingale}
    The sequence $\{W_k\}_{k \geq 1}$ is a martingale difference sequence (MDS) with respect to $\{\mathcal{H}_{k}\}_{k \geq 0}$ and is independent of $\{\xi_k\}$. 
    For $\Gamma_k^W = \mathbb{E}_{k-1} [W_k W_k^\dagger]$, the MDS satisfies the following.
    \begin{assumptionsub}
        \item There exists a finite constant $\rosenthal{W}$ such that        
        \begin{equation*}
        \rosenthal{W} \coloneqq \sup_{k \geq 1} \sqrt{\left\lvert 
        \mathbb{E}\left[  \left\lVert M_k \right\rVert^2 | \mathcal{H}_{k-1} \right]
        \right\rvert_{L_{p/2 }}} + \lVert W_k \rVert_{L_{p}} 
        .
    \end{equation*} 
    When $p < 2$, we assume $\kappa_2^W$ exists and use $\rosenthal{W}$ to denote $\rosn{2}{W}$. 

    \label{ass:martingale:rosenthal}

    \item 
    There exists a finite constant $\beta_{2p}^W$ such that
    \begin{equation*}
        \sup_{k \geq 1}\,  \lVert W_k \rVert_{L_{2p}} \leq \beta_{2p}^W . 
    \end{equation*}
    \label{ass:martingale:moments}

        \item  
        There exists some $\Gamma^W \succ 0$, $C_W > 0$, and $\rho_W \in [0, 1)$ such that
    \begin{equation*}
        \lVert \mathbb{E}\Gamma_k^W - \Gamma^W \rVert_{HS} \leq C_W \rho_W^k .
    \end{equation*}
    \label{ass:martingale:covariance}
    \end{assumptionsub}
\end{assumption}
The condition in \Cref{ass:martingale:rosenthal} is needed to apply the martingale version of the Rosenthal inequality as discussed in \Cref{sec:BDG}.
From \Cref{ass:martingale:moments}, it follows from Jensen's inequality that there exists a constant $\beta_p^W$ such that $\sup_{k \geq 1} \lVert W_k \rVert_{L_p} \leq \beta_p^W$. 




Denote by $\bar{f}(x) = \pi(f(x, \cdot))$ the expectation of $f(x, \xi)$ with respect to the stationary measure $\pi$ for each $x$. 
\begin{assumption}\label{ass:general_sa}
    For \stepsize $\{\gamma_k\}$, the Jacobian $\nabla \bar{f}(x^*)$ at $x^*$ is such that
    \begin{equation}\label{eq:hurwtiz}
        -\bar{A}_a = -\lim_{n \to \infty} \left \{\nabla \bar{f}(x^*) - \frac{a}{2 n \gamma_n} \identity \right\}
    \end{equation}    
    is Hurwitz stable. 
\end{assumption}
The value of the matrix $\bar{A}_a$ is determined by the step size exponent $a$ in \Cref{ass:step_size}, where it is equal to $\nabla \bar{f}(x^*)$ when $\lim_{n \to \infty} (n \gamma_n)^{-1} = 0$. 



The Hurwitz condition on $\bar{A}_a$ ensures a contraction occurs in a weighted norm $\lVert \cdot \rVert$ \citep[Lemma 17]{kaledin2020finite}. 
Specifically, let $Q \succ 0$ be the unique solution to the Lyapunov equation
\begin{equation}\label{eq:lyapunov}
    \bar{A}_a Q + Q \bar{A}_a^\dagger = \identity.
\end{equation}
Denote by $\lVert \cdot \rVert_2$ the Euclidean norm when applied to vectors and the spectral norm when applied to matrices.
Define a weighted vector norm $\lVert \cdot \rVert$ and its induced operator norm as
\begin{equation}\label{eq:weighted_norm}
    \lVert x \rVert = \sqrt{x^\dagger Q x} 
    \qquad 
    \text{and} \qquad 
    \lVert \bar{A}_a \rVert^2 
    =\sup_{v \neq 0} \frac{v^\dagger \bar{A}_a^\dagger Q \bar{A}_a v}{v^\dagger Q v} .     
\end{equation}
For all $\gamma\leq (2 \lVert Q \rVert_2^{2} \lVert \bar{A} \rVert^2)^{-1}$, a contraction occurs:
\begin{equation}\label{eq:contraction}
    \lVert \identity - \gamma \bar{A}_a \rVert \leq 1 - \gamma \lambdaDT ,
    \qquad 
    \lambdaDT = \frac{1}{2 \lVert Q \rVert_2^2} 
    . 
\end{equation}
Similarly, the Hurwitz condition implies exponential stability: there exists a constant $K$ such that for all $t > 0$, 
\begin{equation*}
    \left\lVert e^{-\bar{A}_a t} \right\rVert \leq K e^{-\lambdaCT t} , \qquad \lambdaCT = \frac{1}{2\lVert Q \rVert_2}. 
\end{equation*}


When the \stepsize exponent $a$ is set to $1$, \Cref{ass:general_sa} demands that the constant $\gamma_1$ be sufficiently large. 
A direct consequence of this choice is that the contraction in \eqref{eq:contraction} may not hold for early iterates (e.g., at $n = 1$). 
However, this is only a transient effect; the contraction property is always guaranteed to hold for all $n \geq n_0$ for some finite $n_0$. 
Our analysis focuses on the finite-time behavior past this burn-in period, and we simplify the exposition by proceeding as if this contraction holds for all $n \geq 1$.




Much of the literature in stochastic approximation has focused on establishing asymptotic convergence or $\limsup$ moment bounds of the scaled error 
\begin{equation*}
    y_n = \gamma_n^{-1/2}(x_n - x^*)  .
\end{equation*}
Our goal is to establish convergence rates of $y_n$ to its stationary limit in the Wasserstein-$p$ distance by leveraging moment bounds. 
To that effect, we assume the following can be verified.
\begin{assumption}\label{ass:yn_moment}
    For every $k \geq 1$, $y_k$ has finite $2p$ moments.
\end{assumption}
By \Cref{ass:noise_markov} and \citep[Theorem 17.4.2]{meyn2012markov}, there exists solutions $\Phi, \Phi^A$ to the Poisson equations
\begin{equation}\label{eq:poisson}
    f(x^*, \xi) = \Phi(\xi) - \mathbb{E}[\Phi(\xi_{k+1}) | \xi_k = \xi],
    \quad
    A(\xi) - \mathbb{E}A(\xi) = \Phi^A (\xi) - \mathbb{E}[\Phi^A (\xi_{k+1}) | \xi_k = \xi] 
\end{equation}
such that $\lVert \Phi(\xi) \rVert^{2p} \leq K^\Phi (V(\xi) + 1)$ and $\lVert \Phi^A (\xi)  \rVert^{2p} \leq K^{\Phi^A} (V(\xi)+ 1)$ for all $\xi$. 
The Poisson equation is invariant to additive constants to the solution, and we use $\Phi, \Phi^A$ to denote the centered solutions such that these bounds hold. 
From \eqref{eq:poisson}, we define a sequence $\{M_k\}$ as
\begin{equation}\label{eq:Mk}
    -M_k = [\Phi(\xi_{k}) - (P\Phi)(\xi_{k-1})] + W_k ,
\end{equation}
which is a MDS with respect to $\mathcal{H}_{k-1}$ by construction. 
From \Cref{ass:noise_markov,ass:martingale:covariance}, it holds by \citep[Theorem 3]{srikant2024CLT} that $\Gamma_k \coloneqq \mathbb{E}_{k-1} M_k M_k^\dagger$ converges to some $\Gamma \succ 0$ such that $\lVert \Gamma_k - \Gamma\rVert$ is summable, and $\Gamma$ can be evaluated as
\begin{equation*}
    \Gamma = 
    \mathbb{E}_{\xi_1 \sim \pi} \left[ \left(\Phi(\xi_2) - (P \Phi(\xi_{1}) \right) \left(\Phi(\xi_2) - (P \Phi(\xi_{1}) \right)^\dagger \right]
    + \Gamma^W .
\end{equation*}



To establish rates of convergence, we consider the Wasserstein-$p$ metric which metrizes a topology stronger than weak convergence.
Recall that for any $p \geq 1$, the Wasserstein-$p$ distance is defined as 
\begin{equation}\label{eq:wasserstein_def}
    \mathcal{W}_p \left(X, Y \right) = \inf \lVert X - Y \rVert_{L_p} , 
\end{equation}
where the infimum is over all couplings $(X, Y)$ subject to marginal constraints. 
As is standard in finite-dimensional spaces, the norm $\lVert \cdot \rVert$ defined in \eqref{eq:weighted_norm} is equivalent to the Euclidean norm $\lVert \cdot \rVert_2$. 
Consequently, the Wasserstein-$p$ metric \eqref{eq:wasserstein_def} induced by the cost $\lVert X - Y \rVert^p$ is equivalent to the one induced by the standard Euclidean cost $\lVert X - Y \rVert^p_2$. 
Thus, any convergence result established with respect to one metric holds for the other, up to a constant multiplicative factor of $\sqrt{\lVert Q \rVert_2}$ or $\sqrt{\lVert Q^{-1} \rVert_2}$. 
\begin{assumption}\label{ass:CLT}\label{ass:last}
    For any non-negative integers $k_1, k_2$ such that $I = k_2 - k_1  \geq 1$, the MDS \eqref{eq:Mk} satisfies the central limit theorem uniformly in $k_1$: for some positive valued $I \mapsto \clt(I)$ independent of $k_1$,
    \begin{equation*}
        \mathcal{W}_{p} \left( (I \Gamma)^{-1/2} \sum_{k=k_1}^{k_2 - 1} M_k , Z\right) \leq 
        \sqrt{\lVert Q \rVert_2}  \frac{\clt(I)}{\sqrt{I}} \qquad \text{s.t.} \qquad \lim_{I \to \infty} \frac{\clt(I)}{\sqrt{I}} = 0. 
    \end{equation*}
\end{assumption}


\subsection{Discussion on \texorpdfstring{\Crefrange{ass:step_size}{ass:yn_moment}}{Assumptions}}\label{sec:discussion_assumptions}
\Cref{ass:step_size,ass:jacobian,ass:general_sa} are standard.


We incorporate the MDS noise $\{W_k\}$ to ensure our results apply to algorithms with additive (exogeneous) noise, consistent with the framework in \citep{polyakJuditsky}. 
It is instructive to contrast \Cref{ass:noise_martingale,ass:CLT} with the assumptions of \citep{polyakJuditsky}, who establish an asymptotic CLT for the PR average assuming a Lindeberg-type condition and bounds on the second moment ($\sup_{k \geq 1} \mathbb{E}[W_k W_k^\dagger | \mathcal{H}_{k-1}] < \infty$, $\lim_{k \to \infty} \Gamma_k^W = \Gamma^W$ in probability). 


For \nonasymptotic analysis, stronger assumptions are requisite. 
Generally, the CLT in \Cref{ass:CLT} of the driving noise process requires that the first $p + 2$ moments of the noise is finite \citep{rio_lb}. 
The $2p$ moment condition in \Cref{ass:noise_martingale} imposes no additional restriction beyond the $p + 2$ moment condition requirement when convergence in the $\mathcal{W}_1$ or $\mathcal{W}_2$ distances are of interest. 
The condition is an additional regularity property imposed when convergence in the $p > 2$ Wasserstein distance is of interest, and is satisfied when exponential moments exist as assumed in \citep{borkar25,chen_contractive_concentration}. 






The following remarks follow from \Crefrange{ass:noise_markov}{ass:noise_martingale}, which is used to establish our main results.
\begin{enumerate}
    \item \Cref{ass:noise_markov:drift} ensures that the mean of $V(\xi_k)$ is finite for all $k$, where the drift condition implies $\mathbb{E}[V(\xi_{k+1})] \leq \rho_V \mathbb{E}[V(\xi_k)] + C_V$, which further implies
    \begin{equation*}
        \mathbb{E}[V(\xi_k)] \leq \rho_V^{k-1} \mathbb{E}[V(\xi_1)] + \frac{C_V}{1-\rho_V} . 
    \end{equation*}    
    By \Cref{ass:noise_operator_compatibility}, we have $\sup_{k \geq 1} \lVert f(x^*, \xi_k) \rVert_{L_{2p}} \leq \beta_{2p}^f$ and $\sup_{k \geq 1} \lVert A(\xi_k) \rVert_{L_{2p}} \leq \beta_{2p}^A$. 
    
    

        
    \item
    Recall $\lVert \Phi(\xi) \rVert^{2p} \leq K^\Phi (V(\xi) + 1)$ and $\lVert \Phi^A (\xi) \rVert^{2p} \leq K^{\Phi^A} (V(\xi) + 1)$ as discussed after \eqref{eq:poisson}. 
    Since $\mathbb{E}[V(\xi_k)] < \infty$ for all $k$, there exists some positive constants $\beta_{2p}^\Phi$ and $\beta_{2p}^{\Phi^A}$ such that
    \begin{align*}
        \sup_{k \geq 1} \lVert \Phi(\xi_k) \rVert_{L_{2p}} \leq \beta_{2p}^\Phi, 
        \qquad 
        \sup_{k \geq 1} \lVert \Phi^A (\xi_k) \rVert_{L_{2p}} \leq \beta_{2p}^{\Phi^A} . 
    \end{align*}
    This property is used to apply the Cauchy-Schwartz inequality
    \begin{equation}\label{eq:holder}
        \sup_{k \geq 1} \frac{\lVert \Phi^A(\xi_k) y_k \rVert_{L_p}}{\lVert y_k \rVert_{L_{2p}}} \leq \beta_{2p}^{\Phi^A}         .
    \end{equation}
    When the essential supremum $\lVert \Phi^A \rVert_{L_\infty}$ is finite (as in the case of finite state Markov chains), then we may use H\"{o}lder's inequality $\lVert \Phi^A (\xi_k) y_k \rVert_{L_p} \leq \beta_\infty^{\Phi^A} \lVert y_k \rVert_{L_p}$ instead of \eqref{eq:holder}, which is useful in relaxing the $2p$ moment bound for $\{y_k\}$ to a $p$-moment bound. 



\end{enumerate}
The above remarks imply that there exists constants $\rosenthal{M}, \beta_{2p}^M$ associated with the sequence $\{M_k\}$, defined analogously to the constants $\rosenthal{W}, \beta_{2p}^W$ in \Cref{ass:noise_markov}. 
Similarly, we use $\beta_p^{\Gamma^{-1/2}M}$ to denote the $p$--th moment bound for the normalized sequence $\{\Gamma^{-1/2}M_k\}$. 








We now discuss \Cref{ass:yn_moment}.
The analysis of stochastic approximation algorithms has progressed from establishing almost sure convergence of the unnormalized sequence $\{x_k\}$ to providing rates of convergence. 
It is now a well-established result \citep{gerencser1992rate,kushner_yin,borkar25} that under various assumptions on the noise process and function $f$, the asymptotic rate
\begin{equation}\label{eq:yn_bounds_Lp}
    \limsup_{k \to \infty} \frac{\lVert y_k \rVert_{L_{2p}}}{\lVert \Sigma_a^{1/2} Z \rVert_{L_{2p}}} < \infty
\end{equation}
holds for some $p \geq 1$. 
\Cref{ass:yn_moment} requires that this asymptotic rate of convergence, that $\{x_k\}$ converges to $x^*$ in the $L_p$ norm at rate $\sqrt{\gamma_n}$, is known.



When \eqref{eq:yn_bounds_Lp} holds, one can verify additional conditions to establish the distributional limit
\begin{equation}\label{eq:yn_asymptotic_clt}
    y_n \to \mathcal{N}(0, \Sigma_a), \qquad \bar{A}_a \Sigma_a + \Sigma_a \bar{A}_a^\dagger = \Gamma .
\end{equation}
While \eqref{eq:yn_asymptotic_clt} suggests $\lVert y_n \rVert_{L_{2p}} \to \lVert \Sigma_a^{1/2} Z \rVert_{L_{2p}}$ for some values of $p \geq 1$, upper bounds of the form \eqref{eq:yn_bounds_Lp} are not strong enough to obtain this conclusion. 
One contribution of our work is to show that, under standard conditions for which \Cref{ass:yn_moment} is often verified, convergence in $\mathcal{W}_p$ can be established and the corresponding rate can be obtained. 



We remark that when $\bar{f}$ is a linear map and $\lVert \Phi^A (\xi_k) \rVert_{L_\infty}$ is finite for all $k$, then \Cref{ass:yn_moment} can be relaxed and the results in \Cref{sec:results} can be established with the assumption that $\{y_k\}$ is bounded in $L_p$.











\subsection{Non-asymptotic Central Limit Theorems}\label{sec:nonasymptotic_CLT}
\Cref{ass:CLT} encapsulates the requirement that the noise partial sums converge to a Gaussian limit in $\mathcal{W}_p$. 
We state this as a distinct condition because the implication from the moment bounds (\Crefrange{ass:noise_markov}{ass:noise_martingale}) to the central limit theorem depends on the metric's order $p$.
For $p=1$, existing results \citep{srikant2024CLT} establish that \Crefrange{ass:noise_markov}{ass:noise_martingale} imply \Cref{ass:CLT}, provided the noise possesses finite 3rd moments. 
In this regime, \Cref{ass:CLT} is a consequence of the primitive conditions. 
For $p > 1$, a general martingale CLT in $\mathcal{W}_p$ has not yet been established.
Therefore, \Cref{ass:CLT} serves as a necessary structural hypothesis, which we verify for specific cases such as independent noise.
Concrete instances of this hypothesis, covering both the general martingale case for $p=1$ and independent noise for $p \ge 2$, are established in \Cref{lem:martingale_clt,lem:bonis_weighted} below.

    




For $n \geq 2$, consider a sequence of intervals $\{[k_m, k_{m+1})\}_{m=1}^{N}$ that form a subset of $\{1, \cdots, n-1\}$ with $k_{N+1} = n$.
The length of each interval $[k_m, k_{m+1})$ is denoted by $I_m = k_{m+1} - k_m$, where our results are presented assuming that the interval lengths are increasing in size such that $\gamma_{k_m}I_m \to 0$. 
The cumulative \stepsize in the interval $[k_m, k_{m+1})$ is denoted by $\cumstep_m = \sum_{k=k_m}^{k_{m+1} - 1} \gamma_k$.
Consider the weighted and unweighted sums
\begin{equation}\label{eq:SmZm}
    \hat{Z}_m = (\cumstep_m \Gamma)^{-1/2} \sum_{k=k_m}^{k_{m+1} - 1} \sqrt{\gamma_k} M_k , 
    \qquad 
    S_m = (I_m \Gamma)^{-1/2} \sum_{k=k_m}^{k_{m+1} - 1} M_k . 
\end{equation}
Under appropriate moment conditions on the MDS sequence $\{M_k\}$ and convergence of the conditional covariance $\Gamma_k \to \Gamma$, the sum $S_m$ weakly converges to a Gaussian limit by the martingale CLT \citep{hallp1980martingalelimittheoryanditsapplication}.
In \Cref{ass:CLT}, we assume that the rate of convergence is known. 




A bound on the Wasserstein-$p$ distance between $\hat{Z}_m$ and an appropriate Gaussian distribution is established using a coupling argument and \Cref{ass:CLT}. 
For the following, we use $C_p^{BR}$ to be a constant that depends on $p$, which arises in the martingale Rosenthal inequality as discussed in \Cref{sec:BDG}. 
\begin{lemma}\label{lem:martingale_clt_weighted}
   Let $\{M_k\}$ be a martingale difference sequence associated with $\rosenthal{M}$. 
   If $\mathcal{W}_{p} (S_m, Z)$ is finite, then it holds that
   \begin{equation}\label{eq:martingale_clt_weighted}
        \mathcal{W}_{p} (\hat{Z}_m, Z) 
        \leq 
        C_{p}^{BR} \rosenthal{\Gamma^{-1/2} M} \frac{aI_m}{k_m} 
        +
        \mathcal{W}_{p} (S_m, Z) . 
   \end{equation}
\end{lemma}
\begin{proof}
    See \Cref{proof:martingale_clt_weighted}.
\end{proof}

A bound on $\mathcal{W}_1(S_m, Z)$ is obtained using \citep[Theorem 1]{srikant2024CLT}, and we use \Cref{lem:martingale_clt_weighted} to obtain a bound on $\mathcal{W}_1(\hat{Z}_m, Z)$.
The proofs of the remaining lemmas in this section are provided in \Cref{app:clt_weighted}.
\begin{lemma}\label{lem:martingale_clt}
    Let $\{M_k\}$ be a MDS defined in \eqref{eq:Mk}.
    If $\sup_{k \geq 1} \lVert \Gamma^{-1/2} M_k \rVert_{L_3} \leq \beta_3^{\Gamma^{-1/2} M}$, then for some $C_1 > 0$ and sufficiently large $k_m$,
    \begin{equation*}
    \mathcal{W}_1 (\hat{Z}_m, Z) 
        \leq
    C_1 
    \left(\kappa_2^{\Gamma^{-1/2}M} \frac{a I_m}{k_m} 
    +
    \left(\beta_3^{\Gamma^{-1/2} M} \right)^3 \sqrt{\lVert Q \rVert_2 \lVert \Sigma_a \rVert_2} \frac{d \log I_m}{\sqrt{I_m}}
    \right) .
    \end{equation*}
\end{lemma}
The result holds for any initial distribution, where we used the convergence $\mathbb{E} \Gamma_k \to \Gamma$ by \Cref{ass:noise_markov,ass:noise_martingale}.
When $\{M_k\}$ is a sequence of independent random variables, the central limit theorem can be established for $p \geq 2$ Wasserstein distances \citep{bonis2020steinsmethodnormalapproximation}. 
Using \Cref{lem:martingale_clt_weighted}, we establish the corresponding CLT for the weighted sum $\hat{Z}_m$.
\begin{lemma}\label{lem:bonis_weighted}
    Let $p \geq 2$ and assume $\{M_k\}$ is a mean zero i.i.d. sequence with $\lVert M_k \rVert_{L_{p+2}} \leq \beta_{p+2}^M$. 
    There exists a constant $C_p$ such that 
    \begin{equation*}
        \mathcal{W}_p (\hat{Z}_m, Z) 
        \leq 
         C_p \left(\rosenthal{\Gamma^{-1/2}M}  \frac{a I_m}{k_m} 
+
    \sqrt{\lVert Q \rVert_2}\frac{d^{1/4} (\beta_4^{\Gamma^{-1/2}M})^2 + (\beta_{p+2}^{\Gamma^{-1/2}M})^{1+2/p} }{\sqrt{I_m}}
     \right)
. 
    \end{equation*}
\end{lemma}
The constant $C_p$ in \Cref{lem:bonis_weighted} shares the same asymptotic growth in $p$ as the constant $C_p^{BR}$. 





\Cref{lem:martingale_clt,lem:bonis_weighted} provide concrete settings where \Cref{ass:CLT} is satisfied. When \Cref{ass:CLT} holds, \Cref{lem:martingale_clt_weighted} implies that for any integer sequence $\{I_m\}_{m=1}^N$:
\begin{equation}\label{eq:clt_weighted_master}
    \mathcal{W}_p (\hat{Z}_m, Z) \leq  C_p^{BR} \rosenthal{\Gamma^{-1/2} M} \cdot a \frac{I_m}{k_m} + \frac{\clt(I_m)}{\sqrt{I_m}}.
\end{equation}
While we utilize the CLTs in \Cref{lem:martingale_clt,lem:bonis_weighted} to demonstrate the generality of our framework, Assumption 8 is modular. If the sequence $\{M_k\}$ possesses additional structure, such as independence and a Poincaré inequality, the sharper bounds established in \citep{courtade2018existencesteinkernelsspectral,bonis2024improved} can be used to verify \Cref{ass:CLT} with improved dependencies on the dimension $d$ and order $p$. Consequently, applying our main theorems in these settings simply requires checking \Crefrange{ass:step_size}{ass:yn_moment} alongside the specific conditions of the relevant CLT result.










\subsection{Related Work}
Our work is at the intersection of three major lines of inquiry in stochastic approximation theory: the asymptotic theory of weak convergence, finite-sample moment bounds, and non-asymptotic central limit theorems. 



A central goal of SA theory is to characterize the full distribution of the iterates beyond their almost sure convergence.
The idea of using a diffusion approximation to connect the discrete-time recursion to a limiting Ornstein-Uhlenbeck (OU) process has a rich history, particularly in establishing weak convergence via functional central limit theorems (FCLTs) when the driving noise is a sequence of independent random variables \citep{pezeshki1997strong,joslin2000law,bucklew2002weak}.
The FCLT for stochastic approximation was recently extended to settings where the driving noise is Markovian \citep{borkar25}.
Such convergence results provide a much more complete picture of the dynamics, guaranteeing that the sample path fluctuations can be approximated by a stochastic differential equation (SDE). 
While these results focus on the asymptotic behavior of the sample path, our work leverages this connection to analyze the convergence rate of the last iterate. 
By providing convergence rates in the Wasserstein-$p$ distances, we are able to obtain tail probability bounds for finite $n$. 







In response to the need for finite-sample guarantees, a parallel line of research has focused on deriving explicit moment bounds. 
This approach has been highly successful for stochastic gradient descent \citep{moulines_bach,bach2013non} and has been extended to challenging settings with dependent noise, including Markovian dynamics \citep{srikant19,chen2022finitesampleanalysisnonlinearstochastic,thinhMarkovGD}. 
While this body of work provides invaluable concrete performance bounds, its guarantees are distinct from distributional convergence.
Moment bounds concern specific statistics of the error, and the analyses typically culminate in $\limsup$ bounds on the moments of $y_n$ but not the convergence of the moments. 


\Nonasymptotic CLTs for stochastic approximation have emerged to obtain sharper tail bounds than what can be obtained using Markov's inequality and moment bounds, or convergence of the moments. 
For the last iterates of linear SA, \citet{butyrin2025gaussian} established a finite-time bound in a convex distance which metrizes a topology weaker than the $\mathcal{W}_1$ distance \citep{nourdin2021multivariatenormalapproximationwiener}.
For \nonlinear SA with additive i.i.d. noise, \citet{sheshukova2025gaussianapproximationmultiplierbootstrap} established finite-time bounds for the PR average in the same convex distance.
For linear SA with more general noise structures, convergence rates have been established in the stronger $\mathcal{W}_1$ distance: \citep{srikant2024CLT} for multiplicative Markov noise and \citep{kong2025nonasymptoticclterrorbounds} for (additive) MDS noise.


We develop a \nonasymptotic theory of weak convergence for general nonlinear SA, for both the last iterate and its PR average, established in the strong $\mathcal{W}_p$ metric.
Our approach builds upon the diffusion approximation central to the FCLT, which connects the discrete recursion to a limiting SDE. 
However, rather than establishing weak convergence of the entire sample path in the limit, we utilize the coupling to derive explicit finite-sample bounds in the Wasserstein-$p$ metric. 
This converts the asymptotic qualitative insights of the diffusion framework into the algorithm's performance at a fixed iteration $n$. 



\section{Main Results}\label{sec:results}
Our main result is a general, \nonasymptotic bound on the Wasserstein-$p$ distance between the SA output and its stationary distribution. 


\subsection{Last Iterate Analysis}\label{sec:general}
We show in \Cref{app:error_recursion} that the sequence $\{y_k\}$ satisfies the recursion
\begin{equation}\label{eq:yk}
\begin{split}
    y_{k+1} &= y_k - \gamma_k \bar{A}_a y_k + \sqrt{\gamma_k} M_k + R (\gamma_k, y_k, \xi_k, W_k)  ,
    \qquad 
    y_1 = \frac{x_1 - x^*}{\sqrt{\gamma_1}} ,
\end{split}
\end{equation}
where $M_k$ was defined in \eqref{eq:Mk} and $R$, whose expression is deferred to \eqref{eq:Rk}, encapsulates higher order terms depending on the structure assumed for $f$. 
Neglecting the higher order term, the recursion \eqref{eq:yk} resembles a discretized-form of the Ornstein-Uhlenbeck (OU) process. 








To view the sequence \eqref{eq:yk} as a discretized OU process, one must consider intervals of increasing length \citep{bertsekas_ndp}. 
To this end, we consider for $n \geq 1$ a sequence $\{k_m\}_{m=1}^{N+1}$ that partitions $\{1, \cdots, n\}$ with intervals of length $I_m = k_{m+1} - k_m$ as discussed in \Cref{sec:nonasymptotic_CLT}. 
A sub-sequence $\{y_{k_m}\}$ is obtained as a telescopic sum of the sequence \eqref{eq:yk} over the interval $[k_m, k_{m+1}]$:
\begin{equation}\label{eq:ytildem}
    y_{k_{m+1}} 
    = y_{k_m} - \cumstep_m \bar{A}_a y_{k_m} + \sqrt{\gamma_k \Gamma} \hat{Z}_m + \tilde{R}_m ,     
\end{equation}
where $\cumstep_m = \sum_{k=k_m}^{k_{m+1} - 1} \gamma_k$ is the cumulative \stepsize.
We show that the higher order term is negligible when $I_m$ is chosen so that $\cumstep_m \approx \gamma_{k_m} I_m \to 0$.


Under \Cref{ass:yn_moment}, we use $\{\beta_p^y (m)\}$ and $\{\beta_{2p}^y (m)\} $ to denote sequences such that
\begin{equation}\label{eq:yn_interval_bound}
    \max_{k \in [k_m, k_{m+1})} \lVert y_k \rVert_{L_p}
    \leq \beta_p^y (m) \lVert \Sigma_a^{1/2} Z \rVert_{L_p}    
    \quad 
    \text{and} \quad
    \max_{k \in [k_m, k_{m+1})} \lVert y_k \rVert_{L_{2p}} \leq  
    \beta_{2p}^y (m) \lVert \Sigma_a^{1/2} Z \rVert_{L_{2p}}
    .
\end{equation}
For $a \in (0, 1]$ in \Cref{ass:step_size}, define the sequence $\{E_k^a\}$ as
\begin{equation}\label{eq:Eka_def}
     E_k^a = \left\lVert \sqrt{\frac{\gamma_k}{\gamma_{k+1}}} \left( \identity - \gamma_k \nabla \bar{f}(x^*) \right) - \left(\identity - \gamma_k \bar{A}_a \right) \right\rVert ,
\end{equation}
where $E_k^a = \mathcal{O}(1/k)$ when $a < 1$ and $E_k^a = \mathcal{O}(1/k^2)$ when $a = 1$.
\begin{lemma}\label{lem:general_remainder}
    Suppose \Crefrange{ass:step_size}{ass:yn_moment} hold, and assume that the sequence $\{\beta_{2p}^y(m)\}$ is bounded.
    For any increasing sequence $\{I_m\}$ that satisfies $I_m \gamma_{k_m} \to 0$,
    \begin{equation*}
    \begin{split}
    \Lp{\tilde{R}_m} 
    &\leq \lVert \Sigma_a^{1/2}Z \rVert_{L_p} \cdot I_m E_{k_m}^a \beta_p^y (m) + 4 \beta_p^{\Phi} \sqrt{\gamma_{k_m}}
    + 
    C_p^{BR} \rosenthal{M} \lVert \bar{A}_a \rVert \cumstep_m^{3/2}  \\ 
    &
    + \mathcal{O}\left(\cumstep_{m} I_m E_{k_m}^a 
    + \sqrt{\gamma_{k_m}} \frac{I_m}{k_m} + \sqrt{\gamma_{k_m} \cumstep_m}
    + \cumstep_m^2
    \right)
    .
    \end{split}
    \end{equation*}
\end{lemma}
\begin{proof}
    See \Cref{app:general_remainder}.
\end{proof}
The sequence $\{y_{k_m}\}$ will be compared to the Euler-Maruyama discretization $\{u_m\}_{m=1}^{N+1}$ of an OU process $(U_t)$, described by the equations
\begin{equation*}
\begin{split}
    u_{m+1} &= u_m - \cumstep_m \bar{A}_a u_m + \sqrt{\cumstep_m \Gamma} Z_m, \qquad Z_m \sim \mathcal{N}(0, \mathbb{I})     
    \\ 
    dU_t &= -\bar{A}_a U_t dt + \sqrt{\Gamma} dB_t ,
\end{split}
\end{equation*}
where $(B_t)$ is the standard Brownian motion. 
Since $I_m$ is chosen such that $\cumstep_m \to 0$, the discretized OU process converges to the same limit as the continuous process. 
For the following, we use $U_\infty$ to denote the stationary limit of the OU process.  
The Burkholder-Davis-Gundy (BDG) inequality \citep[Theorem 5.16]{le2016brownian} is used to obtain a bound on the Wasserstein-$p$ distance between the discretized and continuous time OU processes, and we use $C_p^{BDG}$ to denote the constant that arises in the inequality as discussed in \Cref{app:ou_discretization}.



The following uses \Cref{lem:general_remainder} to deduce the bound $\Lp{\tilde{R}_m} = o(\epsilon_m)$, which holds when $I_m$ is such that $I_m \sqrt{\gamma_{k_m}} \to 0$.
\begin{lemma}\label{lem:Wp_recursion}
    Suppose \Crefrange{ass:step_size}{ass:last} hold, and assume $\{I_m\}$ is such that $I_m \sqrt{\gamma_{k_m}} \to 0$. 
    Under the assumptions of \Cref{lem:general_remainder}, the sequence \eqref{eq:ytildem} satisfies  
    \begin{equation}\label{eq:Wp_recursion}
    \begin{split}
    \mathcal{W}_{p} (y_{n}, U_\infty)
    &\leq
    K e^{-\lambdaCT \sum_{m=1}^N \cumstep_m} \mathcal{W}_{p} \left(y_1, U_\infty \right) 
    + \frac{K C_p^{BDG}}{\lambdaCT} \frac{\lVert \bar{A}_a \rVert}{\lambdaDT} \sqrt{\frac{\lVert Q \rVert_2\mathrm{Tr} \Gamma}{3}}\sqrt{\cumstep_N}
    \\ &
    + \frac{1}{\lambdaDT \cumstep_N} \left(
    \sqrt{\lVert \Gamma \rVert \cumstep_N} \left(C_{p}^{BR} \rosenthal{M} \cdot a \frac{I_N}{k_N} +  \frac{\clt(I_N)}{\sqrt{I_N}} \right)
    + \lVert \tilde{R}_N \rVert_{L_p}
    \right)    
    \\ &
    + \mathcal{O}\left(\cumstep_N 
    +
    \frac{1}{\cumstep_N^2 k_N} \left( \sqrt{\cumstep_N} \left(\frac{I_N}{k_N} + \frac{ \clt (I_N)}{\sqrt{I_N}} \right) + \lVert \tilde{R}_N \rVert_{L_p} \right) 
    \right)
    .
    \end{split}
    \end{equation}
\end{lemma}
\begin{proof}
    See \Cref{app:Wp_recursion}.
\end{proof}
We now present our first convergence result.
While the error bound in \Cref{lem:Wp_recursion} holds for general noise sequences that satisfy \Cref{ass:CLT}, the following theorem instantiates this bound for the concrete noise settings established in \Cref{lem:martingale_clt,lem:bonis_weighted}. 
\begin{theorem}\label{thm:general}
    Under the assumptions of \Cref{lem:general_remainder}, the following convergence rates hold depending on the structure of the noise sequence $\{M_k\}$:

    \begin{enumerate}
        \item General Noise: If $\{M_k\}$ satisfies the conditions of \Cref{lem:martingale_clt}, then
        \begin{equation}\label{eq:W1_main}
        \begin{split}
            \mathcal{W}_1 (y_n, U_\infty) 
            =
            \mathcal{O}\left(\gamma_n^{1/6} \left(\log \gamma_n \right)^{1/3}
            + \frac{E_n^a}{\gamma_n} \right) 
            .
        \end{split}
        \end{equation}

        \item Independent Noise Sequence: If $\{M_k\}$ is an i.i.d. sequence satisfying \Cref{lem:bonis_weighted}, then for all $p \geq 2$,
        \begin{equation}\label{eq:Wp_main}
        \mathcal{W}_p (y_n, U_\infty)  
        \leq
        \Upsilon_{p,d} \cdot \Xi_{p, d}\cdot \gamma_n^{1/6}
        +
        \frac{\beta_p^y (N) \lVert \Sigma_a^{1/2}Z\rVert_{L_p}}{\lambdaDT}\cdot \frac{E_{n}^a}{\gamma_n}
        + 
        \mathcal{O}\left( \gamma_n^{1/3} \right), 
        \end{equation}        
        where $\Upsilon_{p,d}$ and $\Xi_{p,d}$ are defined as
        \begin{equation*}
        \begin{split}
    \Upsilon_{p,d} &= 
     \frac{2}{\lambdaDT} \left(\lVert \bar{A}_a \rVert \left(K \frac{C_p^{BDG}}{ \lambdaCT} \sqrt{\frac{\lVert Q \rVert_2\mathrm{Tr} \Gamma}{3}}  + C_p^{BR} \rosenthal{M} \right) \right)^{2/3}   
\\
        \Xi_{p,d} &=
        \left(  C_p \sqrt{\lVert Q \rVert_2 \lVert \Gamma \rVert}
        \left(d^{1/4} (\beta_4^{\Gamma^{-1/2}M})^2 + (\beta_{p+2}^{\Gamma^{-1/2}M})^{1+2/p} \right)  
        + 4 \beta_p^f\right)^{1/3}         
        .
        \end{split}       
        \end{equation*}

    \end{enumerate}    
\end{theorem}
\begin{proof}
    See \cref{app:general}.
\end{proof}



The term $E_n^a/\gamma_n$ dominates the convergence rates in \Cref{thm:general} when $a \in (6/7, 1)$. 
Its presence is necessary because convergence in $\mathcal{W}_p$ for $p \geq 2$ implies the convergence of the first two moments, and the covariance convergence rate is determined by this term. 
To see this, consider the simplest setting of linear stochastic approximation with additive i.i.d. noise: $x_{k+1} = x_k - \gamma_k A x_k + \gamma_k W_k$ with $a < 1$.
The recursion for the covariance $\mathbb{E}[x_k x_k^\dagger]$ satisfies
\begin{align*}
    \mathbb{E} x_{k+1} x_{k+1}^\dagger &= \mathbb{E}x_k x_k^\dagger - \gamma_k \left(A \mathbb{E} x_k x_k^\dagger + \mathbb{E}x_k x_k^\dagger A^\dagger - \gamma_k \Gamma^W \right)
    + \gamma_k^2 A \mathbb{E} x_k x_k^\dagger A^\dagger
    .
\end{align*}
Renormalizing the covariance update with $\Sigma_k^y = \gamma_k^{-1} \mathbb{E}x_k x_k^\dagger$, we obtain
\begin{align*}
    \gamma_{k+1} \Sigma_{k+1}^y &= \gamma_k \Sigma_k^y - \gamma_k^2 \left(A \Sigma_k^y + \Sigma_k^y A^\dagger - \Gamma^W \right) + \gamma_k^3 A \Sigma_k^y A^\dagger 
    \\
    \Rightarrow 
    \gamma_{k+1}\left( \Sigma_{k+1}^y  - \Sigma_a\right) &= \gamma_k (\Sigma_k^y - \Sigma_a)  
    - \gamma_k^2 \left(A (\Sigma_k^y - \Sigma_a) + (\Sigma_k^y - \Sigma_a) A^\dagger - \Gamma^W\right)
    \\ &
    + (\gamma_k - \gamma_{k+1}) \Sigma_a
    + \gamma_k^3 A \Sigma_k^y A^\dagger 
    . 
\end{align*}
Using that $\gamma_k - \gamma_{k+1} = a \gamma_k k^{-1} + \mathcal{O}(\gamma_k k^{-2})$, we obtain
\begin{align*}
    \Sigma_{k+1}^y - \Sigma_a &= 
    \Sigma_k^y - \Sigma_a - \gamma_k (A (\Sigma_k^y - \Sigma_a) + (\Sigma_k^y - \Sigma_a) A^\dagger - \Gamma^W)
    +
    \mathcal{O}\left(\frac{1}{k} + \frac{\gamma_k}{k} + \gamma_k^2 \right) 
    ,
\end{align*}
which can be rearranged to obtain
\begin{align*}
    \frac{(\Sigma_{k+1}^y - \Sigma_a) - (\Sigma_k^y - \Sigma_a)}{\gamma_k}
    &= - (A (\Sigma_k^y - \Sigma_a) + (\Sigma_k^y - \Sigma_a) A^\dagger - \Gamma^W)
    + \mathcal{O}\left(\frac{1}{k \gamma_k}+ \frac{1}{k} + \gamma_k\right) . 
\end{align*}
When the sequence $\{\Sigma^y_k\}$ converges, the rate of convergence is $\lVert \Sigma_n^y - \Sigma_a \rVert = \mathcal{O}(1/n\gamma_n + \gamma_n)$.
A similar calculation shows that the rate of convergence is $\mathcal{O}(1/n^2 \gamma_n)$ when $a = 1$ and $\gamma_1$ is sufficiently large. 
The recursion for the scaled covariance $\Sigma_k^y$ reveals that the convergence of $\Sigma_k^y$ to the stationary covariance $\Sigma_a$ occurs at the rate $\mathcal{O}(E_n^a/\gamma_n)$. 
Since the $\mathcal{W}_p$ metric for $p \geq 2$ controls the convergence of the second moment, our bound in \Cref{thm:general} must reflect this intrinsic rate limit. 




































\subsection{Polyak-Ruppert Averaging}\label{sec:general_pr}
Polyak-Ruppert averaging is used with \stepsize exponent $a < 1$ such that $\bar{A}_a$ in \eqref{eq:hurwtiz} is $\nabla \bar{f}(x^*)$. 
We therefore assume $a < 1$ and use the notation $\bar{A}$ instead of $\bar{A}_a$, since it is a constant independent of $a$ in this range. 


A central limit theorem for the Polyak-Ruppert (PR) average $\bar{x}_n = n^{-1} \sum_{k=1}^n x_k$ was initially obtained by \citet{polyakJuditsky}, where it was shown to achieve the optimal asymptotic mean square convergence
\begin{equation*}
    \lim_{n \to \infty} \left\{n \mathbb{E} (\bar{x}_n - x^*) (\bar{x}_n - x^*)^\dagger - \bar{\Sigma} \right\}= 0 ,
    \qquad \bar{\Sigma} = \bar{A}^{-1} \Gamma \bar{A}^{-\dagger} . 
\end{equation*}
The optimality refers to both the $n^{-1}$ rate in the bound of $\mathbb{E}\lVert \bar{x}_n - x^* \rVert^2$ and the covariance $\bar{\Sigma}$, where $\mathrm{Tr} \bar{\Sigma} \leq \mathrm{Tr} \Sigma_a \leq \mathrm{Tr} \Sigma_1$ for all $a \in (0, 1]$ and $\Sigma_1$ is the asymptotic covariance \eqref{eq:yn_asymptotic_clt} corresponding to $a = 1$. 
The following is a \nonasymptotic CLT for the PR average, where we use $\bar{y}_n = \sqrt{n}(\bar{x}_n - x^*)$ to denote the scaled error of the PR average. 
\begin{theorem}\label{thm:pr_average}
    Under \Cref{ass:step_size} with $a \in (1/2, 1)$ and \Crefrange{ass:jacobian}{ass:last} with $\{\beta_{2p}^y\}$ bounded,
    \begin{equation}\label{eq:pr_average_thm}
    \begin{split}
    \mathcal{W}_{p} \left( \bar{A} \bar{y}_n, \Gamma^{1/2}Z \right) 
    &\leq 
    \mathcal{W}_{p} \left(\frac{1}{\sqrt{n}} \sum_{k=1}^n M_k, \Gamma^{1/2}Z \right)
    \\ 
    &+    
    \frac{\max_{k \leq n+1} \lVert y_k \rVert_{L_p}}{\sqrt{\gamma_n n}}
    +
    \frac{\lip(\nabla \bar{f}) (\max_{k \leq n} \lVert y_k \rVert^2_{L_{2p}})}{1-a} n^{1/2} \gamma_n
    \\ 
    &
    + \beta_{2p}^{\Phi^A}
    \left(\frac{2 \max_{k \leq {n+1}} \lVert y_k \rVert_{L_{2p}} }{a \sqrt{n \gamma_n}} + (\beta_{2p}^f + \beta_{2p}^W) \frac{\gamma_n}{1-a} n^{1/2} \right)
    \\ 
    &+
    \frac{\lip(\nabla f) + \lip(\nabla \bar{f}) }{1-a} \left(\max_{k \leq n} \lVert y_k \rVert^2_{L_{2p}}\right) n^{1/2} \gamma_n
    \\ &
    + \mathcal{O}\left(\sqrt{\gamma_n}\right) .             
    \end{split}
    \end{equation}
\end{theorem}
\begin{proof}
    See \cref{app:pr_average}.
\end{proof}
The first term $\mathcal{W}_{p}(n^{-1/2}\sum_{k=1}^n M_k, \Gamma^{1/2}Z)$ concerns the convergence of a martingale sum, and does not depend on the step size sequence. 
The requirement $a > 1/2$ is imposed so that $n^{1/2} \gamma_n \to 0$ and $\bar{A} \bar{y}_n \to \Gamma^{1/2} Z$ in distribution. 
A step size exponent $a = 2/3$ optimizes the remaining error rates (ignoring constants that depend on $a$), which matches the suggestion in \citep{moulines_bach,srikant2024CLT} and leads to the convergence rate $n^{-1/6}$.
On the other hand, this differs from the $a = 3/4$ suggested by \citet{sheshukova2025gaussianapproximationmultiplierbootstrap} based on an analysis of the convex distance. 
This is a reasonable discrepancy given that the analyses are based on different metrics; since the $\mathcal{W}_p$ metric for any $p \geq 1$ yields a bound on a more diverse class of test functions \citep{nourdin2021multivariatenormalapproximationwiener}, the choice $a = 2/3$ may be needed to achieve the $n^{-1/6}$ rate uniformly over the larger class of test functions. 


The nature of the error bound is significantly different from that in \Cref{thm:general}.
Furthermore, the bound in \Cref{thm:general} depends on the CLT rate of the weighted sum $\hat{Z}_N$ while the bound in \Cref{thm:pr_average} depends on the CLT for the unweighted sum $S_N$. 
Consequently, the dependence of the bound in \Cref{thm:pr_average} on problem constants and $p$ is generally different from that of the bound in \Cref{thm:general}.






\subsection{High Probability Bounds}
A practical application of \nonasymptotic analysis is the construction of confidence intervals or high-probability bounds on the error $\lVert x_k - x^* \rVert$. 
Two methods are commonly used.


\textbf{Moment-Based Bounds:}
In proving moment bounds of the form in \Cref{ass:yn_moment}, one often shows that $\lVert y_n \rVert_{L_p}$ is within a constant multiple of its limit's, i.e., $\lVert y_n \rVert_{L_p} \leq \beta_p^y  \lVert \Sigma^{1/2}Z \rVert_{L_p} + o(1)$ for some time-independent constant $\beta_p^y$.
One can apply Markov's inequality to deduce 
\begin{equation}\label{eq:markov_last}
\begin{split}
    \lVert x_n - x^* \rVert_{L_p}
    &\leq
    \gamma_n^{1/2} \left( \beta_p^y \lVert \Sigma^{1/2} Z \rVert_{L_p} + o(1)\right)  \confint^{-1/p} 
    , 
    \qquad \text{w.p. } \geq 1 - \confint
    .     
\end{split}
\end{equation}
While rigorous, this bound fails to capture the sub-Gaussian nature of the error distribution. 
Any bounds of the form \eqref{eq:markov_last} with a constant $\beta_p^y > 1$ cannot be used to prove convergence. 
The convergence of the $p$ moments implied by Theorem \ref{thm:general} already improves the above estimate by removing the $\beta_p^y$ factor.


We are not aware of finite-time $p > 2$ moment bounds on the PR average $ \bar{x}_n - x^*$, so we discuss the implication of Markov's inequality in combination with Theorem \ref{thm:pr_average}.
Using the implied $p$ moment convergence $\lVert \bar{x}_n - x^* \rVert_{L_p} \leq n^{-1/2} \lVert \bar{\Sigma}^{1/2}Z \rVert_{L_p} + o(n^{-1/2})$, the confidence interval inferred using Markov's inequality is
\begin{equation}\label{eq:markov_pr}
    \lVert \bar{x}_n - x^* \rVert \leq \frac{1}{\sqrt{n}} \left(\lVert \bar{\Sigma}^{1/2} Z \rVert_{L_p}  + o(1)\right)\confint^{-1/p},
    \qquad 
    \text{w.p. }
    \geq 1 - \confint
    .  
\end{equation}




\textbf{Gaussian Approximation:} 
Confidence intervals are often estimated assuming the asymptotic Gaussian approximation $y_n \approx \Sigma_a^{1/2} Z$, which yields estimates of the form
\begin{equation}\label{eq:gaussian_tail_estimates}
\begin{split}
    \lVert x_n - x^* \rVert 
    &\lesssim 
    \gamma_n^{1/2} \left(\mathbb{E}\lVert \Sigma_a^{1/2} Z \rVert 
    + 
   \sqrt{2  \lVert \Sigma_a \rVert \log \frac{2}{\confint}}     \right)
   \\ 
   \lVert \bar{x}_n - x^* \rVert 
   &\lesssim 
   \frac{1}{\sqrt{n}} \left(\mathbb{E}\lVert \bar{\Sigma}^{1/2}Z \rVert + \sqrt{2 \lVert \bar{\Sigma} \rVert \log \frac{2}{\confint}} \right) 
\end{split}
\end{equation}
for some confidence tolerance $1 - \confint \in (0, 1)$, where $\lesssim$ is used to describe that the inequalities are based on approximations. 
While asymptotically valid, this approach can be overly optimistic in finite-time regimes. 
Specifically, standard Gaussian approximations fail to capture the transient, heavy-tailed behavior often exhibited by stochastic approximation algorithms with multiplicative noise (e.g., linear SA). 
By relying solely on the limit distribution, these approximations neglect the convergence rate of the tail probabilities, potentially leading to significant underestimation of risk in early iterations. 
Our Wasserstein bounds address this by explicitly quantifying the transition from this transient phase to the asymptotic Gaussian regime, as demonstrated for linear SA in \Cref{sec:lsa}. 



For linear SA, high probability bounds can be obtained as in \citep[Theorem 1]{durmus2021tighthighprobabilitybounds}. 
However, this utilizes a structure specific to linear SA and the technique can be hard to generalize for nonlinear SA. 
Moreover, the confidence interval can still be significantly larger than the estimate obtained using an asymptotic Gaussian approximation. 




High probability bounds tighter than those in \eqref{eq:markov_last} and \eqref{eq:markov_pr} are implied by \Cref{thm:general,thm:pr_average}, using an argument similar to those discussed in \citep{fang2023moderatedeviations,liu2023wassersteinpboundscentrallimit}. 
For the following result, let $C_{p,1}, C_{p,2} > 0$ be constants such that $\mathcal{W}_p (y_n, \Sigma_a^{1/2} Z) \leq C_{p,1} \gamma_n^{1/6} + C_{p,2} E_n^a/\gamma_n$ under the assumptions of \Cref{thm:general}.
Denote by $\bar{C}_p$ a constant such that $\mathcal{W}_p (\bar{y}_n, (A^{-1} \Gamma A^{-\dagger})^{1/2}Z) \leq \bar{C}_p n^{-1/6}$  which arises in \Cref{thm:pr_average} when the \stepsize is chosen with exponent $a = 2/3$. 
\begin{corollary}\label{cor:high_probability}
    Assume \Crefrange{ass:noise_markov}{ass:CLT} hold for some $p \geq 1$.
    For any $\confint \in (0, 1)$ and sufficiently large $n$, the following error bounds hold with probability $\geq 1 - \confint$.

    Under the assumptions of \Cref{thm:general}, the last iterate achieves the error 
    \begin{equation*}
        \lVert x_n - x^* \rVert 
        \leq 
        \gamma_n^{1/2}\left( \mathbb{E}\lVert \Sigma_a^{1/2} Z \rVert 
        + \sqrt{2  \lVert \Sigma_a \rVert \log \frac{2}{\confint}} \right) + \left(C_{p, 1} \cdot \gamma_n^{2/3} + C_{p, 2} \cdot \gamma_n^{-1/2} E_{n}^a\right) \left( \frac{2}{\confint} \right)^{1/p} .
    \end{equation*}
    Under the assumptions of \Cref{thm:pr_average}, the PR average with \stepsize exponent $a = 2/3$ achieves the error
    \begin{equation*}
        \lVert \bar{x}_n - x^* \rVert 
        \leq 
    \frac{1}{\sqrt{n}} \left( 
    \mathbb{E}\lVert \bar{\Sigma}^{1/2}Z \rVert
    + \sqrt{2 \lVert \bar{\Sigma} \rVert \log \frac{2}{\confint}}
    \right)
    + 
    \frac{\bar{C}_{p}}{n^{2/3}} \left(\frac{2}{\confint} \right)^{1/p}
        .
    \end{equation*}
\end{corollary}
\begin{proof}
    See Section \ref{app:high_probability}.
\end{proof}
The above bounds improve over the Markov inequalities \eqref{eq:markov_last} and \eqref{eq:markov_pr} in several aspects. 
The dependencies on the $p$-moments $\lVert \Sigma^{1/2} Z \rVert_{L_p}$ and $\lVert \bar{\Sigma}^{1/2}Z \rVert_{L_p}$ are replaced by the asymptotic means $\mathbb{E}\lVert \Sigma_a^{1/2}Z \rVert$ and $\mathbb{E}\lVert \bar{\Sigma}^{1/2}Z\rVert$, which are strictly smaller by Jensen's inequality. 
Unlike the $\confint^{-1/p}$ factors that arise in the dominant terms, the estimates in Corollary \ref{cor:high_probability} are multiplied by transient terms, and the confidence interval estimates are asymptotically tight, meaning the Gaussian tails \eqref{eq:gaussian_tail_estimates} are recovered as $n \to \infty$. 


Provided that \Crefrange{ass:step_size}{ass:CLT} hold for sufficiently large values of $p \geq 1$, the bound in \Cref{cor:high_probability} can be optimized by tuning $p$ as a function of $\confint$. 
In \Cref{sec:lsa}, we establish the validity of these conditions for linear stochastic approximation and demonstrate that this optimization yields sharp tail guarantees. 









\section{Applications}\label{sec:applications}

In this section, we apply our main results to two key problems in stochastic approximation to demonstrate how non-asymptotic Wasserstein bounds provide a finer characterization of the error distribution.
In \Cref{sec:lsa}, we investigate linear stochastic approximation (LSA). 
The error distribution in LSA exhibits a notable dichotomy: while classical theorems establish asymptotic normality, recent finite-sample analyses have shown that the error tails are necessarily heavy tailed \citep{chen_contractive_concentration}. 
Our framework provides a quantitative bridge between these two regimes, showing that our high-probability bounds become asymptotically sharper than what is guaranteed by the finite-sample heavy-tailed analysis.


In \Cref{sec:mcmc}, we consider stochastic gradient descent (SGD) with Markovian data and nonlinear objectives. 
Our results provide a quantitative version of the general asymptotic CLT in \citep{borkar25} by establishing an explicit rate of convergence. 
Furthermore, our analysis extends the non-asymptotic CLT for the Polyak-Ruppert average established in \citep{srikant2024CLT} from the linear to the nonlinear setting.



For brevity, we focus our exposition primarily on the last-iterate bounds, noting that analogous results hold for the Polyak-Ruppert averaged iterates via \Cref{thm:pr_average}.





\subsection{Linear Stochastic Approximation: Weibull and Gaussian Tails}\label{sec:lsa}
We first apply our framework to linear stochastic approximation (LSA) with multiplicative i.i.d. noise, a setting that encompasses algorithms such as online least squares and temporal difference (TD) learning. 
A notable feature of LSA is that the error distribution can exhibit heavy tails even when the underlying noise processes are uniformly bounded. 
As shown in \citep{chen_contractive_concentration}, the finite-sample distribution of the iterates necessarily follows a Weibull-type tail, which is heavier than the asymptotic Gaussian limit's. 


Our theory provides a granular, \nonasymptotic analysis that quantifies the transition from the finite-time heavy tail behavior to the asymptotic Gaussian limit.
We demonstrate that the error can be understood as having two components: a transient, heavy-tailed term reflecting the finite-sample risk identified in \citep{chen_contractive_concentration}, and an asymptotically dominant sub-Gaussian term consistent with the CLT. The key result of this section is a quantification of the rate at which this heavy-tailed influence vanishes, providing a precise account of how the error distribution approaches its Gaussian limit.
From a practical standpoint, this refined understanding is useful for constructing confidence intervals and setting stopping criteria. 
Our result moves beyond asymptotic normality approximations, offering a principled way to quantify the transient, heavy-tailed risk inherent in the early phases of the algorithm and to determine the iteration count after which Gaussian-based approximations become reliable.









Our analysis leverages the tail bounds from \citep{chen_contractive_concentration} to verify \Cref{ass:yn_moment}. 
Consider linear stochastic approximation with multiplicative i.i.d. noise
\begin{align*}
    x_{k+1} &= x_k - \gamma_k (A(\xi_k) x_k - b(\xi_k)) , 
\end{align*}
where $\{\xi_k\}$ is an i.i.d. sequence. 
The following is assumed. 
\begin{lsa}\label{ass:lsa}
    Let $\{\xi_k\}$ be an i.i.d. sequence with $\lVert A(\xi_1) \rVert_{L_\infty} < \infty$ and $\lVert b(\xi_1) \rVert_{L_\infty} < \infty$. 
    \begin{assumptionsub}
        \item The matrix $- \mathbb{E}_\pi A(\xi)$ is Hurwitz stable. 
        \item The \stepsize is chosen as $\gamma_n =\gamma_1/n$ for all $n$ and $\gamma_1$ such that \Cref{ass:general_sa} is satisfied. 
    \end{assumptionsub}    
\end{lsa}
Clearly, \Cref{ass:lsa} implies all of \Crefrange{ass:step_size}{ass:last} except for \Cref{ass:yn_moment}. 
\citet[Theorem 2.1 (a)]{chen_contractive_concentration} establish a high probability bound, where for some constants $\alpha > 1$, $D_1 > 0$, and sequences $\{D_{2,n}\}, \{D_{3,n}\}$, there exists some $\gamma_1 > 0$ satisfying \Cref{ass:general_sa} such that
\begin{equation}\label{eq:cmz_high_probability}
    \lVert x_{n} - x^* \rVert
    \leq 
    \gamma_n^{1/2} \sqrt{D_{1}} \lVert x_1 - x^* \rVert \left[\log \frac{\alpha}{\confint} + \max\left\{ D_{2, n}, D_{3, n} \right\} \right]^{\alpha/2}
\end{equation}
with probability $\geq 1 - \confint$. 


The high probability bound in \eqref{eq:cmz_high_probability} is obtained by exploiting a recurrence of the exponential moments, rather than a direct bound on $\lVert y_n \rVert_{L_p}$. 
A standard argument can be used to recover the corresponding moment bound and verify \Cref{ass:yn_moment}. 
The following requires the use of the Gamma function, and we use $\gfunc$ to denote the Gamma function to avoid confusion with the asymptotic covariance $\Gamma$.
\begin{proposition}\label{prop:weibull_moment}
    Suppose a deterministic initialization $x_1$. 
    Under \Cref{ass:lsa} and for sufficiently small $\gamma_1 > 0$, \Cref{ass:yn_moment} is satisfied with $\lVert y_n \rVert_{L_p} \leq \beta_p^y \lVert \Sigma_a^{1/2}Z \rVert_{L_p}$, where 
    \begin{equation}
        \beta_p^y = D_1 \alpha \exp \left(\frac{\max\left\{D_{2, n}, D_{3, n} \right\}}{p}  \right) \frac{\lVert x_1 - x^* \rVert}{\lVert \Sigma_a^{1/2}Z \rVert_{L_p}}  
         \gfunc \left(\frac{\alpha p}{2} + 1\right)^{1/p}
        , 
    \end{equation}
\end{proposition}
\begin{proof}
    The high probability bound \eqref{eq:cmz_high_probability} corresponds to the Weibull tail 
    \begin{align*}
    \mathbb{P}\left( \lVert x_{n} - x^* \rVert > t \right) 
    \leq 
    \alpha \exp \left(\max\{D_{2, n}, D_{3,n}\} \right)
    \cdot 
    \exp \left( 
    - \left(\frac{t}{\sqrt{D_{1}}} \right)^{2/\alpha}
    \right) . 
\end{align*}
The corresponding moments are obtained as 
\begin{align*}
    \mathbb{E} \lVert x_{n} - x^* \rVert^p 
    &\leq 
    \int_0^\infty p t^{p-1} \mathbb{P}(\lVert x_{n} - x^* \rVert > t) dt 
    \\ &\leq 
    \left(\alpha e^{D_{2, n} \vee D_{3, n}} \right) D_{1}^{p/2} \gfunc \left(\frac{\alpha p}{2} +1\right) .         
\end{align*}
Substituting $\lVert y_n \rVert_{L_p} = \sqrt{\gamma_n} \mathbb{E}^{1/p} \lVert x_n - x^* \rVert^p$, we obtain the result. 
\end{proof}


We establish a refined sample complexity that illustrates the interplay between Gaussian tails, which arises due to the central limit theorem, and the Weibull tail.
We first refine the general bound in \Cref{thm:general} to exploit the structure of LSA. 
In the general \nonlinear setting, the bound depends on the $2p$-th moment $(\beta_{2p}^y)$ of the iterates due to the second-order Taylor expansion of the operator $f$, and the Cauchy-Schwartz inequality \eqref{eq:holder}. 
However, for LSA, the operator $f(x, \xi) = A(\xi) x - b(\xi)$ is linear. 
This allows us to avoid the Taylor expansion and replace the Cauchy-Schwartz inequality with H\"{o}lder's inequality, as explained in the discussion following \eqref{eq:holder}, which relaxes the requirement to the $p$-th moment $\beta_p^y$. 
Consequently, for LSA, \Cref{thm:general} holds with all occurrences of $\beta_{2p}^y$ replaced by $\beta_p^y$. 
Furthermore, since $a = 1$, the $E_n^a/\gamma_n$ term decays as $\mathcal{O}(1/n)$, which is negligible compared to the transient $\mathcal{O}(\gamma_n^{1/3})$ rate. 





We now turn to the high-probability guarantee. 
Under \Cref{ass:lsa}, the noise moments $\{\beta_p^M\}, \{\beta_p^\Phi\}, \{\beta_p^{\Phi^A}\}, \{\beta_p^f\}$ are uniformly bounded in $p$. 
We observe that there exists a threshold $n_{d, p}$ such that for all sample sizes $n \geq n_{d, p}$, the bound in \Cref{thm:general} is given by the form
\begin{equation}\label{eq:thm1_application_linear}
    \mathcal{W}_p (y_n, U_\infty) 
    \leq
    C_1 p \gamma_n^{1/6} + C_2 p \beta_p^y \gamma_n^{1/3}
\end{equation}
for some constants $C_1, C_2 > 0$ independent of $p$ and $n$. 
By Stirling's approximation, we obtain $\beta_p^y \sim p^{\alpha/2}$. 
For any $\confint \in (0, 1)$ such that $p^*(\alpha, \delta) = \log (1/\delta)/(1+\alpha/2) \geq 2$, the bound in \Cref{cor:high_probability} is optimized at $p = p^*(\alpha, \delta)$, and the following holds with probability $\geq 1 - \delta$:
\begin{align*}
    \lVert x_n - x^* \rVert\leq 
    \gamma_n^{1/2}\left( \mathbb{E}\lVert \Sigma_a^{1/2} Z \rVert 
    + \sqrt{2  \lVert \Sigma_a \rVert \log \frac{2}{\confint}} \right) + 
    \mathcal{O}\left( \left[1+ \gamma_n^{1/6}\log^{\alpha/2} \frac{1}{\confint}  \right]\gamma_n^{2/3}\log \frac{1}{\confint}
    \right)
    .
\end{align*}
This result provides a granular decomposition of the error bound in \eqref{eq:cmz_high_probability}. 
The bound in \eqref{eq:cmz_high_probability} is globally dominated by the heavy-tailed factor $\log^{\alpha/2} (1/\delta)$, reflecting the inherent sample complexity arising from multplicative noise identified by \citet{chen_contractive_concentration}. 
In contrast, our bound isolates this heavy-tailed behavior within a transient term scaling as $\tilde{\mathcal{O}}(\gamma_n^{2/3} \log^{\alpha/2} (1/\delta))$, while the leading order term scales with the Gaussian tail $\sqrt{\log (2/\delta)}$. 


This structural distinction quantifies the phase transition between the two regimes. 
For small sample sizes, the transient term dominates, recovering the heavy-tailed behavior predicted by \Cref{prop:weibull_moment}. 
However, as $n$ increases, the heavy-tailed component decays faster than the leading Gaussian term. 
This confirms that the heavy-tailed dynamics are strictly transient and are eventually subsumed by the central limit theorem, validating the use of Gaussian approximations once the sample size is sufficiently large. 



\subsection{Stochastic Gradient Descent with Markov Data}\label{sec:mcmc}
We now analyze a class of stochastic optimization problems where gradients are estimated using Markovian data. 
This setting is ubiquitous in large-scale machine learning and statistical physics, where the objective $\mathbb{E}_{\pi}[\ell(x, \xi)]$ is minimized, but the stationary distribution is intractable or accessible only via an ergodic Markov chain $\{\xi_k \}$ with transition kernel $P$.



To establish rigorous finite-sample guarantees in this setting, we focus on irreducible and aperiodic Markov chains that satisfy the DV3 variant of the Donsker-Varadhan condition following \citep{borkar25}. 
This condition is satisfied by many Markov chains, and is strong enough to imply several conditions assumed in our main results. 
Recall that a non-negative function $s: \mathcal{S} \to [0, 1]$ is said to be small if there exists a positive integer $m$, and a non-zero measure $\nu$ on $\mathcal{S}$ such that 
\begin{align*}
    P^m (x, A) \geq s(x) \nu(A) 
\end{align*}
holds for all $x \in \mathcal{S}$ and measurable sets $\mathcal{A} \subset \mathcal{S}$. 
\begin{dv3}\label{ass:DV3}
    There exist functions $\tilde{V}: \mathcal{S} \to \mathbb{R}_+$, $\tilde{V}': \mathcal{S} \to [1, \infty)$, a small function $s: \mathcal{S} \to [0, 1]$, and a constant $b > 0$ such that
    \begin{equation}
        \mathbb{E}\left[\exp(\tilde{V}(\xi_{k+1})) \mid \xi_k = \xi\right] \leq \exp(\tilde{V}(\xi) - \tilde{V}'(\xi) + b s(\xi)).
    \end{equation}
    For every $\tilde{v}_0 > 0$, the sublevel set $\{\xi: \tilde{V}'(\xi) \leq \tilde{v}_0\}$ is small or empty, and $\tilde{V}$ is finite on this sublevel set. 
\end{dv3}
The DV3 condition is satisfied by several important classes of Markov chains.
An irreducible and aperiodic finite Markov chain trivially satisfies \Cref{ass:DV3} with constants $\tilde{V}, \tilde{V}', b, s$. 
On general state spaces, \Cref{ass:DV3} is satisfied by a large class of Markov chains including those obtained from Metropolis-Hastings algorithms \citep{roberts1996geometric,mengersen96,JARNER2000341} and discretizations of diffusion processes with appropriate regularity conditions \citep{mattingly2002ergodicity,cattiaux2010note}. 




The DV3 condition serves as a primitive condition under which our results hold. 
When a Markov chain satisfies \Cref{ass:DV3}, then \Cref{ass:noise_markov:drift} holds with $V = e^{\tilde{V}}$. 
The DV3 condition implies that the chain mixes exponentially fast and has exponential moments. 
Invoking \citep[Theorem 3]{srikant2024CLT}, we obtain that \Cref{ass:CLT} is satisfied. 


The Hurwitz condition in \eqref{eq:hurwtiz} is not sufficient to imply \Cref{ass:yn_moment}.  
To this end, we follow \citep{borkar25} and assume the ODE@$\infty$ condition. 
Specifically, the condition states that the limit $\bar{f}_\infty (x) = \lim_{r \to \infty} r^{-1} \bar{f}(r x)$ exists and the origin is a globally asymptotic stable equilibrium for the ODE $\dot{x} = \bar{f}_\infty (x)$. 
\begin{sgd}\label{ass:sgd}
    The Markov chain $\{\xi_k\}$ satisfies \Cref{ass:DV3}.  
    The function $f$ satisfies \Cref{ass:jacobian,ass:general_sa}, the ODE@$\infty$ condition, and \Cref{ass:noise_operator_compatibility} with $V = e^{\tilde{V}}$. 
    Further, $\bar{f}$ is continuously differentiable and its Jacobian $\nabla \bar{f}$ is bounded. 
\end{sgd}

\Cref{ass:yn_moment} follows from the results of \citep{borkar25}, and we obtain the rate at which $y_n$ converges to its Gaussian limit $\mathcal{N}(0, \Sigma_a)$. 
\begin{theorem}\label{thm:mcmc}
    Consider minimizing a $\mu$-strongly convex function with access to stochastic gradients of the form $f(x, \xi_k)$ subject to \Cref{ass:sgd}. 
    Let $\{\gamma_k\}$ be chosen such that \Cref{ass:step_size} are satisfied with $a > 1/2$ and $\gamma_1 \mu > 1/2$. 

    Then, the CLT holds for the sequences $\{y_k\}$ and $\{\bar{y}_n\}$ with $\mathcal{W}_1(y_n, \Sigma_a^{1/2}Z) = \mathcal{O}(\gamma_n^{1/6} \log \gamma_n^{-1})$ and $\mathcal{W}_1 (\bar{y}_n, \bar{\Sigma}^{1/2} Z) = \mathcal{O}( (\gamma_n n)^{-1/2} + n^{1/2} \gamma_n)$. 
\end{theorem}
\Cref{thm:mcmc} applies to a broad class of stochastic optimization problems where the data stream is generated by Markov Chain Monte Carlo (MCMC) sampling. 
A canonical example is maximum likelihood estimation for latent variable models, where the gradient of the log-likelihood involves an expectation over the posterior of the latent variables \citep{geyer1994convergence}. 
When this posterior is intractable, it is approximated by running a Metropolis-Hastings sampler or unadjusted Langevin algorithm (ULA) targeting the posterior distribution. 
Provided the target density satisfies standard curvature conditions (e.g., log-concavity or dissipative drift), these samplers satisfy the DV3 condition, thereby enabling the rigorous application of our convergence rates to these intractable inference problems.








\section{Extensions and Open Problems}
In this work, we introduced a framework for deriving $\mathcal{W}_p$ error bounds for stochastic approximation. 
This section discusses some potential improvements and limitations of the probabilistic framework considered in this paper. 




\subsection{Wasserstein-$p$ Convergence in CLT for Martingale Differences}
Our analysis crucially relies on \nonasymptotic central limit theorems to bound the distance between the stochastic approximation iterates and their Gaussian limits. 
While sharp rates in the Wasserstein-$p$ distance for all $p \geq 1$ are well-established for sums of independent random variables \citep{bonis2020steinsmethodnormalapproximation,bonis2024improved}, the literature on more general martingale difference sequences is currently limited to $\mathcal{W}_1$.



A feature of our proof technique is its modularity with respect to the underlying CLT. 
We emphasize that the restriction to $\mathcal{W}_1$ is a consequence of the current state of martingale limit theory, not an intrinsic limitation of our analysis. 
Should a quantitative martingale CLT for $\mathcal{W}_p$ with $p \geq 2$ be established, our framework would seamlessly extend these convergence rates to stronger metrics without modification. 
This would provide tighter control over the moments of the error distribution and a significant strengthening for applications requiring high-confidence bounds. 










\subsection{Delayed Gratification: Transient Phase of Polyak-Ruppert Averaging}\label{disc:fast_rates}
A cornerstone result in the theory of stochastic approximation is that Polyak-Ruppert (PR) averaging outperforms standard SA with an asymptotically optimal convergence rate. 
When the last iterate is generated with \stepsize $\gamma_n \propto 1/n$ and the PR average is generated with \stepsize $\gamma_n \propto n^{-2/3}$, \Cref{thm:general,thm:pr_average} imply the bounds
\begin{equation}\label{eq:moment_bounds_discussion}
\begin{split}
    \lVert x_n - x^* \rVert_{L_p} 
    &\leq \sqrt{\gamma_1} \frac{\lVert \Sigma_1^{1/2} Z \rVert_{L_p}}{\sqrt{n}} + \mathcal{O}\left( \frac{1}{n^{2/3}} \right), 
    \\ 
    \lVert \bar{x}_n - x^* \rVert_{L_p} &\leq 
    \frac{\lVert \bar{\Sigma}^{1/2}Z \rVert_{L_p}}{\sqrt{n}} + \mathcal{O}\left(\frac{1}{n^{2/3}} \right) .     
\end{split}    
\end{equation}
The relation $\bar{\Sigma} \preceq \gamma_1 \Sigma_1$ implies asymptotic efficiency of the PR average. 



Asymptotic theory and existing finite-time bounds describe the superiority of Polyak-Ruppert averaging, which is consistent with the bounds in \eqref{eq:moment_bounds_discussion}. 
In practice, however, the last iterate is sometimes observed to yield lower error for a large number of iterations \citep{moulines_bach}. 
Unlike prior works on finite-time bounds such as \citep{haque2023tightfinitetimebounds}, the leading terms in \eqref{eq:moment_bounds_discussion} are exact. 
This suggests that the observation in \citep{moulines_bach} may be attributable to the higher order terms in the bounds in \eqref{eq:moment_bounds_discussion}.
A refined analysis characterizing the critical sample size $n$ required for the averaged estimator to surpass the last iterate remains a subject of theoretical interest.
















\section{Conclusion}
In this paper, we developed a \nonasymptotic framework for analyzing the distributional convergence of \nonlinear stochastic approximation algorithms. 
Our main results provide explicit, finite-sample bounds on the Wasserstein-$p$ distance between the law of the normalized iterates and their limiting Gaussian distribution. 
These bounds serve as a quantitative counterpart to classical asymptotic CLTs and provide a richer characterization of the algorithm's behavior than what is available from moment bounds alone.


A key feature of our framework is its modularity, which allows us to translate \nonasymptotic CLTs for the underlying noise process into guarantees for both last-iterate and Polyak-Ruppert averaged estimators. 
As a direct consequence, our Wasserstein bounds yield high-probability concentration inequalities that are asymptotically sharp and capture the transition from finite-sample, potentially heavy-tailed dynamics to the Gaussian behavior predicted by the central limit theorem. 
We demonstrated the utility of this approach by providing a refined analysis of the error tails in linear stochastic approximation and establishing distributional convergence rates for SGD with Markovian data.








\section*{Acknowledgement}
We thank Professor Siva Theja Maguluri from Georgia Institute of Technology for drawing our attention to the reference \citep{liu2023wassersteinpboundscentrallimit}.


\appendix
\section{Auxiliary Results}
\subsection{Convergence of a Recursive Sequence}\label{app:recursion_convergence}
The following recursion arises frequently, hence we dedicate a result that will be used repeatedly. 
\begin{lemma}\label{lem:recursion}
    Consider a non-negative scalar sequence $\{x_k\}$ that satisfies the recursion
    \begin{equation}\label{eq:recursion_lemma}
        x_{k+1} \leq (1 - \lambda \gamma_k) x_k + \alpha_k , 
    \end{equation}
    where $\lambda \gamma_1 \in (0, 1)$, $\gamma_k \propto k^{-a}$, and $\alpha_k \propto k^{-b}$ for some $0 < a  < 1$ and $b > a$. 
    The solution to \eqref{eq:recursion_lemma} satisfies
    \begin{equation}
        x_{n+1} \leq \frac{1}{\lambda} \frac{\alpha_n}{\gamma_n} 
        + \mathcal{O}\left(\frac{\alpha_n}{\gamma_n^2 n} \right)
        .
    \end{equation} 
\end{lemma}
\begin{proof}
    The result is proved using induction, where we first prove the induction step and prove the base case afterwards. 
    Suppose for some $n \geq 1$ and any constant $C > 1/\lambda$,
    \begin{align*}
        x_{n} \leq C \cdot \frac{\alpha_n}{\gamma_n} . 
    \end{align*}
    From the recursion,
    \begin{align*}
        x_{n+1} \leq (1 - \lambda \gamma_n) x_n + \alpha_n 
        \leq \frac{\alpha_n}{\gamma_n} \left[(1 - \lambda \gamma_n) x_n \cdot \frac{\gamma_n}{\alpha_n} + \gamma_n  \right] .
    \end{align*}
    From the premise $x_n \leq C \cdot  \alpha_n/\gamma_n$ we have
    \begin{align*}
        x_{n+1} \leq  \frac{\alpha_n}{\gamma_n} \left[C(1 - \lambda \gamma_n)  + \gamma_n  \right] 
        = C \frac{\alpha_n}{\gamma_n} + \left(1 - C \lambda \right)\alpha_n 
        .
    \end{align*}
    The statement $x_{n+1} \leq C \alpha_{n+1}/\gamma_{n+1}$ holds if
    \begin{align*}
        C \frac{\alpha_n}{\gamma_n} + \left(1 - C \lambda \right)\alpha_n 
        \leq 
        C\frac{\alpha_{n+1}}{\gamma_{n+1}}
        \Leftrightarrow 
        \frac{1-C\lambda}{C} 
        &\leq 
        \alpha_n^{-1} \left(\frac{\alpha_{n+1}}{\gamma_{n+1}} - \frac{\alpha_n}{\gamma_n} \right)
        \\ &
        =  \frac{1}{\gamma_n} \left(
        \frac{\alpha_{n+1}/\alpha_n}{\gamma_{n+1}/\gamma_n} - 1\right)
    \end{align*}
    Using that
    \begin{align*}
        \frac{\alpha_{n+1}/\alpha_n}{\gamma_{n+1}/\gamma_n} 
        = 1 + \frac{a-b}{n} + \Theta\left(n^{-2} \right), 
    \end{align*}
    we obtain 
    \begin{equation}
        \frac{1 - C \lambda}{C} 
        =
        \mathcal{O}\left(\frac{1}{\gamma_n n} \right)
    \end{equation}
    The upper bound decays to zero, and the induction statement holds for any constant $C > 1/\lambda$ that satisfies the above. 
    By rearranging the inequality, we obtain that the induction statement holds with a more precise constant $C = \lambda^{-1} + \mathcal{O}(1/\gamma_n n)$. 
    

    We now prove the entry point, that there exists some $n \geq 1$ such that $x_n \leq C \alpha_n/\gamma_n$. 
    In proving the above statement, we showed that for any $k \geq 1$ and a constant $C > 1/\lambda$,
    \begin{align*}
        (1 - \lambda \gamma_k) C \frac{\alpha_k}{\gamma_k} + \alpha_k \leq C \frac{\alpha_{k+1}}{\gamma_{k+1}} ,
    \end{align*}
    which implies
    \begin{align*}
        \alpha_k \leq C \frac{\alpha_{k+1}}{\gamma_{k+1}} - (1 - \lambda \gamma_k) C \frac{\alpha_k}{\gamma_k} . 
    \end{align*}
    From the recurrence relation $x_{k+1} \leq (1 - \lambda \gamma_k) x_k + \alpha_k$, we have from the preceding equation:
    \begin{align*}
        x_{k+1} \leq (1 - \lambda \gamma_k) x_k + \left[ C \frac{\alpha_{k+1}}{\gamma_{k+1}} - (1 - \lambda \gamma_k) C \frac{\alpha_k}{\gamma_k}\right] . 
    \end{align*}
    Rearranging, we obtain
    \begin{align*}
        x_{k+1} - C \frac{\alpha_{k+1}}{\gamma_{k+1}}
        \leq (1 - \lambda \gamma_k) \left( x_k - C \frac{\alpha_k}{\gamma_k}\right) . 
    \end{align*}
    This implies that
    \begin{equation}
        x_{n+1} - C \frac{\alpha_{n+1}}{\gamma_{n+1}} \leq \prod_{k=1}^n (1 - \lambda \gamma_k) \left(x_1 - C \frac{\alpha_1}{\gamma_1}\right) 
    \end{equation}
    The statement follows by contradiction: the sequence cannot satisfy $x_{n} > C \alpha_n/\gamma_n$ for all $n \geq 1$. 
    
    
\end{proof}



\subsection{Rosenthal's Inequalities}\label{sec:BDG}
We use the Rosenthal inequality for martingale difference sequences from \citep{pinelis1994optimum}, which stems from the ideas of \citep{rosenthal1970subspaces,burkholder1973distribution}, and refer to the inequality as Burkholder-Rosenthal (BR) inequality. 


We first consider $p \geq 2$. 
For a MDS $\{M_k\}_{k=1}^n$, the BR inequality \citep[Theorem 4.1]{pinelis1994optimum} states that for some constant $C_p^{BR} > 0$,
\begin{align*}
    \mathbb{E}^{1/p}  \left(\max_{k \leq n} \left\lVert \sum_{j=1}^k M_j \right\rVert \right)^p
    \leq C_p^{BR} \left( 
    \left\lvert 
    \left( \sum_{k=1}^n \mathbb{E}[\lVert M_k \rVert^2 | \mathcal{H}_{k-1}]\right)^{1/2}
    \right\rvert_{L_p}
    +
    \left\lvert 
    \max_{k \leq n} \lVert M_k \rVert
    \right\rvert_{L_p}
    \right) . 
\end{align*}
From the inequality $\max_{k \leq n} a_k \leq (\sum_k a_k^p)^{1/p}$ for any $p \geq 1$ and non-negative scalar sequence $\{a_k\}$, we obtain the inequalities
\begin{align*}
    \mathbb{E}^{1/p} \left(\max_{k \leq n} \lVert M_k \rVert \right)^p
    \leq 
    \mathbb{E}^{1/p} \left(\sum_{k=1}^n \lVert M_k \rVert^p\right)
    &\leq 
    \mathbb{E}^{1/p}\left[ \left(
    \max_{k \leq n} \left\lVert \sum_{j=1}^k M_j \right\rVert \right)^p
    \right]
    \\ &
    \leq \left(\sum_{k=1}^n \lVert M_k \rVert^p_{L_p} \right)^{1/p}
    . 
\end{align*}
Substituting for the second term in the upper bound of the BR inequality, we obtain
\begin{equation}\label{eq:BR}
    \left\lVert \sum_{k=1}^n M_k \right\rVert_{L_p}
    \leq C_p^{BR} \left(
        \sqrt{\left\lvert \sum_{k=1}^n \mathbb{E}\left[ \lVert M_k \rVert^2 | \mathcal{H}_{k-1} \right] \right\rvert_{L_{p/2}}}
        + \left(\sum_{k=1}^n \lVert M_k \rVert_{L_p}^p
    \right)^{1/p} \right) .
\end{equation}


We apply the BR inequality for a weighted sum of a MDS. 
For any deterministic sequence $\{\alpha_k\}$, the sequence $\{\alpha_k M_k\}$ is another MDS.
From \eqref{eq:BR} and \Cref{ass:noise_martingale}, we obtain for any $p \geq 2$,
\begin{align*}
    \left\lVert  
    \sum_{k=1}^n \alpha_k M_k
    \right\rVert_{L_p}
    &\leq 
    C_p^{BR} \left(
    \sqrt{\left\lvert \sum_{k=1}^n \alpha_k^2 \mathbb{E}\left[\lVert M_k \rVert^2 | \mathcal{H}_{k-1} \right] \right\rvert_{L_{p/2}}}
    + 
    \left(\sum_{k=1}^n \alpha_k^p \lVert M_k \rVert_{L_p}^p
    \right)^{1/p}\right)
    \\ &
    \leq C_p^{BR}  \sup_{k \geq 1} \left\{ \sqrt{\left\lvert 
        \mathbb{E}\left[  \left\lVert M_k \right\rVert^2 | \mathcal{H}_{k-1} \right]
        \right\rvert_{L_{p/2 }}} + \lVert M_k \rVert_{L_p} \right\} \sqrt{\sum_{k=1}^n \alpha_k^2}
    \\ &
    = C_p^{BR} \rosenthal{M} \sqrt{\sum_{k=1}^n \alpha_k^2} . 
    \numberthis \label{eq:br_inequality}
\end{align*}
By Jensen's inequality, it holds that for any $p \in [1, 2]$,
\begin{equation}
    \left\lVert \sum_{k=1}^n \alpha_k M_k \right\rVert_{L_p} \leq \left\lVert \sum_{k=1}^n \alpha_k M_k \right\rVert_{L_2}
    \leq C_2^{BR} \rosn{2}{M} \sqrt{\sum_{k=1}^n \alpha_k^2} .
\end{equation}
When $p < 2$, we therefore denote by $C_p^{BR}$ the universal constant $C_2^{BR}$. 
We will often use the BR inequality for a partial sum, where for any $k_1 < k_2$, 
\begin{align*}
    \Lp{\sum_{k=k_1}^{k_2 - 1} \alpha_k M_k}
    \leq C_p^{BR} \rosn{p}{M} \sqrt{\sum_{k=k_1}^{k_2 - 1} \alpha_k^2} . 
\end{align*}














\subsection{Nonasymptotic Central Limit Theorems for Weighted Sums}\label{proof:martingale_clt_weighted}
Here we prove \Cref{lem:martingale_clt_weighted}.


Let $\hat{Z}_m = (\cumstep_m \Gamma)^{-1/2}\sum_{k=k_m}^{k_{m+1} - 1} \sqrt{\gamma_k} M_k$ be the weighted average of interest, and $S_m = (I_m\Gamma)^{-1/2} \sum_{k=k_m}^{k_{m+1} - 1} M_k$ the unweighted counterpart.
By the triangle inequality, 
\begin{align*}
    \mathcal{W}_{p}(\hat{Z}_m, Z) 
    &\leq 
    \lVert \hat{Z}_m -  S_m\rVert_{L_{p}} +
    \mathcal{W}_{p} (S_m, Z) 
    .
\end{align*}




We have by the BR inequality that for $\alpha_k = \cumstep_m^{-1/2} \sqrt{\gamma_k} - I_m^{-1/2}$,
\begin{align*}
    \lVert \hat{Z}_m -  S_m\rVert_{L_p} 
    = 
    \left\lVert \sum_{k=k_m}^{k_{m+1} - 1} \alpha_k \Gamma^{-1/2} M_k  \right\rVert_{L_p}
    &\leq C^{BR}_{p} \rosenthal{\Gamma^{-1/2}M} \sqrt{\sum_{k=k_m}^{k_{m+1} - 1} \alpha_k^2} . 
\end{align*}
The sum is evaluated next. 
We have for $\bar{\gamma}_m = \cumstep_m/I_m$ the identities
\begin{equation}\label{eq:sqrt_gamma_difference}
    \sqrt{\sum_{k=k_m}^{k_{m+1} - 1} \alpha_k^2}
    = \sqrt{
    \frac{1}{I_m} \sum_{k=k_m}^{k_{m+1} - 1}\left( \sqrt{\frac{I_m}{\cumstep_m} \gamma_k} - 1 \right)^2}
  = \sqrt{\frac{1}{I_m \bar{\gamma}_m} \sum_{k=k_m}^{k_{m+1}-1} \left(\sqrt{\gamma_k} - \sqrt{\bar{\gamma}_m} \right)^2}
    .    
\end{equation}
Using Taylor's remainder theorem for the summand, we obtain
\begin{align*}
    \sum_{k=k_m}^{k_{m+1} - 1} (\sqrt{\gamma_k} - \sqrt{\gamma_{k_m}})^2
    \leq \frac{a^2}{4}\sum_{k=k_m}^{k_{m+1} - 1} \gamma_{k_m}\frac{(k - k_m)^2}{k_m^2}
    \leq \frac{a^2 I_m^3}{4 k_m^2} \gamma_{k_m} . 
\end{align*}
Substituting for \eqref{eq:sqrt_gamma_difference}, we obtain
\begin{align*}
    \sqrt{\sum_{k=k_m}^{k_{m+1} - 1} \alpha_k^2}
    \leq \frac{a I_m}{2 k_m} \sqrt{\frac{\gamma_{k_m}}{\bar{\gamma}_m}}
    \leq \frac{aI_m}{2 k_m} \left(1 + \frac{I_m}{k_m} \right)^{a/2} 
    \leq \frac{a I_m}{k_m} 
    ,
\end{align*}
where the last inequality uses that $I_m \leq k_m$. 
Summarizing, we have shown that
\begin{equation}\label{eq:weighted_clt}
    \mathcal{W}_p (\hat{Z}_m, Z) 
    \leq 
    C_{p}^{BR} \rosenthal{\Gamma^{-1/2} M} \frac{aI_m}{k_m} 
    +
    \mathcal{W}_{p} (S_m, Z) . 
\end{equation}


\subsection{Proof of \Crefrange{lem:martingale_clt}{lem:bonis_weighted}}\label{app:clt_weighted}
Any weighted combination of a mean-zero iid sequence $\{M_k\}_{k=1}^n$ satisfying $\mathrm{Cov}M_1 = \Gamma$ and $\lVert \sqrt{\gamma_1} M_1 \rVert_{L_p} < \infty$ for $p > 2$ converges asymptotically to a Gaussian distribution as
\begin{equation}
    \lim_{n \to \infty} \frac{1}{\sqrt{\sum_{k=1}^n \gamma_k}} \sum_{k=1}^n \sqrt{\gamma_k} M_k \overset{d.}{=} \mathcal{N}\left(0, \Gamma \right) 
\end{equation}
by the Lyapunov central limit theorem. 
Nonasymptotic CLTs often are established for the unweighted sum $S_m$, and we use \Cref{lem:martingale_clt_weighted} to establish the $\mathcal{W}_p$--CLT for the weighted $\hat{Z}_m$.



For this section, we use $\mathcal{W}_{p, 2}$ to denote the Wasserstein-$p$ distance defined with the Euclidean cost
\begin{align*}
    \mathcal{W}_{p, 2}(X, Y) = \inf \mathbb{E}^{1/p} \lVert X - Y \rVert_{2}^p , 
\end{align*}
where it holds that $\mathcal{W}_p(S_m, Z) \leq \sqrt{\lVert Q \rVert_2} \mathcal{W}_{p, 2} (S_m, Z)$. 







\textbf{Proof of \Cref{lem:martingale_clt}:}
We use \citep[Theorem 1]{srikant2024CLT} to obtain that for some $\tilde{C}_1 > 0$,
\begin{equation}\label{eq:thm1_srikant}
\begin{split}
    \mathcal{W}_{1, 2} \left( 
        S_m, Z
    \right)
    &\leq 
    \tilde{C}_1 \left(\beta_3^{\Gamma^{-1/2} M} \right)^3 \sqrt{\lVert Q \rVert_2 \lVert \Sigma_a \rVert_2} \frac{d \log I_m}{\sqrt{I_m}} \beta_3^{\Gamma^{-1/2} M}     
    \\ & + 
    \frac{\tilde{C}_1}{\sqrt{I_m}} \left\lvert 
    \sum_{k=k_m}^{k_{m+1} - 1} \frac{1}{\sqrt{k_{m+1} - k}}
  \mathrm{Tr} \left( \Gamma^{-1/2}\mathbb{E}\Gamma_k \Gamma^{-1/2} - \identity \right)  \right\rvert
.    
\end{split}
\end{equation}
Define
\begin{align*}
    \Gamma^\xi_k &\coloneqq \mathbb{E}\left[ \left(\Phi(\xi_k) - (P \Phi)(\xi_k) \right)  \left(\Phi(\xi_k) - (P \Phi)(\xi_k) \right)^\dagger \Big| \mathcal{H}_{k-1}  \right],
    \\
    \Gamma^\xi &\coloneqq \mathbb{E}_{\xi_1 \sim \pi} \left[ \left(\Phi(\xi_2) - (P \Phi(\xi_{1}) \right) \left(\Phi(\xi_2) - (P \Phi(\xi_{1}) \right)^\dagger \right]. 
\end{align*}
By \citep[Theorem 3]{srikant2024CLT}, there exists constants $C_\xi > 0$ and $\rho_\xi \in [0, 1)$ such that 
\begin{align*}
    \left\| \mathbb{E}\Gamma_k^\xi - \Gamma^k \right\|_{HS} \leq C_\xi \rho_\xi^k.
\end{align*}
Combining with \Cref{ass:martingale:covariance}, we obtain that
\begin{align*}
    \left\| \mathbb{E}\Gamma_k - \Gamma \right\|_{HS} \leq 
    \left( C_W + C_\xi\right) \left( \rho_W \vee \rho_\xi \right)
    .
\end{align*}
and conclude that for some $\tilde{C}_1 > 0$ and $\rho_M \in (0, 1]$, 
\begin{align*}
    \frac{\tilde{C}_1}{\sqrt{I_m}} \left\lvert 
    \sum_{k=k_m}^{k_{m+1} - 1} \frac{1}{\sqrt{k_{m+1} - k}}
  \mathrm{Tr} \left( \Gamma^{-1/2}\mathbb{E}\Gamma_k \Gamma^{-1/2} - \identity \right)  \right\rvert
    & \leq 
    \frac{\tilde{C}_1'}{\sqrt{I_m}} 
    \sum_{k=k_m}^{k_{m+1} - 1} \frac{\rho_M^k}{\sqrt{k_{m+1} - k}}
    \\ &
    \leq \frac{\tilde{C}_1'}{\sqrt{I_m}} \sum_{j=1}^{I_m} \frac{\rho_M^{k_{m+1} - j}}{\sqrt{j}}
    \\ &
    = \mathcal{O}\left(\rho_M^{k_m}\right) .
\end{align*}
Therefore, when $k_m$ is large, the $\rho_M^{k_m}$ rate can be subsumed by the first term in \eqref{eq:thm1_srikant} to obtain that for some $C_1 > 0$, 
\begin{align*}
    \mathcal{W}_1 (S_m, Z) 
    \leq
    C_1\left(\beta_3^{\Gamma^{-1/2} M} \right)^3 \sqrt{\lVert Q \rVert_2 \lVert \Sigma_a \rVert_2} \frac{d \log I_m}{\sqrt{I_m}}
    ,
\end{align*}
The claimed result is obtained by combining with \Cref{lem:martingale_clt_weighted}. 



\textbf{Proof of \Cref{lem:bonis_weighted}:}
By \citep[Theorem 1]{bonis2020steinsmethodnormalapproximation}, there exists a constant $\tilde{C}_p > 0$ such that
\begin{align*}
    \mathcal{W}_{p, 2} (S_m, Z) 
    &\leq 
    \tilde{C}_p \sqrt{\lVert Q \rVert_2}\frac{d^{1/4} (\beta_4^{\Gamma^{-1/2}M})^2 + (\beta_{p+2}^{\Gamma^{-1/2}M})^{1+2/p} }{\sqrt{I_m}}
\end{align*}
\Cref{lem:bonis_weighted} is obtained by combining with Lemma \ref{lem:martingale_clt_weighted}, using that $\tilde{C}_p$ and $C_p^{BR}$ stem from the same Rosenthal inequality and therefore have the same growth rates with respect to $p$. 








    


\subsection{Error Recursion}\label{app:error_recursion}

We derive the recursion for $y_n = \gamma_n^{-1/2}(x_n - x^*)$ for general stochastic approximation algorithms of the form
\begin{equation}
    x_{k+1} = x_k - \gamma_k (f(x_k, \xi_k) + W_k). 
\end{equation}
Recall $\bar{f}(x) = \mathbb{E}_\pi[f(x, \xi)]$ is the expectation over the stationary distribution of $\{\xi_k\}$. 
From the Poisson equation \eqref{eq:poisson}, the update direction is written as
\begin{align*}
    f(x_k, \xi_k) + W_k 
    &=
    \left[f(x_k, \xi_k) - f(x^*, \xi_k) \right]
    + \left[f(x^*, \xi_k) - \bar{f}(x^*)\right]  + W_k
    \\ &
    =
    \left[f(x_k, \xi_k) - f(x^*, \xi_k) \right]
    + \left[\Phi(\xi_k) - (P\Phi)(\xi_{k-1})\right] 
    \\ &
    +\left[(P\Phi)(\xi_{k-1}) - (P\Phi)(\xi_k) \right]
    + W_k
    \\
    &= \left[f(x_k, \xi_k) - f(x^*, \xi_k) \right]
    + \left[ (P\Phi)(\xi_{k-1}) - (P\Phi)(\xi_{k})\right] - M_k 
    \\ &
    = [f(x_k, \xi_k) - \bar{f}(x_k)]
    - [f(x^*, \xi_k) - \bar{f}(x^*)]
    + \bar{f}(x_k) 
    \\ &
    + \left[ (P\Phi)(\xi_{k-1}) - (P\Phi)(\xi_{k})\right] - M_k 
    .    
\end{align*}
Using the centered function
\begin{equation}\label{eq:gk}
    g(x_k, \xi_k) = f(x_k, \xi_k)  - \bar{f}(x_k)   
\end{equation}
and the Taylor expansion $\bar{f}(x_k) = \nabla \bar{f}(x^*)(x_k - x^*) + (\nabla \bar{f}(\tilde{x}_k) - \nabla \bar{f}(x^*))(x_k - x^*)$ for $\tilde{x}_k = \theta x_k + (1 - \theta) x^*$ with $\theta \in [0, 1]$, we write the update direction as
\begin{align*}
    f(x_k, \xi_k) + W_k 
    &= [g(x_k, \xi_k) - g(x^*, \xi_k)] + \bar{f}(x_k)
    + \left[ (P\Phi)(\xi_{k-1}) - (P\Phi)(\xi_{k})\right] - M_k 
    \\ &
    = \nabla \bar{f}(x^*) (x_k - x^*)
    + (\nabla \bar{f}(\tilde{x}_k) - \nabla \bar{f}(x^*)) (x_k - x^*)
    \\ &
    +[g(x_k, \xi_k) - g(x^*, \xi_k)]
    + \left[ (P\Phi)(\xi_{k-1}) - (P\Phi)(\xi_{k})\right] - M_k 
    .
\end{align*}
The recursion for $\{x_k\}$ is therefore written as
\begin{align*}
    x_{k+1} - x^* &=
    x_k - x^* - \gamma_k \left(\nabla \bar{f}(x^*)(x_k - x^*) - M_k \right)
    \\ &
    -
    \gamma_k (\nabla \bar{f}(\tilde{x}_k) - \nabla \bar{f}(x^*)) (x_k - x^*)
    \\ &
    - \gamma_k \left[ \left(g(x_k, \xi_k) - g(x^*, \xi_k)\right) + \left( (P\Phi)(\xi_{k-1}) - (P\Phi)(\xi_{k})\right) \right]
\end{align*}
Substituting $y_k = \gamma_k^{-1/2}(x_k - x^*)$ and $y_{k+1} = \gamma_{k+1}^{-1/2}(x_{k+1} - x^*)$, we obtain
\begin{equation}\label{eq:yk_draft1}
\begin{split}
    y_{k+1} &= 
    \sqrt{\frac{\gamma_k}{\gamma_{k+1}}} \left(I - \gamma_k \nabla \bar{f}(x^*)\right)y_k + \frac{\gamma_k}{\sqrt{\gamma_{k+1}}} M_k 
    \\ &
    - \sqrt{\frac{\gamma_k}{\gamma_{k+1}}} \cdot \gamma_k (\nabla \bar{f}(\tilde{x}_k) - \nabla \bar{f}(x^*)) y_k 
    \\ &
    - \frac{\gamma_k}{\sqrt{\gamma_{k+1}}} \left[ \left(g(x_k, \xi_k) - g(x^*, \xi_k)\right) + \left[ (P\Phi)(\xi_{k-1}) - (P\Phi)(\xi_{k})\right] \right] .    
\end{split}    
\end{equation}
The operator in the first term is written as
\begin{align*}
    \sqrt{\frac{\gamma_k}{\gamma_{k+1}}} (\mathbb{I} - \gamma_k \nabla \bar{f}(x^*))
    &= \left[\mathbb{I} - \gamma_k \nabla \bar{f}(x^*) \right]
    + \left(\frac{\sqrt{\gamma_k} - \sqrt{\gamma_{k+1}}}{\sqrt{\gamma_{k+1}}} \right) \left[\mathbb{I}
    -  \gamma_k \nabla \bar{f}(x^*) \right]
    . 
\end{align*}
By Taylor's remainder theorem
\begin{align*}
    \frac{\sqrt{\gamma_k} - \sqrt{\gamma_{k+1}}}{\sqrt{\gamma_{k+1}}} = \frac{a}{2k} 
    + \mathcal{O}\left(\frac{1}{k^2}\right) , 
\end{align*}
we obtain that the drift can be expanded as
\begin{align*}
    \sqrt{\frac{\gamma_k}{\gamma_{k+1}}} 
    \left(\identity - \gamma_k \nabla \bar{f}(x^*)\right)     &=
    \left(1 + \frac{a}{2k} + \mathcal{O}\left(\frac{1}{k^2}\right)  \right) \left(\identity - \gamma_k \nabla \bar{f}(x^*) \right) 
    ,
\end{align*}
which can be written as
\begin{align*}
    \left\lVert \sqrt{\frac{\gamma_k}{\gamma_{k+1}}} 
    \left(\identity - \gamma_k \nabla \bar{f}(x^*)\right)      -
    \left(\identity - \gamma_k \nabla \bar{f}(x^*) + \frac{a}{2k} \identity \right) + \frac{a}{2k}   \gamma_k \nabla \bar{f}(x^*) \right\rVert 
    = \mathcal{O}\left(\frac{1}{k^2} \right) .
\end{align*}
When $a < 1$ so that $\gamma_k \gg a/2k$ for large $k$, we obtain 
\begin{equation}\label{eq:a_le_1}
    \left\lVert \sqrt{\frac{\gamma_k}{\gamma_{k+1}}} \left(\identity - \gamma_k \nabla \bar{f}(x^*)\right) 
    - 
    \left(\identity - \gamma_k \nabla \bar{f}(x^*) \right) 
    \right\rVert 
    \leq \frac{a}{2k} \identity + \mathcal{O}\left(\frac{\gamma_k}{k}\right)    
\end{equation}
and when $a = 1$,
\begin{equation}\label{eq:a_eq_1}
    \left\lVert 
    \sqrt{\frac{\gamma_k}{\gamma_{k+1}}} \left(\identity - \gamma_k \nabla \bar{f}(x^*)\right) 
    - 
    \left( \identity - \gamma_k \left(\nabla \bar{f}(x^*) - \frac{a}{2k \gamma_k} \right) \right)
    \right\rVert 
    \leq \frac{\gamma_k}{2k} \lVert \nabla \bar{f}(x^*) \rVert + \mathcal{O}\left(\frac{1}{k^2}\right) .    
\end{equation}
We defined $E_k^a$ in \eqref{eq:Eka_def} as
\begin{align*}
     E_k^a = \left\lVert \sqrt{\frac{\gamma_k}{\gamma_{k+1}}} \left( \identity - \gamma_k \nabla \bar{f}(x^*) \right) - \left(\identity - \gamma_k \bar{A}_a \right) \right\rVert     ,
\end{align*}
where \eqref{eq:a_le_1} demonstrates $E_k^a = \mathcal{O}(1/k)$ when $a < 1$ and \eqref{eq:a_eq_1} demonstrates $E_k^a = \mathcal{O}(1/k^2)$ when $a = 1$.
Substituting $\bar{A}_a$ for the recursion for $\{y_k\}$ in \eqref{eq:yk_draft1}, we obtain 
\begin{equation}
    y_{k+1} = y_k - \gamma_k \bar{A}_a y_k + \sqrt{\gamma_k} M_k + R(\gamma_k, y_k, \xi_k, W_k)    ,
\end{equation}
where
\begin{equation}\label{eq:Rk}
\begin{split}
    R(\gamma_k, y_k, \xi_k, W_k)
    &    
    = 
    \underbrace{\left[ 
    \sqrt{\frac{\gamma_k}{\gamma_{k+1}}}
    \left(\identity - \gamma_k \nabla \bar{f}(x^*) \right)
    - \left( 
    \identity - \gamma_k \bar{A}_a
    \right)
     \right]y_k }_{(a)}
    \\ &
    - \underbrace{\sqrt{\frac{\gamma_k}{\gamma_{k+1}}} \cdot \gamma_k (\nabla \bar{f}(\tilde{x}_k) - \nabla \bar{f}(x^*)) y_k}_{(b)}
    \\ &
    + \underbrace{\sqrt{\gamma_k} \left(\frac{\sqrt{\gamma_k} - \sqrt{\gamma_{k+1}}}{\sqrt{\gamma_{k+1}}} \right)M_k}_{(c)}  
    \\ &
    - \underbrace{\frac{\gamma_k}{\sqrt{\gamma_{k+1}}} 
    \left((P\Phi)(\xi_{k-1}) - (P\Phi)(\xi_{k})\right) }_{(d)} 
    \\ &
    - \underbrace{\frac{\gamma_k}{\sqrt{\gamma_{k+1}}} \left(g(x_k, \xi_k) - g(x^*, \xi_k)\right)}_{(e)}
    . 
\end{split}
\end{equation}
This completes the recursion \eqref{eq:yk}, and (a)--(e) are analyzed in \Cref{app:general_remainder}. 



The sub-sequence \eqref{eq:ytildem} is obtained by summing over the interval $[k_m, k_{m+1})$ as
\begin{align*}
    y_{k_{m+1}} &= y_{k_m} - \sum_{k=k_m}^{k_{m+1} - 1} \gamma_k \bar{A} y_k + \sum_{k=k_m}^{k_{m+1} - 1} \sqrt{\gamma_k} M_k + \sum_{k=k_m}^{k_{m+1} - 1} R(\gamma_k, y_k, \xi_k, W_k) 
    \\&
    = y_{k_m} - \cumstep_m \bar{A} y_{k_m} + \sqrt{\cumstep_m \Gamma} \hat{Z}_m + \sum_{k=k_m}^{k_{m+1} - 1} R(\gamma_k, y_k, \xi_k, W_k)
    \\&
    + \left(\cumstep_m \bar{A} y_{k_m} - \sum_{k=k_m}^{k_{m+1} - 1} \gamma_k \bar{A} y_k \right)
 ,
\end{align*}
from which we define the remainder $\tilde{R}_m$ of the sub-sequence as 
\begin{equation}\label{eq:Rtildem}
\tilde{R}_m = \sum_{k=k_m}^{k_{m+1} - 1} R(\gamma_k, y_k, \xi_k, W_k)
+ \underbrace{\left(\cumstep_m \bar{A} y_{k_m} - \sum_{k=k_m}^{k_{m+1} - 1} \gamma_k \bar{A} y_k \right)}_{(f)}
\end{equation}











We state a result that will be useful in our proofs. 
\begin{lemma}\label{lem:yk_variation}
    Under \Crefrange{ass:step_size}{ass:last}, 
    \begin{equation}
        \lVert y_{k+1} - y_k \rVert_{L_{2p}} 
        \leq 
        \sqrt{\gamma_k} (\beta_{2p}^f + \beta_{2p}^W)
        +
        \mathcal{O}\left(\gamma_k \lVert y_k \rVert_{L_{2p}}+ \frac{\sqrt{\gamma_k}}{k} \left( \beta_{2p}^f + \beta_{2p}^W  \right) \right)  
        . 
    \end{equation}
\end{lemma}
\begin{proof}
    By definition of $y_k = \gamma_k^{-1/2}(x_k - x^*)$, we have
    \begin{align*}
        y_{k+1} - y_k = \gamma_{k+1}^{-1/2}(x_{k+1} - x^*) - \gamma_k^{-1/2}(x_k - x^*) . 
    \end{align*}
    Substituting the recursion for $x_{k+1}$ on the right hand side,
    \begin{align*}
        y_{k+1} - y_k 
        &= \gamma_{k+1}^{-1/2}((x_k - x^*) - \gamma_k (f(x_k, \xi_k) - W_k)) - \gamma_k^{-1/2}(x_k - x^*) 
        \\ &
        = (\gamma_{k+1}^{-1/2} - \gamma_k^{-1/2}) (x_k - x^*) - \gamma_{k+1}^{-1/2} \gamma_k (f(x_k, \xi_k) - W_k) . 
    \end{align*}
    By the Minkowski inequality, 
    \begin{align*}
        &\lVert y_{k+1} - y_k \rVert_{L_{2p}} 
        \\ \leq& 
        \lvert \gamma_{k+1}^{-1/2} - \gamma_k^{-1/2} \rvert \sqrt{\gamma_k} \lVert y_k \rVert_{L_{2p}} + \frac{\gamma_k}{\sqrt{\gamma_{k+1}}} \lVert f(x_k, \xi_k) - W_k \rVert_{L_{2p}} 
        \\ 
        \leq& \frac{a}{2k} \lVert y_k \rVert_{L_{2p}} + \left(1 +\mathcal{O}\left(\frac{1}{k}\right) \right)  \sqrt{\gamma_k} \left\lVert f(x_k, \xi_k) - f(x^*, \xi_k) + f(x^*, \xi_k) - W_k \right\rVert_{L_{2p}}
        ,
    \end{align*}
    where we used the mean value theorem to obtain $\lvert \gamma_{k+1}^{-1/2} - \gamma_k^{-1/2} \rvert \leq \frac{a}{2 k \sqrt{\gamma_k}}$ and $\gamma_k/\sqrt{\gamma_{k+1}} = (1 + \mathcal{O}(1/k)) \sqrt{\gamma_k}$. 
    Applying the Minkowski inequality and Lipschitz continuity of $f$, 
    \begin{align*}
        \lVert y_{k+1} - y_k \rVert_{L_{2p}}
        &\leq \frac{a}{2k} \lVert y_k \rVert_{L_{2p}} + \left(1 + \mathcal{O}\left(\frac{1}{k}\right)\right) \sqrt{\gamma_k} \left( \sqrt{\gamma_k} \mathrm{Lip}(f) \lVert y_k \rVert_{L_{2p}} + \beta_{2p}^f + \beta_{2p}^W  \right) , 
    \end{align*}
    where we used $\lVert f(x^*, \xi_k) \rVert \leq \beta_{2p}^f$. 
    Pulling out the dominant term $\sqrt{\gamma_k}(\beta_{2p}^f + \beta_{2p}^W)$, we obtain the result. 
\end{proof}





\section{Supplement to Section \ref{sec:general}}





\subsection{Proof of Lemma \ref{lem:general_remainder}}\label{app:general_remainder}
Here we obtain the rate of decay of the remainder $\lVert \tilde{R}_m\rVert_{L_p}$ by obtaining bounds on terms (a)--(e) in \eqref{eq:Rk} and (f) in \eqref{eq:Rtildem}. 
The proof reveals that the bound holds for partial sums of the form
\begin{align*}
    \max_{k \in [k_m, k_{m+1})} \Lp{\sum_{j=k_m}^{k}
    R(\gamma_j, y_j, \xi_j, W_j)
    } 
\end{align*}
of (a)--(e). 
The proof is the same, and we present the proof only for the sum over the interval $[k_m, k_{m+1})$. 
The bound for the partial sum is used to obtain a bound for the (f) term in \eqref{eq:Rtildem}. 



We apply the Minkowski inequality and Taylor's remainder theorems to (a) and (b).
\textbf{Term (a):}
A bound on the norm of (a) was obtained in \eqref{eq:a_le_1} and \eqref{eq:a_eq_1}.
By sub-additivity of $\lVert \cdot \rVert$,
\begin{align*}
    \left\lVert \sum_{k=k_m}^{k_{m+1} -1} (a) \right\rVert_{L_p}
    \leq \sum_{k=k_m}^{k_{m+1} - 1} E_k^a \lVert y_k \rVert_{L_p}
    \leq 
    I_m E_{k_m}^{a}\max_{k \in [k_m, k_{m+1})]} \lVert y_k \rVert_{L_p} 
    .
\end{align*}
Substituting \eqref{eq:yn_interval_bound}, we have
\begin{equation}\label{eq:abound}
    \left\lVert \sum_{k=k_m}^{k_{m+1} -1} (a) \right\rVert_{L_p}
    \leq 
    I_m E_{k_m}^{a}  \beta_p^y (m) \lVert \Sigma_a^{1/2} Z \rVert_{L_p} 
    .
\end{equation}
 


\textbf{Term (b):}
By Lipschitz continuity of $\nabla \bar{f}$, 
\begin{align*}
    \lVert  (\nabla\bar{f}(\tilde{x}_k) - \nabla \bar{f}(x^*)) y_k \rVert_{l_p}
    &\leq 
    \lip(\nabla \bar{f}) \mathbb{E}^{1/p}\left(\lVert \tilde{x}_k - x^* \rVert^p \lVert y_k \rVert^p \right)
    \\ &
    = \sqrt{\gamma_k} \lip(\nabla \bar{f}) \lVert y_k \rVert_{L_{2p}}^2 
    . 
\end{align*}
Using $(1 + x)^r \leq 1 + r x$ for $x > 0$ and $r \in [0, 1]$ to obtain $\sqrt{\gamma_k/\gamma_{k+1}} \leq 1 + a/2k$, we apply the Minkowski inequality to obtain
\begin{equation}\label{eq:bbound}
\begin{split}
    \Lp{\sum_{k=k_m}^{k_{m+1} - 1} (b)}
    \leq  
    \left(1 + \mathcal{O}\left( \frac{1}{k_m}\right) \right) \lip(\nabla \bar{f}) \lVert \Sigma_a^{1/2} Z \rVert_{L_{2p}} 
    \cdot \left(\beta_{2p}^y (m) \right)^2 \sqrt{\gamma_{k_m}}\cumstep_m
     .
\end{split}
\end{equation}


\textbf{Term (c):}
Let $\alpha_k = \sqrt{\frac{\gamma_k}{\gamma_{k+1}}} (\sqrt{\gamma_k} - \sqrt{\gamma_{k+1}})$. 
Since $\{M_k\}$ is a MDS, the BR inequality yields
\begin{align*}
    \Lp{\sum_{k=k_m}^{k_{m+1} - 1} \alpha_k M_k } 
    &\leq 
    C_p^{BR} \rosenthal{M} \sqrt{\sum_{k=k_m}^{k_{m+1} - 1} \alpha_k^2} . 
\end{align*}
Using the mean-value theorem 
\begin{equation*}
    \lvert \sqrt{\gamma_k} - \sqrt{\gamma_{k+1}} \rvert 
    \leq 
    \frac{a}{2} \sup_{t \in [k, k+1]} \frac{\sqrt{\gamma_t}}{t}
    = 
    \frac{a}{2k} \sqrt{\gamma_k} ,    
\end{equation*}
and that $\gamma_k/\gamma_{k+1}$ is monotonically decreasing, we obtain 
\begin{align*}
    \sum_{k=k_m}^{k_{m+1} - 1} \alpha_k^2 
    = \sum_{k=k_m}^{k_{m+1}-1} \frac{\gamma_k}{\gamma_{k+1}} \left\lvert \sqrt{\gamma_k} - \sqrt{\gamma_{k+1}}\right\rvert^2
    &\leq 
    \frac{a^2}{4} \sum_{k=k_m}^{k_{m+1} - 1} \frac{\gamma_{k}}{\gamma_{k+1} } \cdot \frac{\gamma_k}{k^2} 
    \\ &\leq\frac{a^2}{4} \frac{\gamma_{k_m}}{\gamma_{k_{m+1}}} \frac{1}{k_m^2} \sum_{k=k_m}^{k_{m+1} - 1} \gamma_k 
    \\ &=\frac{a^2}{4} \frac{\gamma_{k_m}}{\gamma_{k_{m+1}}} \cdot \frac{\cumstep_m}{k_m^2}
    \\ &
    \leq a^2 \frac{\cumstep_m}{k_m^2}
    .
\end{align*}
Therefore, the norm of the weighted sum is bounded as
\begin{equation}\label{eq:cbound}
    \Lp{\sum_{k=k_m}^{k_{m+1} - 1} (c)}
    \leq C_p^{BR}\rosenthal{M}\cdot a\frac{\sqrt{\cumstep_m}}{k_m}      
\end{equation}


\textbf{Term (d):}
Let $\alpha_k = \gamma_k/\sqrt{\gamma_{k+1}}$ and $\Psi$ be defined as $\Psi (\xi) = (P \Phi)(\xi)$ for every $\xi$, so that (d) is given by $\alpha_k (\Psi(\xi_{k-1}) - \Psi(\xi_k)) $.
The following identity is obtained from the summation by parts formula:
\begin{equation}\label{eq:d_sumbyparts}
\begin{split}
    \sum_{k=k_m}^{k_{m+1} - 1} \alpha_k (\Psi(\xi_{k-1}) - \Psi(\xi_{k}))
    & =  
    \alpha_{k_m} \Psi(\xi_{k_m - 1})
    - \alpha_{k_{m+1}} \Psi(\xi_{k_{m+1} - 1}) 
    \\ &+ \sum_{k=k_m + 1}^{k_{m+1}}(\alpha_k - \alpha_{k-1}) \Psi(\xi_{k-1}) ,    
\end{split}    
\end{equation}
From the mean value theorem applied to $\alpha(t) = (t+1)^{a/2}/t^a$, we obtain
\begin{align*}
    \left\lvert \alpha(t) - \alpha(t-1) \right\rvert 
    \leq 
    \max_{s \in [t-1, t]} \lvert \alpha'(s) \rvert 
    \leq 
    \frac{2a}{(t-1)^{1 + a/2}} , 
\end{align*}
which implies $\lvert \alpha_k - \alpha_{k-1} \rvert \leq 2a \sqrt{\gamma_{k-1}}/(k-1)$. 
Using that $\alpha_k$ is monotonic with $\alpha_k \leq 2\sqrt{\gamma_k}$, the $L_p$ norm of \eqref{eq:d_sumbyparts} is bounded using the Minkowski inequality as
\begin{align*}
    \Lp{\sum_{k=k_m}^{k_{m+1} - 1} \alpha_k (\Psi(\xi_{k-1}) - \Psi(\xi_{k}))}
    &\leq \beta_p^\Phi \left(2\alpha_{k_m}
    + \sum_{k=k_m+1}^{k_{m+1}} \lvert \alpha_k - \alpha_{k-1} \rvert 
    \right)
    \\ &
    \leq 
    \beta_p^\Phi \left(4 \sqrt{\gamma_{k_m}}
    + 2a \sum_{k=k_m+1}^{k_{m+1}} \frac{\sqrt{\gamma_{k-1}}}{k-1}
    \right)
    \\ &
    \leq 
    4\beta_p^\Phi \left(\sqrt{\gamma_{k_m}} + \frac{a}{2} \frac{\sqrt{\gamma_{k_m}}}{k_m} I_m \right)
    .
\end{align*}
The $\sqrt{\gamma_{k_m}}$ rate dominates, and we have
\begin{equation}\label{eq:dbound}
    \left\lVert \sum_{k=k_m}^{k_{m+1} - 1} (d) \right\rVert_{L_p} 
    \leq \left(1 + \mathcal{O}\left( \frac{I_m}{k_m}\right) \right) 4 \beta_p^\Phi \sqrt{\gamma_{k_m}}  , 
\end{equation}
where $I_m/k_m \to 0$. 


\textbf{Term (e):}
Consider the Taylor expansion (mean value theorem)
\begin{align*}
    g(x_k, \xi_k) - g(x^*, \xi_k) 
    &= \nabla g(\tilde{x}_k, \xi_k) (x_k - x^* )
    \\ 
    &=  \nabla g(x^*, \xi_k) (x_k - x^*)  + (\nabla g(\tilde{x}_k, \xi_k) - \nabla g(x^*, \xi_k)) (x_k - x^* )
    \\ 
    &= (A(\xi_k) - \mathbb{E}A(\xi))  (x_k - x^*)
    +  (\nabla g(\tilde{x}_k, \xi_k) - \nabla g(x^*, \xi_k)) (x_k - x^*)
    \numberthis \label{eq:e_taylor}
    . 
\end{align*}
The absolute value of the last term is bounded using the Lipschitz continuity of $\nabla g$:
\begin{equation}\label{eq:e_last}
    \left\lvert (\nabla g(\tilde{x}_k, \xi_k) - \nabla g(x^*, \xi_k)) (x_k - x^*) \right\rvert 
    \leq \lip (\nabla g) \lVert x_k - x^* \rVert^2
    = \lip(\nabla g) \gamma_k \lVert y_k \rVert^2 .     
\end{equation}
Let $\alpha_k = \gamma_k/\sqrt{\gamma_{k+1}} \leq \sqrt{2 \gamma_k}$. 
Using \eqref{eq:e_taylor}, we express the $L_p$ norm of the sum of $(e)$ as
\begin{align*}
&\Lp{\sum_{k=k_m}^{k_{m+1} - 1}\alpha_k \left(g(x_k, \xi_k) - g(x^*, \xi_k)\right)}
\\ 
\leq&
\Lp{\sum_{k=k_m}^{k_{m+1} - 1} \alpha_k (A(\xi_k) - \mathbb{E}A(\xi))  (x_k - x^*) }
\\ + &
\Lp{\sum_{k=k_m}^{k_{m+1} - 1} \alpha_k (\nabla g(\tilde{x}_k, \xi_k) - \nabla g(x^*, \xi_k)) (x_k - x^*) }
\\ 
\leq &
\Lp{\sum_{k=k_m}^{k_{m+1} - 1} \alpha_k \sqrt{\gamma_k} (A(\xi_k) - \mathbb{E}A(\xi)) y_k}
+
\sqrt{2} \lip (\nabla g) \sqrt{\gamma_{k_m}} \cumstep_m (\beta_{2p}^y(m) \lVert \Sigma_a^{1/2}Z \rVert_{L_{2p}})^2
\numberthis \label{eq:e_high_level}
,
\end{align*}
where we used \eqref{eq:e_last} and $x_k - x^* = \sqrt{\gamma_k} y_k$.
We use the following bound on the first term.
\begin{lemma}\label{lem:e_bound}
    Let $\alpha_k = \gamma_k/\sqrt{\gamma_{k+1}}$.
    Under the assumptions of \Cref{lem:yk_variation}, it holds that
    \begin{equation}\label{eq:e_intermediate}
    \begin{split}
        &\Lp{\sum_{k=k_m}^{k_{m+1} - 1} \alpha_k \sqrt{\gamma_k} (A(\xi_k) - \mathbb{E}A(\xi))  y_k}
        = \mathcal{O}\left(\gamma_{k_m} \sqrt{I_m} \right)
        .
    \end{split}
    \end{equation}
\end{lemma}
\begin{proof}
    The proof is deferred to \Cref{app:e_state_dependent}.
\end{proof}
For $I_m$ increasing, \eqref{eq:e_intermediate} is the dominant rate in \eqref{eq:e_high_level} and therefore
\begin{equation}\label{eq:ebound}
    \left\lVert \sum_{k=k_m}^{k_{m+1} - 1} (e) \right\rVert_{L_p}
    = \mathcal{O}(\gamma_{k_m} \sqrt{I_m})
    .     
\end{equation}



Using the bounds \eqref{eq:abound}, \eqref{eq:bbound}, \eqref{eq:cbound}, \eqref{eq:dbound}, and \eqref{eq:ebound}, we obtain the following bound on the $L_p$ norm of the sum of \eqref{eq:Rk}:
\begin{equation}\label{eq:abcde_detailed}
\begin{split}
    &\Lp{\sum_{k=k_m}^{k_{m+1} - 1} R(\gamma_k, y_k, \xi_k, M_k)}
    \\ 
    \leq & 
        I_m E_{k_m}^{a}  \beta_p^y (m) \lVert \Sigma_a^{1/2} Z \rVert_{L_p} 
    \\ +&
        \left(1 + \mathcal{O}\left( \frac{1}{k_m}\right) \right) \lip(\nabla \bar{f}) \lVert \Sigma_a^{1/2} Z \rVert_{L_{2p}} 
    \cdot \left(\beta_{2p}^y (m) \right)^2 \sqrt{\gamma_{k_m}}\cumstep_m
     \\ 
     + &
      C_p^{BR} \rosenthal{M} \cdot a\frac{\sqrt{\cumstep_m}}{k_m}   
      \\ 
      + & \left(1 + \mathcal{O}\left( \frac{I_m}{k_m}\right) \right) 4 \beta_p^\Phi \sqrt{\gamma_{k_m}}  
      \\ 
    + & 
    \mathcal{O}(\gamma_{k_m} \sqrt{I_m})
    .         
\end{split}
\end{equation}
The rate $\sqrt{\gamma_{k_m}} \cumstep_m$ in the second line of the bound is negligible relative to $\sqrt{\gamma_{k_m}}$, and the third line is $\mathcal{O}(1/n)$ which is negligible. 
The $\gamma_{k_m} \sqrt{I_m} \approx \sqrt{\gamma_{k_m} \cumstep_m}$ rate in the last line decays faster than $\sqrt{\gamma_{k_m}}$, and \eqref{eq:abcde_detailed} can be written as
\begin{equation}\label{eq:abcde}
\begin{split}
    \Lp{\sum_{k=k_m}^{k_{m+1} - 1} R(\gamma_k, y_k, \xi_k, W_k)}
    &\leq 
     \lVert \Sigma_a^{1/2}Z \rVert_{L_p} \cdot I_m E_{k_m}^a \beta_p^y (m)
    +
    4    \beta_p^\Phi \sqrt{\gamma_{k_m}}
    \\ &+
    \mathcal{O}\left( \sqrt{\gamma_{k_m}} \left( \frac{I_m}{k_m} + \sqrt{\cumstep_m}\right) \right)
    .    
\end{split}    
\end{equation}


\textbf{Term (f):}
We start with the identity
\begin{align*}
    \cumstep_m y_{k_m} - \sum_{k=k_m}^{k_{m+1} - 1} \gamma_k y_k 
    &=
    \sum_{k=k_m}^{k_{m+1} - 1} \gamma_k (y_{k_m} - y_k) 
    \\ &
    = -\sum_{k=k_m}^{k_{m+1} - 1} \gamma_k \sum_{j=k_m}^{k-1} \left( \gamma_j \bar{A} y_j - \sqrt{\gamma_j} M_j - R_j
    \right)
    ,
\end{align*}
where $R_j = R(\gamma_j, y_j, \xi_j, W_j)$. 
By the Minkowski inequality, we have
\begin{align*}
    \Lp{    \cumstep_m y_{k_m} - \sum_{k=k_m}^{k_{m+1} - 1} \gamma_k y_k }
    &\leq 
    \Lp{\sum_{k=k_m}^{k_{m+1} - 1}
    \gamma_k \sum_{j=k_m}^{k - 1} \gamma_j \bar{A} y_j 
    }
    +
    \Lp{\sum_{k=k_M}^{k_{m+1} - 1} \gamma_k \sum_{j=k_m}^{k - 1} \sqrt{\gamma_j} M_j }
    \\ &
    + 
    \Lp{\sum_{k=k_m}^{k_{m+1} - 1} \gamma_k \sum_{j=k_m}^{k-1} R_j} . 
\end{align*}
The first term is bounded using the Minkowski inequality, where 
\begin{align*}
        \Lp{\sum_{k=k_m}^{k_{m+1} - 1}
    \gamma_k \sum_{j=k_m}^{k - 1} \gamma_j \bar{A} y_j 
    }
    &\leq\lVert \bar{A} \rVert \sup_{j \in [k_m, k_{m+1})} \lVert y_j \rVert_{L_p} \sum_{k=k_m}^{k_{m+1} - 1} \gamma_k \sum_{j=k_m}^{k-1} \gamma_j 
    \\ &
    \leq \lVert \bar{A} \rVert \lVert \Sigma_a^{1/2} Z \rVert_{L_p} \beta_p^y (m) \cumstep_m^2 
    \numberthis \label{eq:f1}
    . 
\end{align*}
Therefore, $\Lp{\sum_{k=k_m}^{k_{m+1} - 1}
    \gamma_k \sum_{j=k_m}^{k - 1} \gamma_j \bar{A} y_j 
    } = \mathcal{O}(\cumstep_m^2)$. 
The martingale sum is first rearranged to
\begin{align*}
    \Lp{\sum_{k=k_M}^{k_{m+1} - 1} \gamma_k \sum_{j=k_m}^{k - 1} \sqrt{\gamma_j} M_j }
    &=
    \Lp{
    \sum_{j=k_m}^{k_{m+1} - 2} \left( \sum_{k=j+1}^{k_{m+1} - 1} \gamma_k \right) \sqrt{\gamma_j} M_j
    } . 
\end{align*}
Using $\alpha_j = (\sum_{k=j+1}^{k_{m+1} - 1} \gamma_k) \sqrt{\gamma_j}$ and following previous arguments, we obtain
\begin{align*}
    \Lp{\sum_{j=k_m}^{k_{m+1} - 2}\alpha_j M_j }
    \leq 
    C_p^{BR} \rosenthal{M}\sqrt{\sum_{j=k_m}^{k_{m+1} - 1} \alpha_j^2}
    &=
    C_p^{BR} \rosenthal{M}\sqrt{\sum_{j=k_m}^{k_{m+1} - 1}
    \gamma_j \left(\sum_{k=j+1}^{k_{m+1} - 1} \gamma_k \right)^2
    }
    \\ &
    \leq C_p^{BR} \rosenthal{M}\cumstep_m^{3/2} . 
    \numberthis \label{eq:f2}
\end{align*}
Therefore, the second term is of order $\mathcal{O}(\cumstep_m^{3/2})$. 


The last term is evaluated by first rearranging using the summation by parts formula
\begin{align*}
    \sum_{k=k_m}^{k_{m+1} - 1} \gamma_k \sum_{j=k_m}^{k-1} R_j
    = 
    \sum_{j=k_m}^{k_{m+1} - 2} \left(\sum_{k=j+1}^{k_{m+1} - 1} \gamma_k \right) R_j 
    &=
    \sum_{j=k_m}^{k_{m+1} - 2}\left(\sum_{k=j+1}^{k_{m+1} - 1} \gamma_k \right) \left(\sum_{i=k_m}^{j} R_i - \sum_{i=k_m}^{j-1} R_i \right)
    .
\end{align*}
Let $\alpha_j = \sum_{k=j+1}^{k_{m+1} - 1} \gamma_k$ and $B_j = \sum_{i=k_m}^{j} R_i$, so that the above is written as
\begin{align*}
    \sum_{j=k_m}^{k_{m+1} - 2} \alpha_j \left(B_j - B_{j-1} \right) 
    &=
    \alpha_{k_{m+1} - 2} B_{k_{m+1} - 3} - \sum_{j=k_m}^{k_{m+1} - 2} (\alpha_{j+1} - \alpha_j) B_j
    \\ &
    = \gamma_{k_{m+1} - 1} B_{k_{m+1} - 2} + \sum_{j=k_m}^{k_{m+1} - 3} \gamma_{j+1} B_j . 
\end{align*}
Taking the $L_p$ norm and using that the bound \eqref{eq:abcde} holds for the partial sum $\max_{j \in [k_m, k_{m+1})} \lVert B_j \rVert = \mathcal{O}(I_m E_{k_m}^a + \sqrt{\gamma_{k_m}})$ (as discussed in the beginning of \Cref{app:general_remainder}), we obtain 
\begin{align*}
    \Lp{    \sum_{k=k_m}^{k_{m+1} - 1} \gamma_k \sum_{j=k_m}^{k-1} R_j}
    &= \mathcal{O}\left( 
     \left(I_m E_{k_m}^a +  \sqrt{\gamma_{k_m}} \right) \left( \gamma_{k_{m+1} - 1} 
     + \sum_{j=k_m+1}^{k_{m+1} - 2} \gamma_j \right)
     \right)
     \\ &
     = \mathcal{O}\left(\left(I_m E_{k_m}^a +  \sqrt{\gamma_{k_m}} \right) \cumstep_m  \right). 
     \numberthis \label{eq:f3}
\end{align*}
Combining \eqref{eq:f1}--\eqref{eq:f3}, we obtain 
\begin{align*}
    \Lp{(f)}
    &\leq 
    \lVert \bar{A} \rVert \left(\lVert \bar{A} \rVert \lVert \Sigma_a^{1/2}Z \rVert_{L_p} \beta_p^y (m) \cumstep_m^2 + C_p^{BR} \rosenthal{M}\cumstep_m^{3/2} 
    + \mathcal{O}\left( \left(I_m E_{k_m}^a +\sqrt{\gamma_{k_m}}\right) \cumstep_m\right)\right)
    .
\end{align*}
Using that $\sqrt{\gamma_{k_m}} \cumstep_m \ll \cumstep_m^{3/2}$ when $I_m$ is an increasing sequence, we obtain
\begin{equation}\label{eq:f_bound}
    \lVert (f) \rVert_{L_p} 
    \leq C_p^{BR}\rosenthal{M} \lVert \bar{A} \rVert \cumstep_m^{3/2} + \mathcal{O}\left(
    \cumstep_m^2 + (I_m E_{k_m}^a + \sqrt{\gamma_{k_m}} )\cumstep_m \right) . 
\end{equation}
To obtain the result in \Cref{lem:general_remainder}, we combine \eqref{eq:abcde} with \eqref{eq:f_bound} to obtain
\begin{align*}
    \Lp{\tilde{R}_m} 
    &\leq \lVert \Sigma_a^{1/2}Z \rVert_{L_p} \cdot I_m E_{k_m}^a \beta_p^y (m) + 4 \beta_p^{\Phi} \sqrt{\gamma_{k_m}}
    + 
    C_p^{BR} \rosenthal{M}\lVert \bar{A} \rVert \cumstep_m^{3/2}  \\ 
    &
    + \mathcal{O}\left(
    I_N E_n^a \cumstep_m +
    \sqrt{\gamma_{k_m}} \frac{I_m}{k_m} + \sqrt{\gamma_{k_m} \cumstep_m}
    + \cumstep_m^2
    \right)
    .
\end{align*}



\subsection{Proof of \Cref{lem:Wp_recursion}}\label{app:Wp_recursion}
The result is proved by a comparison argument. 
Recall the two recursions
\begin{equation}\label{eq:thm1_dt_recursions}
\begin{split}
    y_{k_{m+1}} &= y_{k_m} - \cumstep_m \bar{A} y_m + \sqrt{\cumstep_m \Gamma} \hat{Z}_m + \tilde{R}_m, 
    \\ 
    u_{m+1} &= u_m - \cumstep_m \bar{A} u_m + \sqrt{\cumstep_m \Gamma} Z_m , \quad u_1 = y_1
\end{split}
\end{equation}
where $\{Z_m\}$ is an iid standard Gaussian sequence and $\Gamma$ is the asymptotic covariance. 
Using the triangle inequality,
\begin{equation}\label{eq:Wp_triangle_inequality}
    \mathcal{W}_{p} (y_{k_{m+1}},  u_{m+1})
    \leq (1 - \lambdaDT \cumstep_m) \mathcal{W}_{p} (y_{k_m},  u_{k_m})
    + \lVert \tilde{R}_m \rVert_{L_p} 
    + \sqrt{\cumstep_m \lVert \Gamma \rVert} \mathcal{W}_{p} (\hat{Z}_m, Z_m),     
\end{equation}
where we used that $\lVert I - \cumstep_m \bar{A} \rVert \leq 1 - \lambdaDT \cumstep_m$. 
By \Cref{lem:martingale_clt_weighted}, we have \eqref{eq:clt_weighted_master} where
\begin{align*}
    \mathcal{W}_{p} (\hat{Z}_m, Z_m) 
    \leq
    \left(\mathbb{E} \mathcal{W}_{p}^{p} (\hat{Z}_m, Z_m | \mathcal{H}_{k_m})\right)^{1/{p}}
    \leq 
     C_{p}^{BR} \rosenthal{M} \cdot a \frac{I_m}{k_m} + \frac{\clt(I_m)}{\sqrt{I_m}} , 
\end{align*}
where we used that the CLT rates in \Cref{app:clt_weighted} hold uniformly over $\mathcal{H}_{k_m}$.
By \Cref{lem:general_remainder}, $\Lp{\tilde{R}_m} = \mathcal{O}(I_m E_{k_m}^a \beta_p^y(m) + \sqrt{\gamma_{k_m}} + \cumstep_m^{3/2})$ decays at a polynomial rate faster than $\cumstep_m$ whenever $\sqrt{\gamma_{k_m}} \ll \cumstep_m$, which requires $I_m \sqrt{\gamma_{k_m}} \to 0$.
For this choice, the recursion in \eqref{eq:Wp_triangle_inequality} is solved using \Cref{lem:recursion}:
\begin{equation}\label{eq:yN_uN}
\begin{split}
    \mathcal{W}_{p} (y_{k_{N+1}}, u_{N+1}) 
    &\leq
    \frac{1}{\lambdaDT \cumstep_N} \left(
    \sqrt{\lVert \Gamma \rVert \cumstep_N} \left(C_{p}^{BR} \rosenthal{M} \cdot a \frac{I_N}{k_N} + \frac{\clt(I_N)}{\sqrt{I_N}} \right)
    + \lVert \tilde{R}_N \rVert_{L_p}
    \right)    
    \\ &
    + \mathcal{O}\left(\frac{1}{\cumstep_N^2 k_N} \left( \sqrt{\cumstep_N} \left(\frac{I_N}{k_N} + \frac{ \clt (I_N)}{\sqrt{I_N}} \right) + \lVert \tilde{R}_N \rVert_{L_p} \right) \right)
\end{split}
\end{equation}
Let $\tau_1 = 0$ and consider the interpolant times $\tau_{m+1} = \tau_m + \cumstep_m$. 
The distance between the discrete-time OU process and its continuous time counterpart 
\begin{equation}\label{eq:ct_ou_sde}
    dU_t = - \bar{A}_a U_t dt + \sqrt{\Gamma} dB_t , \quad U_{\tau_1} \overset{d.}{=} u_1
\end{equation}
is described.
\begin{lemma}\label{lem:ou_discretization}
    For $p \geq 1$ such that $\lVert U_{\tau_1} \rVert_{L_p}$ is finite, there exists a constant $C_p^{BDG} > 0$ such that the solutions to \eqref{eq:thm1_dt_recursions} and \eqref{eq:ct_ou_sde} at interpolation times $\tau_{m+1} - \tau_m = \cumstep_m$ satisfy
    \begin{equation}\label{eq:ou_discretization}
    \mathcal{W}_p (u_{N+1}, U_{\tau_{N+1}}) 
    \leq 
    K \frac{C_p^{BDG}}{\lambdaDT} 
    \frac{\lVert \bar{A}_a \rVert }{\lambdaCT} \sqrt{\frac{\lVert Q \rVert_2\mathrm{Tr} \Gamma}{3}} \sqrt{\cumstep_m}
    + \mathcal{O}\left(\cumstep_m \right) 
    .
    \end{equation}
\end{lemma}
\begin{proof}
    See Section \ref{app:ou_discretization}.
\end{proof}
The convergence of the continuous time OU process to its stationary distribution is proved for completeness.
A simple calculation shows that the stationary limit is a Gaussian distribution $\mathcal{N}(0, \Sigma_a)$, where $\Sigma_a$ is defined in \eqref{eq:yn_asymptotic_clt}.
\begin{lemma}\label{lem:ct_ou_convergence}
    Let $U_\infty \sim \mathcal{N}(0, \Sigma_a)$. 
    There exists a constant $K > 0$ such that the solution to \eqref{eq:ct_ou_sde} with initial condition $\lVert U_0 \rVert_{L_p} <\infty $ satisfies
    \begin{equation}\label{eq:ct_ou_convergence}
        \mathcal{W}_p \left(U_t, U_\infty  \right)
        \leq 
        K e^{-\lambdaCT t} \mathcal{W}_p (U_0, U_\infty) . 
    \end{equation}
\end{lemma}
\begin{proof}
    See Section \ref{app:ct_ou_convergence}. 
\end{proof}
\textbf{Convergence rate:}
Combining Eqs. \eqref{eq:yN_uN}, \eqref{eq:ou_discretization}, and \eqref{eq:ct_ou_convergence} with $U_0 = y_1$, we obtain
\begin{equation}
\begin{split}
    \mathcal{W}_{p} (y_{k_{N+1}}, U_\infty)
    &\leq 
    \mathcal{W}_{p} (U_{\tau_{N+1}}, U_\infty) + \mathcal{W}_{p} (u_{N+1}, U_{\tau_{N+1}})  + \mathcal{W}_{p} (y_{k_{n+1}}, u_{N+1}) 
    \\
    &\leq
    K e^{-\lambdaCT \sum_{m=1}^N \cumstep_m} \mathcal{W}_p \left(y_1, U_\infty \right) 
    +  K\frac{C_p^{BDG}}{\lambdaDT} 
    \frac{\lVert \bar{A}_a \rVert }{\lambdaCT} \sqrt{\frac{\lVert Q \rVert_2\mathrm{Tr} \Gamma}{3}} \sqrt{\cumstep_N}
    \\ &
    + \frac{1}{\lambdaDT \cumstep_N} \left(
    \sqrt{\lVert \Gamma \rVert \cumstep_N} \left(C_{p}^{BR} \rosenthal{M} \cdot a \frac{I_N}{k_N} + \frac{\clt(I_N)}{\sqrt{I_N}} \right)
    + \lVert \tilde{R}_N \rVert_{L_p}
    \right)    
    \\ &
    + \mathcal{O}\left(\cumstep_N 
    +
    \frac{1}{\cumstep_N^2 k_N} \left( \sqrt{\cumstep_N} \left(\frac{I_N}{k_N} + \frac{ \clt (I_N)}{\sqrt{I_N}} \right) + \lVert \tilde{R}_N \rVert_{L_p} \right) 
    \right)
\end{split}
\end{equation}



\subsection{Proof of Theorem \ref{thm:general}}\label{app:general}
We substitute for $\lVert \tilde{R}_N \rVert_{L_p}$ the bound in \Cref{lem:general_remainder}. 
The $I_N$ in \Cref{lem:Wp_recursion} is chosen to (approximately) minimize the error bound. 
Substituting the bound in \Cref{lem:general_remainder} for the error bound in \Cref{lem:Wp_recursion}, we obtain
\begin{equation}\label{eq:thm1_all_terms}
\begin{split}
    \mathcal{W}_{1}(y_n, U_\infty) 
    &\leq 
    K e^{-\lambdaCT \sum_{m=1}^N \cumstep_m} \mathcal{W}_1 \left(y_1, U_\infty \right) 
    + K\frac{C_p^{BDG}}{\lambdaDT} 
    \frac{\lVert \bar{A}_a \rVert }{\lambdaCT} \sqrt{\frac{\lVert Q \rVert_2\mathrm{Tr} \Gamma}{3}} \sqrt{\cumstep_N}
    \\ &
    + \frac{1}{\lambdaDT} 
    \sqrt{\frac{\lVert \Gamma \rVert}{\cumstep_N}}  \left( C_p^{BR} \rosenthal{\Gamma^{-1/2} M} \cdot a \frac{I_m}{k_m} + \frac{\clt(I_m)}{\sqrt{I_m}} \right)
    \\ &
    + \frac{1}{\lambdaDT \cumstep_N} \left(          \lVert \Sigma_a^{1/2}Z \rVert_{L_p} \cdot I_N E_{k_N}^a \beta_p^y (N)
    +
    4    \beta_p^\Phi \sqrt{\gamma_{k_N}}
    + 
        C_p^{BR} \rosenthal{ M} \lVert \bar{A}_a \rVert \cumstep_N^{3/2} \right)
    \\ &
    + \mathcal{O}\left(\cumstep_N 
    +
    \frac{1}{\cumstep_N^2 k_N} \left( \sqrt{\cumstep_N} \left(\frac{I_N}{\sqrt{\cumstep_N} k_N} + \frac{ \clt(I_N)}{\sqrt{\gamma_{k_N}} I_N} \right) + \lVert \tilde{R}_N \rVert_{L_p} \right)
    \right)
    \\ &
    +     
    \mathcal{O}\left( \frac{1}{\cumstep_N} \left( \gamma_{k_N} I_N E_{k_N}^a + \sqrt{\gamma_{k_N}} \left( \cumstep_N + \frac{I_N}{k_N} + \sqrt{\cumstep_N}\right) 
    + I_N^2 \gamma_{k_N}^2
    \right)\right)     
\end{split}
\end{equation}
The $\mathcal{O}(\cdot)$ terms were previously verified to be asymptotically negligible when $I_m$ is increasing at a rate such that $\cumstep_m \approx \gamma_{k_m} I_m$ is decreasing.
Ignoring the transient big-$\mathcal{O}$ terms, we obtain that the bound is dominated by
\begin{equation}\label{eq:thm1_dominant}
\begin{split}
     &K\frac{C_p^{BDG}}{\lambdaDT} 
    \frac{\lVert \bar{A}_a \rVert }{\lambdaCT} \sqrt{\frac{\lVert Q \rVert_2\mathrm{Tr} \Gamma}{3}} \sqrt{\cumstep_N}
     + \frac{1}{\lambdaDT} 
    \sqrt{\frac{\lVert \Gamma \rVert}{\cumstep_N}} \left( C_p^{BR} \rosenthal{\Gamma^{-1/2} M} \cdot a \frac{I_m}{k_m} + \frac{\clt(I_m)}{\sqrt{I_m}}\right)
    \\  +& 
     \frac{1}{\lambdaDT \cumstep_N} \left(          \lVert \Sigma_a^{1/2}Z \rVert_{L_p} \cdot I_N E_{k_N}^a \beta_p^y (N)
    +
    4    \beta_p^\Phi \sqrt{\gamma_{k_N}}
    + 
        C_p^{BR} \rosenthal{M}\lVert \bar{A}_a \rVert \cdot \cumstep_N^{3/2} \right) 
        .    
\end{split}    
\end{equation}
Since $\cumstep_N \sim \gamma_{k_N} I_N$, the $I_N E_{k_N}^a/\cumstep_N$ does not affect the choice of $I_N$. 
The second rate $I_N/k_N\sqrt{\cumstep_N}$ is of order $n^{-1 + a/2} I_N^{1/2} \ll \sqrt{\cumstep_N}$ and is negligible. 
Grouping the dominant terms by their rates and neglecting the common term $1/\lambdaDT$, we aim to optimize over $I_N$ the bound
\begin{equation}\label{eq:thm1_dominant_step2}
\begin{split}
    &\sqrt{\gamma_n I_N} \left(  \lVert \bar{A}_a \rVert \left(
    K\frac{C_p^{BDG}}{\lambdaDT \lambdaCT} \sqrt{\frac{\lVert Q \rVert_2\mathrm{Tr} \Gamma}{3}} 
    + C_p^{BR} \rosenthal{M}\right)\right)
    \\ + & 
    \frac{1}{\sqrt{\gamma_n} I_N} \left(\sqrt{\lVert \Gamma \rVert} \clt(I_N) + 4 \beta_p^{\Phi} \right)
    .         
\end{split}
\end{equation}
The bound is optimized over $I_N$, treating $\clt(I_N)$ as a constant $\dummy$ independent of $I_N$, which yields
\begin{equation}\label{eq:INstar}
    I_N^* = \left(\frac{\sqrt{\lVert \Gamma \rVert} \dummy  + 4 \beta_p^{\Phi}}{\lVert \bar{A}_a \rVert \left(K\frac{C_p^{BDG}}{\lambdaDT \lambdaCT} \sqrt{\frac{\lVert Q \rVert_2\mathrm{Tr} \Gamma}{3}}  + C_p^{BR} \rosenthal{M} \right)} \right)^{2/3} \gamma_n^{-2/3} .     
\end{equation}
For simplicity, we treat $I_N^*$ as a continuous variable during the optimization; since the optimal choice grows with $n$, rounding it to the nearest integer does not affect the asymptotic rate of the bound. 
Using this choice of $I_N^*$, we obtain that the error bound in \Cref{lem:Wp_recursion} is 
\begin{equation}\label{eq:optimized_Wp}
\begin{split}
    \mathcal{W}_p (y_n, U_\infty) 
    &\leq
    \frac{2}{\lambdaDT} \left(\lVert \bar{A}_a \rVert \left(K\frac{C_p^{BDG}}{\lambdaDT \lambdaCT} \sqrt{\frac{\lVert Q \rVert_2\mathrm{Tr} \Gamma}{3}}  + C_p^{BR} \rosenthal{M} \right) \right)^{2/3}  
    \\ &
    \cdot \left( \sqrt{\lVert \Gamma \rVert} \clt(I_N^*) + 4 \beta_p^\Phi\right)^{1/3} \gamma_n^{1/6}
    \\ 
    &+
    \frac{\lVert \Sigma_a^{1/2}Z\rVert_{L_p}}{\lambdaDT} \beta_p^y (N) \frac{E_{k_N}^a}{\gamma_n}
    + 
    \mathcal{O}\left( \gamma_n^{1/3} \right).     
\end{split}
\end{equation}
The bounds in \Cref{thm:general} are obtained by instantiating \eqref{eq:optimized_Wp} for the following special cases. 


\textbf{General Noise:}
The first part uses the martingale CLT \Cref{lem:martingale_clt}: for a universal constant $C>0$, the bound holds for $p = 1$ and $\mathrm{CLT}(I_N^*) = \mathcal{O}(\log \gamma_n^{-1})$.
Substituting, we obtain
\begin{align*}
    \mathcal{W}_1 (y_n, U_\infty) 
    = \mathcal{O}\left(\gamma_n^{1/6} \left(\log \gamma_n \right)^{1/3} + \frac{E_n^a}{\gamma_n} \right) .
\end{align*}
To ensure our choice of $I_N^*$ is valid, we must verify that $I_N^* \in [1, n]$ for sufficiently large $n$. 
The choice of $I_N^*$ is a solution to an equation of the form 
\begin{align*}
    I_N^* = \left(D_1 \log I_N^* + D_2 \right)^{2/3} \gamma_n^{-2/3} , 
\end{align*}
where $D_1, D_2$ are problem dependent constants independent of $n$. 
When $n$ is sufficiently large, then 
\begin{align*}
    \lim_{n \to \infty} \frac{\log I_N^*}{\log \gamma_n^{-1}} 
    = 
    \frac{2}{3}\lim_{n \to \infty} \left\{ \frac{\log \left( D_1 \log I_N^* + D_2\right)}{\log \gamma_n^{-1}} + 1\right\}
    = \frac{2}{3} , 
\end{align*}
which confirms that the choice of $I_N^*$ is valid when $n$ is sufficiently large. 




\textbf{Independent Noise:}
When $\{W_k\}$ and $\{\xi_k\}$ are independent sequences and satisfy the conditions in \Cref{lem:bonis_weighted}, then \Cref{ass:CLT} holds with $\mathrm{CLT}(I_N) = C_p \sqrt{\lVert Q \rVert_2} (d^{1/4} (\beta_{p+2}^{\Gamma^{-1/2}M})^2 + (\beta_{p+2}^{\Gamma^{-1/2}M})^{1+2/p})$, which is independent of $I_N$.
Substituting for \eqref{eq:INstar}, we obtain that for any $p \geq 2$, \eqref{eq:optimized_Wp} becomes
\begin{align*}
    \mathcal{W}_p (y_n, U_\infty) 
    &\leq
    \frac{1}{\lambdaDT} \left(\lVert \bar{A}_a \rVert \left(K\frac{C_p^{BDG}}{\lambdaDT \lambdaCT} \sqrt{\frac{\lVert Q \rVert_2\mathrm{Tr} \Gamma}{3}}  + C_p^{BR} \beta_p^M \right) \right)^{2/3}  
    \\ 
    &\cdot 
    \left( C_p \sqrt{\lVert Q \rVert_2 \lVert \Gamma \rVert}
        \left(d^{1/4} (\beta_4^{\Gamma^{-1/2}M})^2 + (\beta_{p+2}^{\Gamma^{-1/2}M})^{1+2/p} \right)  
        + 4 \beta_p^\Phi  \right)^{1/3} \gamma_n^{1/6}
    \\ 
    &+
    \frac{\lVert \Sigma_a^{1/2}Z\rVert_{L_p}}{\lambdaDT} \beta_p^y (N) \frac{E_{k_N}^a}{\gamma_n}
    + 
    \mathcal{O}\left( \gamma_n^{1/3} \right).     
\end{align*}
It is easy to check that \eqref{eq:INstar} lies within the interval $[1, n]$ for all sufficiently large $n$, hence is a valid choice. 
Moreover, the solution $\Phi$ to the Poisson equation $f(x^*, \xi)  = (I - P)\Phi(\xi) $ simplifies to $f(x^*, \xi)$ for i.i.d. noise, and the result is presented using $\beta_p^f$ in place of $\beta_p^\Phi$.







\subsection{Proof of Theorem \ref{thm:pr_average}}\label{app:pr_average}
Note that $a < 1$ is assumed for Polyak-Ruppert averaging, and so $\bar{A} = \nabla f(x^*)$. 
Recall the recursion \eqref{eq:sa}, which is rearranged as
\begin{align*}
    x_{k+1} &=
    x_k - \gamma_k \bar{A}(x_k - x^*) 
    + \gamma_k (\bar{A}(x_k - x^*) - \bar{f}(x_k))
    + \gamma_k (\bar{f}(x_k) - f(x_k, \xi_k))
    - \gamma_k W_k .
\end{align*}
For the analysis of Polyak-Ruppert averaging, it is convenient to rearrange the recursion as
\begin{align*}
    \gamma_k \bar{A} (x_k - x^*) = 
    \left[x_k - x_{k+1} \right] + \gamma_k \left[\bar{A}(x_k - x^*) - \bar{f}(x_k)  \right] 
    + \gamma_k \left[ \bar{f}(x_k) - f(x_k, \xi_k)\right]
    - \gamma_k W_k
    . 
\end{align*}
Substituting $g(x_k, \xi_k) = f(x_k, \xi_k) - \bar{f}(x_k)$ defined in \eqref{eq:gk} and dividing by $\gamma_k$, we obtain
\begin{align*}
    \bar{A}(x_k - x^*) 
    &=
    \gamma_k^{-1}(x_k - x_{k+1}) 
    + (\bar{A}(x_k - x^*) - \bar{f}(x_k))
    + g(x_k, \xi_k) 
    - W_k
    \\ &=
    M_k + \gamma_k^{-1}(x_k - x_{k+1}) 
    + (\bar{A}(x_k - x^*) - \bar{f}(x_k))
    - (g(x_k, \xi_k) - g(x^*, \xi_k))
    \\ &
    - (g(x^*, \xi_k) + M_k + W_k)
    . 
\end{align*}
Summing over $k = 1, 2, \cdots, n$ and normalizing by $\sqrt{n}$, we obtain
\begin{equation}\label{eq:pr_decomposition}
\begin{split}
    \bar{A} \bar{y}_n 
    &= 
     \underbrace{\frac{1}{\sqrt{n}} \sum_{k=1}^n M_k }_{(a)}
     +
    \underbrace{\frac{1}{\sqrt{n}}  \sum_{k=1}^{n} \gamma_k^{-1}(x_k - x_{k+1}) }_{(b)}
    + 
    \underbrace{\frac{1}{\sqrt{n}} \sum_{k=1}^n \left(\bar{A} (x_k - x^*) - \bar{f}(x_k) \right)}_{(c)}
    \\ &
    - \underbrace{\frac{1}{\sqrt{n}} \sum_{k=1}^n (g(x_k, \xi_k) - g(x^*, \xi_k))}_{(d)} 
    -
    \underbrace{\frac{1}{\sqrt{n}} \sum_{k=1}^n (g(x^*, \xi_k) + M_k + W_k)}_{(e)}
    .     
\end{split}    
\end{equation}
The terms (b)--(e) are analyzed separately. 


\textbf{Term (b):}
We apply the summation by parts formula to obtain
\begin{align*}
    \frac{1}{\sqrt{n}} \sum_{k=1}^n \gamma_k^{-1}(x_k - x_{k+1})
    &= \sqrt{\frac{1}{\gamma_1 n} } y_1 - \sqrt{\frac{1}{\gamma_{n+1}(n+1)}} y_{n+1} + \frac{1}{\sqrt{n}} \sum_{k=1}^n (\gamma_k^{-1} - \gamma_{k+1}^{-1})\sqrt{\gamma_k} y_k         .
\end{align*}
By the Minkowski inequality, 
\begin{align*}
    \Lp{\frac{1}{\sqrt{n}} \sum_{k=1}^n \gamma_k^{-1}(x_k - x_{k+1})}
    &
    \leq 
    \frac{1}{\sqrt{\gamma_n n}} \lVert y_{n+1} \rVert_{L_p} + 
    \frac{\sup_{k \leq n} \lVert y_k \rVert_{L_p}}{\sqrt{n}} \sum_{k=1}^n \lvert \gamma_k^{-1} - \gamma_{k+1}^{-1}\rvert \sqrt{\gamma_k} 
    \\ &
    + \mathcal{O}\left(\frac{1}{\sqrt{n}} \right) . 
\end{align*}
Using $\lvert \gamma_k^{-1} - \gamma_{k+1}^{-1} \rvert \leq (\gamma_k k)^{-1}$ and $\sum_{k=1}^n k^{-1}/\sqrt{\gamma_k} \leq 1/\sqrt{\gamma_k} + \mathcal{O}(1/n\sqrt{\gamma_k})$, we obtain
\begin{equation}\label{eq:pr_b_bound}
    \Lp{\frac{1}{\sqrt{n}} \sum_{k=1}^n \gamma_k^{-1}(x_k - x_{k+1})}
    \leq 
    \frac{\max_{k \leq n+1} \lVert y_k \rVert_{L_p}}{\sqrt{\gamma_n n}}
    + \mathcal{O}\left(\frac{1}{\sqrt{n}}\right) . 
\end{equation}

\textbf{Term (c):}
We use Taylor's remainder theorem, where for $\tilde{x} = \theta x + (1-\theta) x^*$ with $\theta \in [0, 1]$,
\begin{align*}
    \bar{f}(x) &= \bar{f}(x^*) + \nabla \bar{f}(\tilde{x}) (x - x^*)
    \\ &
    = \bar{f}(x^*) + \nabla \bar{f}(x^*)^\dagger (x - x^*)
    + (\nabla \bar{f}(\tilde{x}) - \nabla \bar{f}(x^*))^\dagger (x - x^*). 
\end{align*}
Using uniform Lipschitz continuity of $\bar{f}$ in \Cref{ass:general_sa} and $\bar{f}(x^*) = 0$, we obtain
\begin{equation}\label{eq:remainder_barf}
    \left\lVert \bar{f}(x_k) - (\bar{A}(x_k - x^*)) \right\rVert 
    \leq \lip(\nabla \bar{f}) \lVert x_k - x^* \rVert^2 .     
\end{equation}
Therefore,
\begin{align*}
    \Lp{\frac{1}{\sqrt{n}} \sum_{k=1}^n (\bar{A}(x_k - x^*) - \bar{f}(x_k))}
    &\leq 
    \frac{1}{\sqrt{n}}\lip(\nabla \bar{f}) \sum_{k=1}^n \gamma_k \lVert y_k \rVert_{L_{2p}}^2     
    \\ &
    \leq \frac{\lip(\nabla \bar{f}) (\sup_{k \leq n} \lVert y_k \rVert^2_{L_{2p}})}{1-a} n^{1/2} \gamma_n + \mathcal{O}\left(\frac{1}{\sqrt{n}}\right). 
\end{align*}
To summarize,
\begin{equation}\label{eq:pr_c_bound}
    \Lp{(c)} \leq 
    \frac{\lip(\nabla \bar{f}) (\max_{k \leq n} \lVert y_k \rVert^2_{L_{2p}})}{1-a} n^{1/2} \gamma_n + \mathcal{O}\left(\frac{1}{\sqrt{n}}\right). 
\end{equation}

\textbf{Term (d):}
Following the steps in \eqref{eq:e_high_level}, we obtain for the first term
\begin{align*}
    &\Lp{\frac{1}{\sqrt{n}} \sum_{k=1}^n (g(x_k, \xi_k) - g(x^*, \xi_k))}
    \\ 
    \leq &
    \frac{1}{\sqrt{n}} \Lp{\sum_{k=1}^n \sqrt{\gamma_k} (A(\xi_k) - \bar{A}) y_k}
    +
    \frac{\lip(\nabla g)}{\sqrt{n}} \Lp{\sum_{k=1}^n \gamma_k \lVert y_k \rVert_{L_{2p}}^2
    }
    \\ 
    \leq &
    \frac{1}{\sqrt{n}} \Lp{\sum_{k=1}^n \sqrt{\gamma_k} (A(\xi_k) - \bar{A}) y_k}
    +
    \frac{\left(\lip(\nabla g) \max_{k \leq n} \lVert y_k \rVert^2_{L_{2p}} \right)}{1-a} n^{1/2} \gamma_n
    + \mathcal{O}\left(\frac{1}{\sqrt{n}} \right) 
    \numberthis \label{eq:pr_d_bound_first}
    .
\end{align*}
The first term in the bound \eqref{eq:pr_d_bound_first} is analyzed next using the Poisson equation
\begin{align*}
    A(\xi_k) - \bar{A} = \Phi^A (\xi_k) - (P \Phi^A) (\xi_k) . 
\end{align*}
The sum is decomposed as
\begin{equation}\label{eq:pr_d_first_term}
\begin{split}
    \frac{1}{\sqrt{n}} \sum_{k=1}^n \sqrt{\gamma_k} (A(\xi_k) - \bar{A}) y_k
    &=
    \frac{1}{\sqrt{n}}\sum_{k=1}^n (\sqrt{\gamma_k} \Phi^A (\xi_k) y_k - \sqrt{\gamma_{k+1}} \Phi^A (\xi_{k+1}) y_{k+1})
    \\ &+
    \frac{1}{\sqrt{n}} \sum_{k=1}^n \Phi^A (\xi_{k+1}) (\sqrt{\gamma_{k+1}} y_{k+1} - \sqrt{\gamma_k} y_k)
    \\ &
    + \frac{1}{\sqrt{n}} \sum_{k=1}^n \sqrt{\gamma_k} (\Phi^A(\xi_{k+1}) - (P\Phi^A)(\xi_{k})) y_k .     
\end{split}    
\end{equation}
After telescoping, the $L_p$ norm of the first term in \eqref{eq:pr_d_first_term} is bounded using the triangle inequality and the Cauchy-Schwartz inequality \eqref{eq:holder} as 
\begin{align*}
    &\Lp{\frac{1}{\sqrt{n}}\sum_{k=1}^n (\sqrt{\gamma_k} \Phi^A (\xi_k) y_k - \sqrt{\gamma_{k+1}} \Phi^A (\xi_{k+1}) y_{k+1})}
    \\ 
    \leq&
    \frac{2}{\sqrt{n}} \beta_{2p}^{\Phi^A} \left(\sqrt{\gamma_1} \lVert y_1 \rVert_{L_{2p}} + \sqrt{\gamma_{n+1}} \lVert y_{n+1} \rVert_{L_{2p}} \right)
    \\ 
    =& \mathcal{O}\left(\frac{1}{\sqrt{n}} \right) 
    \numberthis \label{eq:pr_d_identity_first_bound}
    . 
\end{align*}
The second term in \eqref{eq:pr_d_first_term} is bounded using Minkowski's inequality
\begin{align*}
    &\Lp{\frac{1}{\sqrt{n}} \sum_{k=1}^n \Phi^A (\xi_{k+1}) (\sqrt{\gamma_{k+1}} y_{k+1} - \sqrt{\gamma_k} y_k)}
    \\
    \leq
    &\frac{\beta_{2p}^{\Phi^A}}{\sqrt{n}} \sum_{k=1}^n \left\lVert (\sqrt{\gamma_{k+1}}y_{k+1} - \sqrt{\gamma_k} y_k\right\rVert_{L_{2p}}
    \\ 
    \leq& \frac{\beta_{2p}^{\Phi^A}}{\sqrt{n}} \sum_{k=1}^n 
    \left(\lvert \sqrt{\gamma_{k+1}} - \sqrt{\gamma_k} \rvert \lVert y_{k+1} \rVert_{L_{2p}}
    + \sqrt{\gamma_k}\lVert y_{k+1} - y_k\rVert_{L_{2p}}
    \right)
    \\ 
    \leq& 
    \frac{\beta_{2p}^{\Phi^A}}{\sqrt{n}} 
    \sum_{k=1}^n \left(\frac{\max_{k \leq {n+1}} \lVert y_k \rVert_{L_{2p}}}{\sqrt{\gamma_k} k} + \gamma_k (\beta_{2p}^f + \beta_{2p}^W) \right)
    + 
    \mathcal{O}\left(\frac{1}{\sqrt{n}} \right)
    ,
\end{align*}
where we used \Cref{lem:yk_variation} for the bound $\lVert y_{k+1} - y_k \rVert_{L_{2p}} \leq \sqrt{\gamma_k} (\beta_{2p}^f + \beta_{2p}^W) + \mathcal{O}(\gamma_k)$. 
Computing the sum, we obtain the bound on $\Lp{(b)}$:
\begin{equation}\label{eq:pr_d_identity_second_bound}
\begin{split}
    &\Lp{\frac{1}{\sqrt{n}} \sum_{k=1}^n \Phi^A (\xi_{k+1}) (\sqrt{\gamma_{k+1}} y_{k+1} - \sqrt{\gamma_k} y_k)}
    \\ 
    \leq &
        \beta_{2p}^{\Phi^A}
    \left(\frac{2 \max_{k \leq {n+1}} \lVert y_k \rVert_{L_{2p}} }{a \sqrt{n \gamma_n}} + (\beta_{2p}^f + \beta_{2p}^W) \frac{\gamma_n}{1-a} n^{1/2} \right)
    + \mathcal{O}\left(\frac{1}{\sqrt{n}} \right) . 
\end{split}
\end{equation}
Lastly, the third term in \eqref{eq:pr_d_first_term} is a MDS sum and can be bounded using the BR inequality
\begin{align*}
    \Lp{\frac{1}{\sqrt{n}} \sum_{k=1}^n \sqrt{\gamma_k} (\Phi^A (\xi_{k+1}) - (P\Phi^A) (\xi_k)) y_k }
    &= \mathcal{O}\left(\frac{1}{\sqrt{n}} \sqrt{\sum_{k=1}^n \gamma_k} \right)
    \\ &
    = \mathcal{O}\left( \sqrt{\gamma_n} \right).
    \numberthis \label{eq:pr_d_identity_third_bound}
\end{align*}
Substituting equations \eqref{eq:pr_d_identity_first_bound}, \eqref{eq:pr_d_identity_second_bound}, and \eqref{eq:pr_d_identity_third_bound} for the first term in \eqref{eq:pr_d_bound_first}, we obtain
\begin{equation}\label{eq:pr_d_bound}
\begin{split}
    \Lp{(d)}
    &\leq 
    \beta_{2p}^{\Phi^A}
    \left(\frac{2 \max_{k \leq {n+1}} \lVert y_k \rVert_{L_{2p}} }{a \sqrt{n \gamma_n}} + (\beta_{2p}^f + \beta_{2p}^W) \frac{\gamma_n}{1-a} n^{1/2} \right)
    \\ 
    &+
    \frac{\left(\lip(\nabla g) \max_{k \leq n} \lVert y_k \rVert^2_{L_{2p}} \right)}{1-a} n^{1/2} \gamma_n
    + \mathcal{O}\left(\sqrt{\gamma_n}\right) .     
\end{split}
\end{equation}




\textbf{Term (e):}
Recall $g(x^*, \xi_k) = f(x^*, \xi_k) - \bar{f}(x^*)$. 
We again use the Poisson equation
\begin{align*}
    g(x^*, \xi_k) + M_k + W_k
    &=
    f(x^*, \xi_k) - \bar{f}(x^*) + M_k + W_k
    \\ &
    = \Phi(\xi_k) - (P \Phi)(\xi_k) - (\Phi(\xi_k) - (P\Phi)(\xi_{k-1}))
    \\ &
    = (P \Phi)(\xi_{k-1}) - (P\Phi)(\xi_k)  . 
\end{align*}
The sum is a telescopic series, and we obtain
\begin{equation}\label{eq:pr_e_bound}
    \frac{1}{\sqrt{n}} \Lp{\sum_{k=1}^n \left( g(x^*, \xi_k) + M_k + W_k\right) }
    =
    \frac{1}{\sqrt{n}} \Lp{(P\Phi)(\xi_0) - (P\Phi)(\xi_n)}
    = \mathcal{O}\left(\frac{1}{\sqrt{n}} \right) .     
\end{equation}
Using the triangle inequality for the $\mathcal{W}_p$ distance, we obtain from \eqref{eq:pr_decomposition} the bound
\begin{equation}
\begin{split}
    &\mathcal{W}_{p} \left( \bar{A} \bar{y}_n, \Gamma^{1/2}Z \right) 
    \\ 
    \leq
    &
    \mathcal{W}_{p} \left(\frac{1}{\sqrt{n}} \sum_{k=1}^n M_k, \Gamma^{1/2} Z \right)
    +
    \Lp{(a)} + \Lp{(b)}
    + \Lp{(c)}
    + \Lp{(d)}
    + \Lp{(e)} .     
\end{split}
\end{equation}
Substituting the respective bounds \eqref{eq:pr_b_bound}, \eqref{eq:pr_c_bound}, \eqref{eq:pr_c_bound}, \eqref{eq:pr_d_bound}, \eqref{eq:pr_e_bound}, we obtain 
\begin{equation}
\begin{split}
    \mathcal{W}_{p} \left( \bar{A} \bar{y}_n, \Gamma^{1/2}Z \right) 
    &\leq 
    \mathcal{W}_{p} \left(\frac{1}{\sqrt{n}} \sum_{k=1}^n M_k, \Gamma^{1/2}Z \right)
    \\ 
    &+    
    \frac{\max_{k \leq n+1} \lVert y_k \rVert_{L_p}}{\sqrt{\gamma_n n}}
    +
    \frac{\lip(\nabla \bar{f}) (\max_{k \leq n} \lVert y_k \rVert^2_{L_{2p}})}{1-a} n^{1/2} \gamma_n
    \\ 
    &
    + \beta_{2p}^{\Phi^A}
    \left(\frac{2\max_{k \leq {n+1}} \lVert y_k \rVert_{L_{2p}} }{a \sqrt{n \gamma_n}} + (\beta_{2p}^f + \beta_{2p}^W) \frac{\gamma_n}{1-a} n^{1/2} \right)
    \\ 
    &+
    \frac{\left(\lip(\nabla g) \max_{k \leq n} \lVert y_k \rVert^2_{L_{2p}} \right)}{1-a} n^{1/2} \gamma_n
    \\ &
    + \mathcal{O}\left(\sqrt{\gamma_n}\right) .     
\end{split}
\end{equation}
The statement is presented with $\lip(\nabla g) = \lip(\nabla f - \nabla \bar{f}) \leq \lip(\nabla f) + \lip(\nabla \bar{f})$.




\subsection{Proof of Corollary \ref{cor:high_probability}}\label{app:high_probability}
        Let $Z \sim \mathcal{N}(0, \identity)$ be a Gaussian variable and consider the optimal coupling $(y_k, \Sigma_a^{1/2} Z)$. 
        By the triangle inequality and union bound, 
    \begin{align*}
        \mathbb{P}\left( \lVert y_k \rVert \geq \dummy \right) 
        &\leq 
        \mathbb{P}\left( \lVert \Sigma_a^{1/2} Z \rVert + \lVert y_k - \Sigma_a^{1/2} Z \rVert \geq \dummy \right) 
        \\ &
        \leq 
        \inf_{\theta \in (0, \dummy)} \left\{ \mathbb{P}\left(\lVert \Sigma_a^{1/2} Z \rVert \geq \dummy - \theta \right) + \mathbb{P}\left( \lVert y_k - \Sigma_a^{1/2} Z \rVert \geq \theta \right) \right\}. 
        \numberthis \label{eq:c1_union}
    \end{align*}
    Observing that $\lVert \Sigma_a^{1/2} Z \rVert$ is $\sqrt{\lVert \Sigma_a\rVert}$--Lipschitz continuous, we use the concentration for Lipschitz functions of Gaussian variables \citep[Theorem 5.6]{boucheron12} to obtain
    \begin{align*}
        \mathbb{P}\left( \lVert \Sigma_a^{1/2} Z \rVert \geq \mathbb{E}\lVert \Sigma_a^{1/2} Z \rVert + t \right) 
        \leq \exp \left(-\frac{t^2}{2\lVert \Sigma_a\rVert}\right) , \quad \forall t > 0. 
    \end{align*}
    Substituting for \eqref{eq:c1_union} and applying Markov's inequality, we obtain
\begin{equation}\label{eq:wasserstein_tail}
    \mathbb{P}(\lVert y_n \rVert \geq \dummy )
    \leq \inf_{\theta \in (0, \dummy)}  \left\{ \exp \left(-\frac{\left(\dummy - \theta - \mathbb{E}\lVert \Sigma_a^{1/2} Z \rVert\right)^2}{2 \lVert \Sigma_a \rVert} \right)
    + \frac{\mathcal{W}_p^{p} \left(y_n, \Sigma_a^{1/2} Z\right)}{\theta^p} \right\}.             
\end{equation}
From \Cref{thm:general}, we have that $\mathcal{W}_p (y_n, \Sigma^{1/2} Z) \leq C_{p,1} \gamma_n^{1/6} + C_{p, 2} \gamma_n^{-1} E_n^a$ for some $C_{p,1}, C_{p,2}$. 
For sufficiently large $\dummy > 0$, both terms in \eqref{eq:wasserstein_tail} are $\leq \delta/2$ when $\theta$ is chosen to satisfy
\begin{align*}
    \theta \leq \dummy - \left(\mathbb{E}\lVert \Sigma_a^{1/2}Z \rVert + \sqrt{2 \lVert \Sigma_a \rVert \log \frac{2}{\delta}} \right) 
    , \qquad 
    \theta \geq \mathcal{W}_p (y_n, \Sigma_a^{1/2} Z) \left( \frac{1}{\delta} \right)^{1/p}
    .
\end{align*}
For the choice
\begin{align*}
    \dummy = \left(\mathbb{E}\lVert \Sigma_a^{1/2}Z \rVert + \sqrt{2 \lVert \Sigma_a \rVert \log \frac{2}{\delta}} \right)  +
    \mathcal{W}_p (y_n, \Sigma_a^{1/2} Z) \left( \frac{2}{\delta} \right)^{1/p} , 
\end{align*}
we obtain that with probability $\geq 1 - \delta$,
\begin{align*}
    \lVert y_n \rVert 
    &\leq \left(\mathbb{E}\lVert \Sigma_a^{1/2} Z \rVert + \sqrt{2 \lVert \Sigma_a \rVert \log \frac{2}{\delta}} + \mathcal{W}_p (y_n, \Sigma_a^{1/2} Z) \left(\frac{2}{\delta} \right)^{1/p} \right)        
    \\     
    &\leq \left(\mathbb{E}\lVert \Sigma_a^{1/2} Z \rVert + \sqrt{2 \lVert \Sigma_a \rVert \log \frac{2}{\delta}} + \left(C_{p, 1} \gamma_n^{1/6} + C_{p, 2} \frac{E_n^a}{\gamma_n} \right) \left(\frac{2}{\delta} \right)^{1/p} \right)         .
\end{align*}
Multiplying both sides by $\sqrt{\gamma_n}$, we obtain the first result. 


The second result follows using the same steps with the coupling $(\bar{y}_n, \bar{\Sigma}^{1/2}Z)$. 
Substituting $\mathcal{W}_p (\bar{y}_n, (\bar{A}^{-1} \Gamma \bar{A}^{-\dagger})^{-1/2} Z) \leq \bar{C}_p n^{-1/6}$, which holds when $a = 2/3$ to optimize the rate in \Cref{thm:pr_average}, we obtain
\begin{align*}
    \lVert \bar{y}_n \rVert 
    \leq 
    \left( 
    \mathbb{E}\lVert \bar{\Sigma}^{1/2}Z \rVert
    + \sqrt{2 \lVert \bar{\Sigma} \rVert \log \frac{2}{\delta}}
    \right)
    + 
    \frac{\bar{C}_{p}}{n^{1/6}} \left(\frac{2}{\delta} \right)^{1/p} . 
\end{align*}
Dividing both sides by $\sqrt{n}$, we obtain the result. 



\section{Proof of Auxiliary Lemmas}

\subsection{Proof of \Cref{lem:e_bound}}\label{app:e_state_dependent}
Recall from \eqref{eq:poisson} the solution $\Phi^{A}$ to the Poisson equation
\begin{align*}
    A(\xi_k) - \mathbb{E}A(\xi) = 
    \Phi^A (\xi_k) - \mathbb{E}_k \Phi^A (\xi_{k+1})
    .
\end{align*}
The quantity of interest is the $L_p$ norm of
\begin{align*}
    &\sum_{k=k_m}^{k_{m+1} - 1} \alpha_k \sqrt{\gamma_k} (A(\xi_k) - \mathbb{E}A) y_k 
    \\=& 
    \sum_{k=k_m}^{k_{m+1} - 1} \alpha_k \sqrt{\gamma_k} (\Phi^A (\xi_k) - \Phi^A(\xi_{k+1})) y_k 
    + \sum_{k=k_m}^{k_{m+1} - 1} \alpha_k \sqrt{\gamma_k} \left(\Phi^A (\xi_{k+1}) - \mathbb{E}_{k} \Phi^A (\xi_{k+1}) \right) y_k 
    \\ 
    = 
    &
    \sum_{k=k_m}^{k_{m+1} - 1} (\alpha_k \sqrt{\gamma_k} \Phi^A(\xi_k) y_k - \alpha_{k+1} \sqrt{\gamma_{k+1}} \Phi^A(\xi_{k+1}) y_{k+1}) \numberthis \label{eq:e1}
    \\ + &
    \sum_{k=k_m}^{k_{m+1} - 1} \Phi^A (\xi_{k+1}) (\alpha_{k+1} \sqrt{\gamma_{k+1}} y_{k+1} - \alpha_k \sqrt{\gamma_k} y_k ) \numberthis \label{eq:e2}
    \\ +
    &\sum_{k=k_m}^{k_{m+1} - 1} \alpha_k \sqrt{\gamma_k} \left(\Phi^A (\xi_{k+1}) - \mathbb{E}_{k} \Phi^A (\xi_{k+1}) \right) y_k 
    \numberthis \label{eq:e3}
    .
\end{align*}
The $L_p$ norm of \eqref{eq:e1}--\eqref{eq:e3} are bounded separately to obtain the result. 



The norm of \eqref{eq:e1} is bounded by telescoping the sum and applying the Minkowski inequality and the Cauchy-Schwartz inequality \eqref{eq:holder}:
\begin{align*}
    &\Lp{    \sum_{k=k_m}^{k_{m+1} - 1} (\alpha_k \sqrt{\gamma_k} \Phi^A(\xi_k) y_k - \alpha_{k+1} \sqrt{\gamma_{k+1}} \Phi^A(\xi_{k+1}) y_{k+1})}
    \\ 
    =&
    \Lp{\alpha_{k_m} \sqrt{\gamma_{k_m}} \Phi^A (\xi_{k_m}) y_{k_m} - \alpha_{k_{m+1}} \sqrt{\gamma_{k_{m+1}}} \Phi^A (\xi_{k_{m+1}}) y_{k_{m+1}}}
    \\ \leq & 
    \beta_{2p}^{\Phi^A} \lVert \Sigma_a^{1/2} Z \rVert_{L_{2p}}
    \left( \alpha_{k_m} \sqrt{\gamma_{k_m}} \beta_{2p}^y (m)  + \alpha_{k_{m+1}} \sqrt{\gamma_{k_{m+1}}} \beta_{2p}^y (m+1)\right) 
    \\ 
    = &\mathcal{O}\left(\gamma_{k_m}  \right) , 
\end{align*}
where we substituted $\alpha_k = \gamma_k/\sqrt{\gamma_{k+1}} = \mathcal{O}(\sqrt{\gamma_k})$. 



For \eqref{eq:e2}, we use the identity
\begin{align*}
    \alpha_{k+1} \sqrt{\gamma_{k+1}} y_{k+1} - \alpha_k \sqrt{\gamma_k} y_k 
    &=
    \alpha_{k+1} \sqrt{\gamma_{k+1}} (y_{k+1} - y_k) + (\alpha_{k+1} \sqrt{\gamma_{k+1}} - \alpha_k \sqrt{\gamma_k}) y_k 
\end{align*}
and substitute to obtain
\begin{align*}
    &\Lp{    \sum_{k=k_m}^{k_{m+1} - 1} \Phi^A (\xi_{k+1}) (\alpha_{k+1} \sqrt{\gamma_{k+1}} y_{k+1} - \alpha_k \sqrt{\gamma_k} y_k )}
    \\
    \leq &
    \Lp{\sum_{k=k_m}^{k_{m+1} - 1} \alpha_{k+1} \sqrt{\gamma_{k+1}}\Phi^A(\xi_{k+1})  \left( y_{k+1} -  y_k   \right)}
    \\ 
    + &
    \Lp{\sum_{k=k_m}^{k_{m+1}}
    (\alpha_{k+1} \sqrt{\gamma_{k+1}} - \alpha_k \sqrt{\gamma_k}) \Phi^A (\xi_{k+1}) y_k 
    }
    \numberthis \label{eq:e_ub_yrecursion}
\end{align*}
Since $\alpha_k \sqrt{\gamma_k} \approx \gamma_k$ is a power decay sequence with $\lvert \gamma_{k+1} - \gamma_k\rvert \leq \gamma_k/k$, we have $\lvert \alpha_{k+1} \sqrt{\gamma_{k+1}} - \alpha_k \sqrt{\gamma_k}\rvert \leq C \gamma_k/k$ for some $C > 0$ and by the Minkowski inequality,
\begin{align*}
    \Lp{\sum_{k=k_m}^{k_{m+1}}
    (\alpha_{k+1} \sqrt{\gamma_{k+1}} - \alpha_k \sqrt{\gamma_k}) \Phi^A (\xi_{k+1}) y_k }
    &\leq 
    C \sum_{k=k_m}^{k_{m+1}} \frac{\gamma_k}{k} \beta_{2p}^{\Phi^A} \beta_{2p}^y (m) \lVert \Sigma_a^{1/2}Z  \rVert_{L_{2p}}
    \\ &
    \leq
    C \frac{\cumstep_m}{k_m} \beta_{2p}^{\Phi^A} \beta_{2p}^y(m) \lVert \Sigma_a^{1/2} Z \rVert_{L_{2p}} . 
    \numberthis \label{eq:temp1}
\end{align*}
For the first term in the upper bound \eqref{eq:e_ub_yrecursion}, we use that $\lvert \alpha_k \sqrt{\gamma_k} \rvert \leq \sqrt{2} \gamma_k$, the Minkowski inequality, and \Cref{lem:yk_variation}:
\begin{align*}
    &\Lp{\sum_{k=k_m}^{k_{m+1} - 1} \alpha_{k+1} \sqrt{\gamma_{k+1}} \Phi^A (\xi_{k+1}) (y_{k+1} - y_k)}
    \\ 
    \leq&
    \sqrt{2} \sum_{k=k_m}^{k_{m+1} - 1} \gamma_{k+1} \beta_{2p}^{\Phi^A} \lVert y_{k+1} - y_k\rVert_{L_2p}
    \\ 
    \leq& 
    \sqrt{2} \cumstep_m \beta_{2p}^{\Phi^A} \left(\sqrt{\gamma_{k_m}} (\beta_{2p}^f + \beta_{2p}^W)
    + \mathcal{O}\left(\gamma_{k_m} \beta_{2p}^y (m) + \frac{\sqrt{\gamma_{k_m}}}{k_m} \right)\right) . 
\end{align*}
The $\cumstep_m/k_m$ rate in \eqref{eq:temp1} is negligible and we have the following upper bound on \eqref{eq:e2}:
\begin{equation}
    \eqref{eq:e2} \leq 
        \sqrt{2} \cumstep_m \beta_p^{\Phi^A} \left(\sqrt{\gamma_{k_m}} (\beta_p^f + \beta_p^W)
    + \mathcal{O}\left(\gamma_{k_m} \beta_{2p}^y (m) + \frac{\sqrt{\gamma_{k_m}}}{k_m} \right)\right)  
\end{equation}
Next we obtain a bound on the $L_p$ norm of \eqref{eq:e3}. 
The quantity is a sum of a martingale difference sequence, since
\begin{align*}
    \mathbb{E}\left[\alpha_k \sqrt{\gamma_k} (\Phi^A (\xi_{k+1}) - \mathbb{E}_k \Phi^A (\xi_{k+1}) )y_k \Big| \mathcal{H}_k \right] = 0 . 
\end{align*}
Applying the BR inequality as in \eqref{eq:br_inequality}, we have that 
\begin{align*}
&\Lp{\sum_{k=k_m}^{k_{m+1} - 1} \alpha_k \sqrt{\gamma_k} (\Phi^A (\xi_{k+1}) - P \Phi^A (\xi_k) ) y_k}
\leq   
     C_p^{BR}    \sqrt{\sum_{k=k_m}^{k_{m+1} - 1}\alpha_k^2 \gamma_k}
     \\ &
     \cdot      \sup_{k \geq 1}\Big\{
\sqrt{ \left\lvert \mathrm{Tr} \mathbb{E}\left[
    (\Phi^A (\xi_{k+1}) - P \Phi^A(\xi_k)) y_k y_k^\dagger (\Phi^A (\xi_{k+1}) - P \Phi^A(\xi_k))^\dagger \Big| \mathcal{H}_{k-1} 
    \right] \right\rvert_{L_{p/2}}} \\ 
    &+
    \left\lVert \left(\Phi^A(\xi_{k+1}) - P \Phi^A (\xi_{k})\right) y_k \right\rVert_{L_p}
    \Big\}
    .
\end{align*}
The two quantities inside the $\sup_{k \geq 1}$ are bounded: For any $p \geq 2$, we have by the cyclic property of the trace
\begin{align*}
    &\left\lvert \mathrm{Tr} \mathbb{E}\left[
    (\Phi^A (\xi_{k+1}) - P \Phi^A(\xi_k)) y_k y_k^\dagger (\Phi^A (\xi_{k+1}) - P \Phi^A(\xi_k))^\dagger \Big| \mathcal{H}_{k-1} 
    \right]\right\rvert_{L_{p/2}}
    \\ = &
    \left\lvert \mathrm{Tr} y_k^\dagger \mathbb{E}\left[ 
    (\Phi^A(\xi_{k+1} ) - P \Phi^A(\xi_k) ) (\Phi^A(\xi_{k+1} ) - P \Phi^A(\xi_k) ) ^\dagger | \mathcal{H}_{k-1}
    \right] y_k^\dagger \right\rvert_{L_{p/2}}
    \\ \leq &
    \left\lvert 
    \left\lVert \mathbb{E} \left[(\Phi^A(\xi_{k+1} ) - P \Phi^A(\xi_k) ) (\Phi^A(\xi_{k+1} ) - P \Phi^A(\xi_k) ) ^\dagger | \mathcal{H}_{k-1} \right] \right\rVert \lVert y_k \rVert^2
    \right\rvert_{L_{p/2}}
    \\ \leq &
    \beta_p^{\Phi^A} \lVert y_k \rVert^2_{L_p} . 
\end{align*}
The other quantity in the $\sup$ is bounded using the Cauchy-Schwartz and the triangle inequality $\lVert (\Phi^A(\xi_{k+1}) - P \Phi^A(\xi_k))y_k \rVert_{L_p} \leq 2\beta_{2p}^{\Phi^A} \beta_{2p}^y$. 
Together we have by $\alpha_k = \gamma_k/\sqrt{\gamma_{k+1}}$ the bound
\begin{equation}
    \lVert \eqref{eq:e3}\rVert_{L_p} = \mathcal{O}\left(\sqrt{\sum_{k=k_m}^{k_{m+1} - 1} \gamma_k^{2} }\right) 
    = \mathcal{O}\left(\gamma_{k_m} \sqrt{I_m}\right) . 
\end{equation}


Recall the rates for \eqref{eq:e1}--\eqref{eq:e3} are respectively $\mathcal{O}(\gamma_{k_m} )$, $\mathcal{O}(\cumstep_m \sqrt{\gamma_{k_m}})$, and $\mathcal{O}(\gamma_{k_m}\sqrt{I_m} )$. 
The $\mathcal{O}(\gamma_{k_m}\sqrt{I_m} )$ rate in \eqref{eq:e3} dominates, and we have 
\begin{align*}
    &\left\lVert 
    \sum_{k=k_m}^{k_{m+1} - 1} \alpha_k \sqrt{\gamma_k} (A(\xi_k) - \mathbb{E}A(\xi)) y_k 
    \right\rVert_{L_p}
    = \mathcal{O}(\gamma_{k_m} \sqrt{I_m}) . 
\end{align*}

\subsection{Proof of Lemma \ref{lem:ou_discretization}}\label{app:ou_discretization}
Convergence of the Euler-Maruyama discretization $\{u_m\}$ of the Ornstein-Uhlenbeck process $(U_t)$ is proved when using a diminishing \stepsize sequence $\{\cumstep_m\}$. 
The processes $\{u_m\}$ and $(U_t)$ are described by the equations
\begin{align*}
    u_{m+1} &= u_m - \cumstep_m \bar{A}_a u_m + \sqrt{\cumstep_m \Gamma} Z_m 
    \\ 
    d U_t &= - \bar{A}_a U_t dt + \sqrt{\Gamma} dB_t . 
\end{align*}
The SDE is solved in the interval $[\tau_m, \tau_{m+1}]$, where the solution is given by
\begin{align*}
    U_{\tau_{m+1}} &= e^{-\bar{A}_a \cumstep_m} U_{\tau_m}
    +
    \int_{\tau_m}^{\tau_{m+1}} e^{-\bar{A}_a (\tau_{m+1} - t)} \sqrt{\Gamma} dB_t . 
\end{align*}
Subtracting the continuous time and discrete time equations, we have
\begin{align*}
    U_{\tau_{m+1}} - u_{m+1} 
    &= 
    e^{-\bar{A}_a \cumstep_m} U_{\tau_m} - (\identity - \cumstep_m \bar{A}_a)u_m 
    + 
    \int_{\tau_m}^{\tau_{m+1}} e^{-\bar{A}_a (\tau_{m+1} - t)} \sqrt{\Gamma} dB_t  - 
    \sqrt{\cumstep_m \Gamma} Z_m 
    \\ &
    \underset{(*)}{=} (\identity - \cumstep_m \bar{A}_a) (U_{\tau_m} - u_m)
    + 
    \sum_{j=2}^\infty \frac{(-\cumstep_m \bar{A}_a)^j}{j!} U_{\tau_m}
    \\ &
    +
    \int_{\tau_m}^{\tau_{m+1}} \left(e^{-\bar{A}_a (\tau_{m+1} - t)} - \identity \right)  \sqrt{\Gamma} dB_t
    ,
\end{align*}
where $(*)$ used a synchronized coupling between the two Brownian motions driving the discrete-time and continuous time processes, and the Taylor expansion of the matrix exponential:
\begin{align*}
    \left\lVert 
    \sum_{j=2}^\infty \frac{(-\cumstep_m \bar{A}_a)^j}{j!} U_{\tau_m}
    \right\rVert_{L_p}
    \leq 
    \sum_{j=2}^\infty \frac{\cumstep_m^j}{j!} \lVert \bar{A}_a \rVert \lVert U_{\tau_m} \rVert_{L_p}
    &=
    \lVert \bar{A}_a \rVert \lVert U_{\tau_m} \rVert_{L_p} \cdot \left( e^{\cumstep_m} - 1 - \cumstep_m
    \right)
    \\ &
    \leq 
    \frac{e}{2} \cumstep_m^2\lVert \bar{A}_a \rVert \lVert U_{\tau_m} \rVert_{L_p} ,
\end{align*}
where we used $\lvert e^x - 1 - x \rvert \leq x^2 e^x / 2 \leq e x^2/2$ when $x \in [0, 1]$. 
For the synchronous coupling above, it therefore holds that
\begin{equation}\label{eq:OU_recursion1}
\begin{split}
    \lVert U_{\tau_{m+1}} - u_{m+1} \rVert_{L_p}
    &\leq
    (1 - \lambdaDT \cumstep_m) \lVert U_{\tau_m} - u_m \rVert_{L_p}
    + 
    \left\lVert     \int_{\tau_m}^{\tau_{m+1}} \left(e^{-\bar{A}_a (\tau_{m+1} - t)} - \identity \right)  \sqrt{\Gamma} dB_t\right\rVert_{L_p}
    \\ &
    + \frac{e}{2} \cumstep_m^2\lVert \bar{A}_a \rVert \lVert U_{\tau_m} \rVert_{L_p} . 
\end{split}
\end{equation}
The quantity $\Lp{U_{\tau_m}}$ is finite for all $m \geq 1$ when the initial condition $U_{\tau_1}$ is in $L_p$. 


The $L_p$ norm of the It\^{o} integral is bounded using the Burkholder-Davis-Gundy (BDG) inequality \citep[Theorem 5.16]{le2016brownian}. 
We first relate the weighted norm to the Euclidean norm and then apply the BDG inequality, which is evaluated as
\begin{align*}
    &\mathbb{E}\left[\sup_{t \in [\tau_m, \tau_{m+1}]} \left\lVert  
    \int_{\tau_m}^{t} \left(e^{-\bar{A}_a (\tau_{m+1} - t)} - \identity \right)  \sqrt{\Gamma} dB_t
    \right\rVert^p 
    \right]
    \\ \leq &
    \lVert Q \rVert_2^{p/2}
    \mathbb{E}\left[\sup_{t \in [\tau_m, \tau_{m+1}]} \left\lVert  
    \int_{\tau_m}^{t} \left(e^{-\bar{A}_a (\tau_{m+1} - t)} - \identity \right)  \sqrt{\Gamma} dB_t
    \right\rVert^p_2 
    \right]
    \\
    \leq & 
    \left( \lVert Q \rVert_2^{1/2}C_p^{BDG} \right)^p \left( 
    \int_{\tau_m}^{\tau_{m+1}}\mathrm{Tr}\left(e^{-\bar{A}_a (\tau_{m+1} - t)} - \identity \right) \Gamma \left(e^{-\bar{A}_a (\tau_{m+1} - t)} - \identity \right)^\dagger  dt
    \right)^{p/2} 
    \\ 
    = &
    \left( \lVert Q \rVert_2^{1/2}C_p^{BDG} \right)^p \left( 
    \int_{0}^{\cumstep_m}\mathrm{Tr}\left(e^{-\bar{A}_a t} - \identity \right) \Gamma \left(e^{-\bar{A}_a t} - \identity \right)^\dagger  dt
    \right)^{p/2} 
\end{align*}
for some constant $C_p^{BDG} > 0$. 
Using the property
\begin{align*}
    \mathrm{Tr}\left(e^{-\bar{A}_a t} - \identity \right) \Gamma \left(e^{-\bar{A}_a t} - \identity \right)^\dagger
    &\leq 
    \lVert  e^{-\bar{A}_a t} - \identity  \rVert_2^2 \mathrm{Tr} \Gamma  
\end{align*}
and the identity $e^{-\bar{A}_a t} - \identity = \int_0^1 (t \bar{A}_t) e^{-s t \bar{A}_a} ds $, we obtain for some constant $K > 0$ the bounds
\begin{align*}
    \lVert e^{-\bar{A}_a t} - \identity \rVert
    \leq 
    \lVert \bar{A}_a \rVert t \int_0^1 \lVert e^{-st \bar{A}_a} \rVert  ds 
    &\leq 
    K \lVert \bar{A}_a  \rVert t \int_0^1 e^{-\lambdaCT t s} ds
    \\ &
    \leq
    K \frac{\lVert \bar{A}_a \rVert t}{\lambdaCT} 
\end{align*}
for all $t > 0$. 
Substituting, we obtain 
\begin{align*}
&\mathbb{E}\left[\sup_{t \in [\tau_m, \tau_{m+1}]} \left\lVert  
    \int_{\tau_m}^{t} \left(e^{-\bar{A}_a (\tau_{m+1} - t)} - \identity \right)  \sqrt{\Gamma} dB_t
    \right\rVert^p 
    \right]
    \\ 
    \leq & 
        \left(K \lVert Q \rVert_2^{1/2}C_p^{BDG} \right)^p \left( 
    \int_{0}^{\cumstep_m}\left(\frac{\lVert \bar{A}_a \rVert t}{\lambdaCT} \right)^2 \mathrm{Tr} \Gamma  dt
    \right)^{p/2}     
    \\ 
    = &
    \left(K \frac{\lVert Q \rVert_2^{1/2}C_p^{BDG}}{\sqrt{3}} \right)^p
    \left(\frac{\lVert \bar{A}_a \rVert }{\lambdaCT} \sqrt{\mathrm{Tr} \Gamma} \right)^p
    \cumstep_m^{3p/2} . 
\end{align*}
From \eqref{eq:OU_recursion1}, we obtain the recursion
\begin{equation}
\begin{split}
    \lVert U_{\tau_{m+1}} - u_{m+1} \rVert_{L_p}
    &\leq (1 - \lambdaDT \cumstep_m) \lVert U_{\tau_m} - u_m \rVert_{L_p}
    + 
    K \frac{\lVert Q \rVert_2^{1/2}C_p^{BDG}}{\sqrt{3}}
    \left(\frac{\lVert \bar{A}_a \rVert }{\lambdaCT} \sqrt{\mathrm{Tr} \Gamma} \right)
    \cumstep_m^{3/2} 
    \\ &
    + \frac{e}{2} \cumstep_m^2\lVert \bar{A}_a \rVert \lVert U_{\tau_m} \rVert_{L_p}
    .
\end{split}
\end{equation}
Solving the recursion, we obtain 
\begin{align*}
    \lVert U_{\tau_{N+1}} - u_{N+1} \rVert_{L_p}
    \leq 
    K \frac{C_p^{BDG}}{\lambdaDT} 
    \frac{\lVert \bar{A}_a \rVert }{\lambdaCT} \sqrt{\frac{\lVert Q \rVert_2\mathrm{Tr} \Gamma}{3}} \sqrt{\cumstep_m}
    + \mathcal{O}\left(\cumstep_m \right) 
    .
\end{align*}
Therefore, we conclude
\begin{equation}
    \mathcal{W}_p (U_{\tau_{N+1}}, u_{N+1})
    \leq
    K \frac{C_p^{BDG}}{\lambdaDT} 
    \frac{\lVert \bar{A}_a \rVert }{\lambdaCT} \sqrt{\frac{\lVert Q \rVert_2\mathrm{Tr} \Gamma}{3}} \sqrt{\cumstep_m}
    + \mathcal{O}\left(\cumstep_m \right) 
    .
\end{equation}




\subsection{Proof of Lemma \ref{lem:ct_ou_convergence}}\label{app:ct_ou_convergence}
Let $U_\infty$ be a stationary limit of the OU process $(U_t)$. 
For a synchronized coupling of the driving Brownian processes, the difference between the states satisfies
\begin{align*}
    d (U_t - U_\infty) &= - \bar{A}_a (U_t - U_\infty)  dt . 
\end{align*}
Therefore, 
\begin{align*}
    \mathcal{W}_p (U_t, U_\infty) \leq  \left\lVert e^{-A_a t} (U_0 - U_\infty) \right\rVert_{L_p}
    \leq K e^{-\lambdaCT t} \mathcal{W}_p (U_0, U_\infty) .
\end{align*}











\bibliographystyle{plainnat} 
\bibliography{bibliography}     

\end{document}
